% Niveaux d'énergie de l'atome d'hydrogène — Physique Tle E — Cours signé OnBuch+
% Programme officiel MINESEC. Compiler avec : tectonic P17-niveaux-energie-atome-hydrogene.tex
\newcommand{\DOCMATIERE}{Physique}
\newcommand{\DOCNIVEAU}{Tle E}
\newcommand{\DOCMODULE}{Module 4 : Onde, matière et transformations nucléaires}
\newcommand{\DOCLECON}{P17}
\newcommand{\DOCTITRE}{Niveaux d'énergie de l'atome d'hydrogène}
% OnBuch+ — préambule « cours signés » ENRICHI (compilé avec tectonic / XeLaTeX)
% Système visuel « Pop » d'OnBuch+ : fond crème (jamais blanc), bordures encre
% épaisses, ombres « dures » décalées, titres Archivo Black en capitales.
% Version MATHÉMATIQUES : graphes (pgfplots/tikz), boîtes pédagogiques
% (définition, propriété, méthode, attention, le savais-tu, exercice résolu,
% exercice, TP, bilan), page de garde et signature OnBuch+ ancrée en bas.
%
% Macros attendues dans main.tex AVANT \input{preamble} :
%   \DOCMATIERE  \DOCNIVEAU  \DOCMODULE  \DOCLECON (numéro)  \DOCTITRE
\documentclass[11pt]{article}

\usepackage{fontspec}
\usepackage[french]{babel}
\usepackage[a4paper,margin=2cm,top=3.4cm,bottom=2.8cm,headheight=1.4cm,footskip=1.3cm]{geometry}
\usepackage{amsmath,amssymb,amsfonts}
\usepackage{mathtools} % \xmapsto, \coloneqq... souvent utilisés par les modèles
\usepackage{cancel} % simplifications barrées (bilans de forces)
% \abs{z} (module, valeur absolue) et \norm{u} (norme), utilisés par les modèles
\providecommand{\abs}{}\let\abs\relax\DeclarePairedDelimiter{\abs}{\lvert}{\rvert}
\providecommand{\norm}{}\let\norm\relax\DeclarePairedDelimiter{\norm}{\lVert}{\rVert}
\usepackage[table]{xcolor} % [table] : fournit aussi \rowcolors (alternance de lignes)
\usepackage{pagecolor}
\usepackage{tikz}
\usetikzlibrary{arrows.meta,positioning,calc,shapes.geometric,shapes.misc,shapes.arrows,decorations.pathmorphing,decorations.markings,patterns,shadows,mindmap,trees,fit,backgrounds,matrix,intersections,angles,quotes}
\usepackage{pgfplots}
\usepackage[version=4]{mhchem} % équations nucléaires \ce{^{235}_{92}U}
\usepackage[european,siunitx]{circuitikz} % schémas électriques
\pgfplotsset{compat=1.18}
\usepgfplotslibrary{fillbetween,groupplots} % aires entre courbes (calcul intégral)
\usepackage{enumitem}
\usepackage{fancyhdr}
% [most] charge toutes les bibliothèques tcolorbox (listings, minted, raster,
% documentation...) et gonfle inutilement la mémoire TeX sur un document long
% (~300 boîtes) : on ne charge que ce qui est réellement utilisé.
\usepackage[skins,breakable]{tcolorbox}
\usepackage{graphicx}
\usepackage{colortbl}
\usepackage{booktabs}
\usepackage{tabularx}
\usepackage{diagbox} % case d'en-tête diagonale des tableaux à double entrée
% Colonnes extensibles centrée / gauche / droite (C, L, R), souvent utilisées par les modèles
\newcolumntype{C}{>{\centering\arraybackslash}X}
\newcolumntype{L}{>{\raggedright\arraybackslash}X}
\newcolumntype{R}{>{\raggedleft\arraybackslash}X}
\usepackage{array}
\usepackage{multicol}
\usepackage{siunitx}
% parse-numbers=false : accepte aussi \SI{2,5 \times 10^{-3}}{...}
\sisetup{locale=FR,output-decimal-marker={,},group-minimum-digits=4,parse-numbers=false}
\DeclareSIUnit\ohm{\ensuremath{\symup{\Omega}}} % U+2126 absent de la police
\DeclareSIUnit\division{div} % réglages d'oscilloscope (ms/div, V/div)
\usepackage{qrcode}
% \degree : commande fréquemment utilisée par les modèles pour un angle ou une
% température en dehors de \SI{}{} (non fournie par les paquets chargés).
\providecommand{\degree}{^{\circ}}
% \qed (fin de démonstration), utilisé par les modèles sans amsthm
\providecommand{\qed}{\hfill$\square$}
% \barre : double barre des valeurs interdites dans un tableau de variations (tkz-tab)
\providecommand{\barre}{\ensuremath{\|}}
% environnement proof (amsthm non chargé) utilisé par les modèles
\ifdefined\proof\else\newenvironment{proof}[1][Démonstration]{\par\noindent\textit{#1.}\ }{\qed\par}\fi
\usepackage{lastpage}
\usepackage{hyperref}
\hypersetup{hidelinks}

% ============================================================
% Palette « Pop » OnBuch+ — mêmes hex que la classe P (plus_ui.dart)
% ============================================================
\definecolor{poporange}{HTML}{F59321}
\definecolor{popdark}{HTML}{E07A0C}
\definecolor{popcream}{HTML}{FFF6E8}
\definecolor{popink}{HTML}{1C1714}
\definecolor{poppurple}{HTML}{7A5AE0}
\definecolor{popgreen}{HTML}{2E9E6A}
\definecolor{poppink}{HTML}{E0457B}
\definecolor{popblue}{HTML}{2F7DE1}
\definecolor{popgold}{HTML}{C98A12}
\definecolor{popmuted}{HTML}{8A7F76}
\definecolor{popline}{HTML}{EADFD2}
\definecolor{popgray}{HTML}{8A7F76}
\colorlet{popgrey}{popgray}
% Alias tolérants (le modèle nomme parfois les couleurs différemment)
\colorlet{popwhite}{white}
\colorlet{popblack}{popink}
\colorlet{popred}{poppink}
\colorlet{popyellow}{popgold}
% Teintes claires pour fonds de tableaux / zones
\colorlet{popblueL}{popblue!10!white}
\colorlet{poporangeL}{poporange!14!white}
\colorlet{popgreenL}{popgreen!12!white}
\colorlet{poppurpleL}{poppurple!12!white}
\colorlet{poppinkL}{poppink!12!white}
\colorlet{popgoldL}{popgold!14!white}

\pagecolor{popcream}

% ============================================================
% Typographie
% ============================================================
\defaultfontfeatures{Ligatures=TeX}
\setmainfont{PlusJakartaSans-Medium.ttf}[Path=../../../fonts/,BoldFont=PlusJakartaSans-Bold.ttf]
% Maths en sans-serif (Fira Math) avec chiffres et lettres droites de Plus
% Jakarta, pour que les formules se fondent dans le texte.
\usepackage[math-style=ISO,bold-style=ISO]{unicode-math}
\setmathfont{FiraMath-Regular.otf}[Path=../../../fonts/]
\setmathfont{PlusJakartaSans-Medium.ttf}[Path=../../../fonts/,range={up/num,up/{Latin,latin},\mathup}]
\usepackage{icomma} % virgule décimale sans espace en mode math
% Caractères mathématiques Unicode absents des polices de texte (Plus Jakarta,
% Archivo Black) : rendus par leur équivalent mathématique, en texte comme en
% formule — sinon ils s'affichent en carré vide (ex. « Arithmétique dans ℤ »).
\usepackage{newunicodechar}
\newunicodechar{ℝ}{\ensuremath{\mathbb{R}}}
\newunicodechar{ℤ}{\ensuremath{\mathbb{Z}}}
\newunicodechar{ℂ}{\ensuremath{\mathbb{C}}}
\newunicodechar{ℕ}{\ensuremath{\mathbb{N}}}
\newunicodechar{ℚ}{\ensuremath{\mathbb{Q}}}
\newunicodechar{𝔻}{\ensuremath{\mathbb{D}}}
\newunicodechar{α}{\ensuremath{\alpha}}
\newunicodechar{β}{\ensuremath{\beta}}
\newunicodechar{γ}{\ensuremath{\gamma}}
\newunicodechar{δ}{\ensuremath{\delta}}
\newunicodechar{ε}{\ensuremath{\varepsilon}}
\newunicodechar{θ}{\ensuremath{\theta}}
\newunicodechar{λ}{\ensuremath{\lambda}}
\newunicodechar{μ}{\ensuremath{\mu}}
\newunicodechar{σ}{\ensuremath{\sigma}}
\newunicodechar{τ}{\ensuremath{\tau}}
\newunicodechar{φ}{\ensuremath{\varphi}}
\newunicodechar{ψ}{\ensuremath{\psi}}
\newunicodechar{ω}{\ensuremath{\omega}}
\newunicodechar{Σ}{\ensuremath{\Sigma}}
% majuscules produites par \MakeUppercase dans les titres (α-aminés -> Α)
\newunicodechar{Α}{\ensuremath{\alpha}}
\newunicodechar{Β}{\ensuremath{\beta}}
\newunicodechar{Γ}{\ensuremath{\Gamma}}
\newunicodechar{Δ}{\ensuremath{\Delta}}
\newunicodechar{Ε}{\ensuremath{\varepsilon}}
\newunicodechar{Θ}{\ensuremath{\Theta}}
\newunicodechar{Λ}{\ensuremath{\Lambda}}
\newunicodechar{Μ}{\ensuremath{\mu}}
\newunicodechar{Τ}{\ensuremath{\tau}}
\newunicodechar{Φ}{\ensuremath{\Phi}}
\newunicodechar{Ψ}{\ensuremath{\Psi}}
\newunicodechar{Ω}{\ensuremath{\Omega}}
\newunicodechar{↦}{\ensuremath{\mapsto}}
\newunicodechar{∈}{\ensuremath{\in}}
\newunicodechar{∉}{\ensuremath{\notin}}
\newunicodechar{∧}{\ensuremath{\wedge}}
\newunicodechar{⇔}{\ensuremath{\Leftrightarrow}}
\newunicodechar{⟺}{\ensuremath{\Longleftrightarrow}}
\newunicodechar{⇒}{\ensuremath{\Rightarrow}}
\newunicodechar{⟹}{\ensuremath{\Longrightarrow}}
\newunicodechar{∘}{\ensuremath{\circ}}
\newunicodechar{∪}{\ensuremath{\cup}}
\newunicodechar{∩}{\ensuremath{\cap}}
\newunicodechar{≡}{\ensuremath{\equiv}}
\newunicodechar{⊂}{\ensuremath{\subset}}
\newunicodechar{⊥}{\ensuremath{\perp}}
\newunicodechar{∃}{\ensuremath{\exists}}
\newunicodechar{∀}{\ensuremath{\forall}}
\newunicodechar{∛}{\ensuremath{\sqrt[3]{\,}}}
\newunicodechar{∜}{\ensuremath{\sqrt[4]{\,}}}
\newunicodechar{⃗}{\ensuremath{{}^{\to}}}
\newunicodechar{ⁿ}{\ifmmode^{n}\else\textsuperscript{\ensuremath{n}}\fi}
\newunicodechar{ˣ}{\ifmmode^{x}\else\textsuperscript{\ensuremath{x}}\fi}
\newunicodechar{ᵇ}{\ifmmode^{b}\else\textsuperscript{\ensuremath{b}}\fi}
\newunicodechar{ʳ}{\ifmmode^{r}\else\textsuperscript{\ensuremath{r}}\fi}
\newunicodechar{ᵘ}{\ifmmode^{u}\else\textsuperscript{\ensuremath{u}}\fi}
\newunicodechar{ʸ}{\ifmmode^{y}\else\textsuperscript{\ensuremath{y}}\fi}
\newunicodechar{ᵃ}{\ifmmode^{a}\else\textsuperscript{\ensuremath{a}}\fi}
\newunicodechar{ᵉ}{\ifmmode^{e}\else\textsuperscript{\ensuremath{e}}\fi}
\newunicodechar{ᵏ}{\ifmmode^{k}\else\textsuperscript{\ensuremath{k}}\fi}
\newunicodechar{ᵖ}{\ifmmode^{p}\else\textsuperscript{\ensuremath{p}}\fi}
\newunicodechar{ᴺ}{\ifmmode^{N}\else\textsuperscript{\ensuremath{N}}\fi}
\newunicodechar{ᶜ}{\ifmmode^{c}\else\textsuperscript{\ensuremath{c}}\fi}
\newunicodechar{ʰ}{\ifmmode^{h}\else\textsuperscript{\ensuremath{h}}\fi}
\newunicodechar{ᵠ}{\ifmmode^{q}\else\textsuperscript{\ensuremath{q}}\fi}
\newunicodechar{ₙ}{\ifmmode_{n}\else\textsubscript{\ensuremath{n}}\fi}
\newunicodechar{ₐ}{\ifmmode_{a}\else\textsubscript{\ensuremath{a}}\fi}
\newunicodechar{ᵢ}{\ifmmode_{i}\else\textsubscript{\ensuremath{i}}\fi}
\newunicodechar{ᵣ}{\ifmmode_{r}\else\textsubscript{\ensuremath{r}}\fi}
\newunicodechar{ₖ}{\ifmmode_{k}\else\textsubscript{\ensuremath{k}}\fi}
\newunicodechar{ᵦ}{\ifmmode_{\beta}\else\textsubscript{\ensuremath{\beta}}\fi}
\newunicodechar{ₓ}{\ifmmode_{x}\else\textsubscript{\ensuremath{x}}\fi}
\newfontfamily\popdisplay{ArchivoBlack-Regular.ttf}[Path=../../../fonts/]
\newfontfamily\poplogo{SpaceGrotesk-Bold.ttf}[Path=../../../fonts/]
\newfontfamily\popsemi{PlusJakartaSans-SemiBold.ttf}[Path=../../../fonts/]
\newfontfamily\popxbold{PlusJakartaSans-ExtraBold.ttf}[Path=../../../fonts/]
\newfontfamily\popxboldls{PlusJakartaSans-ExtraBold.ttf}[Path=../../../fonts/,LetterSpace=3.0]

\setlist[enumerate]{leftmargin=1.7em,itemsep=5pt,topsep=4pt}
\setlist[itemize]{leftmargin=1.7em,itemsep=4pt,topsep=4pt,label={\color{poporange}\textbullet}}
\setlength{\parindent}{0pt}
\setlength{\parskip}{5pt}
\renewcommand{\arraystretch}{1.25}

% pgfplots : style Pop par défaut
\pgfplotsset{
  every axis/.append style={axis line style={popink,line width=1pt},tick style={popink},
    grid style={popline,line width=0.5pt},label style={font=\small},tick label style={font=\footnotesize}},
  popcurve/.style={poporange,line width=1.8pt,no markers,smooth},
}
% Même style disponible dans un \draw TikZ simple (hors pgfplots)
\tikzset{popcurve/.style={poporange,line width=1.8pt}}
\tikzset{popgrid/.style={popline,very thin}}
\colorlet{popcurve}{poporange} % le nom du style est parfois utilisé comme couleur

% ============================================================
% Pill / squiggle
% ============================================================
\newcommand{\poppill}[3]{%
  \tikz[baseline=-0.55ex]\node[draw=popink,line width=1.2pt,rounded corners=2.6pt,%
    fill=#2,inner xsep=6.5pt,inner ysep=3pt]{%
    \popxboldls\fontsize{7}{7}\selectfont\color{#3}\MakeUppercase{#1}};%
}
\newcommand{\popsquiggle}[1]{%
  \tikz[baseline=-0.4ex]{
    \draw[#1,line width=2.6pt,line cap=round]
      plot [smooth] coordinates {(0,0.09) (0.34,0.01) (0.68,0.21) (1.02,0.01) (1.36,0.10)};
  }%
}
% Mot-clé surligné (marqueur orange sous le texte)
\newcommand{\cle}[1]{{\popxbold\color{popdark}#1}}

% ============================================================
% En-tête / pied
% ============================================================
\pagestyle{fancy}
\fancyhf{}
\renewcommand{\headrulewidth}{0pt}
\renewcommand{\footrulewidth}{0pt}
\lhead{\raisebox{-0.15ex}{\poplogo\fontsize{13}{13}\selectfont\color{popink}OnBuch\color{poporange}+}}
\rhead{\poppill{\DOCMATIERE\ \textbullet\ \DOCNIVEAU\ \textbullet\ Leçon \DOCLECON}{white}{popink}}
\lfoot{\popxbold\fontsize{7.6}{7.6}\selectfont\color{popmuted}\MakeUppercase{Cours signé OnBuch+}\ \textbullet\ {\popsemi\DOCTITRE}}
\rfoot{\tikz[baseline=-0.6ex]\node[rounded corners=8.5pt,fill=popink,minimum height=17pt,minimum width=17pt,inner xsep=5pt,inner sep=0]{\popxbold\fontsize{8}{8}\selectfont\color{popcream}\thepage\,/\,\pageref*{LastPage}};}

% ============================================================
% Titres
% ============================================================
\newcommand{\chaptitre}[2]{%
  \begin{tcolorbox}[enhanced,colback=poporange,colframe=popink,boxrule=2.6pt,arc=11pt,
      drop shadow=popink,left=15pt,right=15pt,top=12pt,bottom=11pt,before skip=3pt,after skip=18pt]
    \poppill{#2}{popink}{popcream}\par\vspace{9pt}
    {\raggedright\hyphenpenalty=10000\exhyphenpenalty=10000\popdisplay\fontsize{22}{24}\selectfont\color{popcream}\MakeUppercase{#1}\par}
  \end{tcolorbox}
}
% Section numérotée : pastille orange avec le numéro + titre Archivo
\newcounter{popsec}
\newcommand{\coursec}[1]{%
  \stepcounter{popsec}\setcounter{popsub}{0}%
  \par\vspace{16pt}\noindent
  \begin{minipage}{\linewidth}
  \tikz[baseline=-0.7ex]\node[draw=popink,line width=1.6pt,fill=poporange,rounded corners=4pt,minimum size=22pt,inner sep=0]{\popdisplay\fontsize{12}{12}\selectfont\color{popcream}\thepopsec};%
  \hspace{8pt}{\popdisplay\fontsize{15}{17}\selectfont\color{popink}\MakeUppercase{#1}}\par
  \vspace{4pt}\noindent{\color{poporange}\rule{36pt}{3.6pt}}
  \end{minipage}\par\nopagebreak\vspace{8pt}
}
\newcounter{popsub}
\newcommand{\courssub}[1]{%
  \stepcounter{popsub}%
  \par\vspace{9pt}\noindent{\popxbold\fontsize{11.5}{13}\selectfont\color{popink}{\color{poporange}\thepopsec.\thepopsub}\hspace{6pt}#1}\par\nopagebreak\vspace{4pt}
}

% ============================================================
% Boîtes pédagogiques — même recette : fond blanc, bordure épaisse colorée,
% ombre dure encre, badge plein. Titre optionnel [#1].
% ============================================================
\tcbset{popbase/.style={breakable,enhanced,colback=white,boxrule=2.3pt,arc=9pt,
  shadow={3pt}{-3pt}{0pt}{popink},left=14pt,right=14pt,top=11pt,bottom=11pt,before skip=11pt,after skip=15pt,
  fontupper=\popsemi\fontsize{10.6}{14.5}\selectfont\color{popink},
  fontlower=\popsemi\fontsize{10.6}{14.5}\selectfont\color{popink}}}
% Coche dessinée (le glyphe ✓ n'existe pas dans les polices du document)
\newcommand{\popcheck}[1]{\tikz[baseline=-0.5ex]\draw[#1,line width=1.6pt,line cap=round,line join=round](-0.13,0.0)--(-0.04,-0.09)--(0.13,0.1);}
\providecommand{\coche}{\popcheck{popgreen}}
\providecommand{\resultat}[1]{\par\smallskip\noindent\fcolorbox{popgreen}{popgreenL}{\parbox{\dimexpr\linewidth-2\fboxsep-2\fboxrule\relax}{\textbf{Résultat.} #1}}\par\smallskip}
\newcommand{\popboxhead}[3]{\poppill{#1}{#2}{white}\ifx\relax#3\relax\else\hspace{6pt}{\popxbold\fontsize{10.6}{12}\selectfont\color{popink}#3}\fi\par\vspace{7pt}}

\newtcolorbox{aretenir}[1][]{popbase,colframe=popgreen,before upper={\popboxhead{À retenir}{popgreen}{#1}}}
\newtcolorbox{exemplebox}[1][]{popbase,colframe=poporange,before upper={\popboxhead{Exemple}{poporange}{#1}}}
\newtcolorbox{definition}[1][]{popbase,colframe=popblue,before upper={\popboxhead{Définition}{popblue}{#1}}}
\newtcolorbox{propriete}[1][]{popbase,colframe=poppurple,before upper={\popboxhead{Propriété}{poppurple}{#1}}}
\newtcolorbox{methode}[1][]{popbase,colframe=popgold,before upper={\popboxhead{Méthode}{popgold}{#1}}}
\newtcolorbox{attention}[1][]{popbase,colframe=poppink,before upper={\popboxhead{Attention}{poppink}{#1}}}
\newtcolorbox{experience}[1][]{popbase,colframe=popblue,colback=popblueL,before upper={\popboxhead{Expérience}{popblue}{#1}}}
% « Le savais-tu ? » : carte violette pleine (contexte, culture, Cameroun)
\newtcolorbox{savaistu}[1][]{popbase,colback=poppurple,colframe=popink,
  fontupper=\popsemi\fontsize{10.4}{14.2}\selectfont\color{white},
  before upper={\poppill{Le savais-tu\,?}{white}{poppurple}\ifx\relax#1\relax\else\hspace{6pt}{\popxbold\color{white}#1}\fi\par\vspace{7pt}}}
% Exercice résolu : énoncé en haut, \tcblower puis la solution sur fond crème
\newcounter{popexo}\newcounter{popexr}
\newtcolorbox{exoresolu}[1][]{popbase,colframe=poporange,colbacklower=popcream,
  segmentation style={popink,line width=1.2pt,dashed},
  before upper={\stepcounter{popexr}\popboxhead{Exercice résolu \thepopexr}{poporange}{#1}},
  before lower={\poppill{Solution}{popink}{popcream}\par\vspace{6pt}}}
% Exercice (non résolu en place) — la correction est dans la partie Corrigés
\newtcolorbox{exercice}[1][]{popbase,colframe=popink,
  before upper={\stepcounter{popexo}\popboxhead{Exercice \thepopexo}{popink}{#1}}}
% Corrigé
\newtcolorbox{corrige}[1][]{popbase,colframe=popgreen,colback=popgreenL,
  before upper={\popboxhead{Corrigé}{popgreen}{#1}}}
% Objectifs / prérequis / bilan
\newtcolorbox{objectifs}{popbase,colframe=popink,colback=poporangeL,
  before upper={\popboxhead{Objectifs de la leçon}{popink}{}}}
\newtcolorbox{prerequis}{popbase,colframe=popink,colback=popblueL,
  before upper={\popboxhead{Prérequis}{popblue}{}}}
\newtcolorbox{bilan}{popbase,colframe=popink,colback=white,boxrule=2.8pt,
  before upper={\popboxhead{Fiche bilan}{poporange}{L'essentiel à connaître pour l'examen}}}
% Figure encadrée avec légende
\newtcolorbox{popfigure}[1][]{enhanced,breakable=false,colback=white,colframe=popline,boxrule=1.4pt,arc=8pt,
  left=10pt,right=10pt,top=10pt,bottom=8pt,before skip=10pt,after skip=12pt,halign=center,
  lower separated=false,halign lower=center,
  fontlower=\popsemi\fontsize{9}{11.5}\selectfont\color{popmuted},
  #1}
\newcommand{\legende}[1]{\par\vspace{6pt}{\popsemi\fontsize{9}{11.5}\selectfont\color{popmuted}#1}}

% ============================================================
% Signature OnBuch+
% ============================================================
\newcommand{\poplogoinline}[1]{{\poplogo\fontsize{#1}{#1}\selectfont\color{popink}OnBuch\color{poporange}+}}
\newcommand{\signatureonbuch}{%
  \begin{tcolorbox}[enhanced,colback=popink,colframe=popink,boxrule=2.2pt,arc=10pt,
      drop shadow=poporange,left=14pt,right=14pt,top=10pt,bottom=10pt,before skip=0pt,after skip=0pt]
    \tikz[baseline=-0.7ex]\node[circle,fill=popgreen,draw=popcream,line width=1.3pt,minimum size=22pt,inner sep=0]{\popcheck{white}};%
    \hspace{9pt}\begin{minipage}[c]{0.62\linewidth}
      {\popxbold\fontsize{12}{13}\selectfont\color{popcream}Signé\ \textbullet\ {\poplogo OnBuch\color{poporange}+}}\par\vspace{2pt}
      {\raggedright\popsemi\fontsize{8}{10.5}\selectfont\color{popline}\MakeUppercase{Équipe pédagogique OnBuch+}\\\MakeUppercase{Conforme au programme officiel MINESEC}\par}
    \end{minipage}\hfill
    {\poplogo\fontsize{22}{22}\selectfont\color{popcream}OnBuch\color{poporange}+}
  \end{tcolorbox}%
}

% ============================================================
% Page de garde
% #1 titre · #2 matière · #3 niveau · #4 module · #5 n° leçon · #6 durée ·
% #7 description / objectifs (texte ou itemize)
% ============================================================
\newcommand{\pagegarde}[7]{%
  \thispagestyle{empty}
  \newgeometry{a4paper,margin=1.7cm,top=1.6cm,bottom=1.4cm}%
  \begin{tikzpicture}[remember picture,overlay]
    \fill[poporange!18] ([xshift=-3.2cm,yshift=-3.6cm]current page.north east) circle (4.6cm);
    \fill[poppurple!14] ([xshift=2.4cm,yshift=7.6cm]current page.south west) circle (3.8cm);
  \end{tikzpicture}%
  {\poplogo\fontsize{30}{30}\selectfont\color{popink}OnBuch\color{poporange}+}\hfill
  \raisebox{0.6ex}{\poppill{Cours signé \textbullet\ Premium}{popink}{popcream}}\par
  \vspace{1pt}\popsquiggle{poppurple}\par
  \vspace{8pt}
  \begin{tcolorbox}[enhanced,colback=poporange,colframe=popink,boxrule=2.8pt,arc=13pt,
      drop shadow=popink,left=18pt,right=18pt,top=12pt,bottom=13pt,after skip=10pt]
    \poppill{#2 \textbullet\ #3}{popink}{popcream}\hspace{5pt}\poppill{#4}{white}{popink}\par\vspace{8pt}
    {\popdisplay\fontsize{14}{14}\selectfont\color{popink}LEÇON #5}\par\vspace{4pt}
    {\raggedright\hyphenpenalty=10000\exhyphenpenalty=10000\popdisplay\fontsize{25}{27}\selectfont\color{popcream}\MakeUppercase{#1}\par}
    \vspace{8pt}
    {\popxbold\fontsize{9.5}{9.5}\selectfont\color{popink}\MakeUppercase{Durée conseillée : #6}}
  \end{tcolorbox}
  \begin{tcolorbox}[enhanced,colback=white,colframe=popink,boxrule=2.2pt,arc=10pt,
      drop shadow=popink,left=16pt,right=16pt,top=10pt,bottom=10pt,after skip=8pt]
    \poppill{Dans cette leçon}{poppurple}{white}\par\vspace{5pt}
    \popsemi\fontsize{9.3}{12.6}\selectfont\color{popink}#7
  \end{tcolorbox}
  \noindent\begin{minipage}[t]{0.487\linewidth}
    \begin{tcolorbox}[enhanced,colback=popgreen,colframe=popink,boxrule=2.2pt,arc=10pt,
        drop shadow=popink,left=11pt,right=11pt,top=10pt,bottom=10pt,height=3.1cm,valign=center]
      \tcbox[colback=white,colframe=popink,boxrule=1.3pt,arc=5pt,boxsep=0pt,nobeforeafter,
        left=3pt,right=3pt,top=3pt,bottom=3pt,valign=center]{\qrcode[height=1.9cm]{https://play.google.com/store/apps/details?id=com.luvvix.onbuch&hl=fr}}%
      \hspace{8pt}\begin{minipage}[c]{0.5\linewidth}
        {\popdisplay\fontsize{11}{12.5}\selectfont\color{white}\MakeUppercase{Télécharge OnBuch}}\par\vspace{4pt}
        {\raggedright\popsemi\fontsize{8.4}{11}\selectfont\color{white}Gratuit \textbullet\ Google Play\\Tuteur IA, annales, concours, résultats\par}
      \end{minipage}
    \end{tcolorbox}
  \end{minipage}\hfill
  \begin{minipage}[t]{0.487\linewidth}
    \begin{tcolorbox}[enhanced,colback=poppurple,colframe=popink,boxrule=2.2pt,arc=10pt,
        drop shadow=popink,left=11pt,right=11pt,top=10pt,bottom=10pt,height=3.1cm,valign=center]
      \tikz[baseline=-0.7ex]\node[circle,fill=white,draw=popink,line width=1.3pt,minimum size=1.6cm,inner sep=0]{\popdisplay\fontsize{24}{24}\selectfont\color{poppurple}+};%
      \hspace{8pt}\begin{minipage}[c]{0.6\linewidth}
        {\popdisplay\fontsize{11}{12.5}\selectfont\color{white}\MakeUppercase{Passe à OnBuch+}}\par\vspace{4pt}
        {\raggedright\popsemi\fontsize{8.4}{11}\selectfont\color{white}Tous les cours signés, Tuteur IA illimité, fiches PDF — onglet OnBuch+ de l'app\par}
      \end{minipage}
    \end{tcolorbox}
  \end{minipage}
  \vspace{8pt}
  \popcontenu
  \vfill
  \signatureonbuch
  \restoregeometry
  \newpage
}

% Rangée « Ce cours contient » (page de garde)
\newcommand{\poptile}[3]{% #1 couleur #2 gros texte #3 légende
  \begin{tcolorbox}[enhanced,colback=#1,colframe=popink,boxrule=2pt,arc=9pt,shadow={2.5pt}{-2.5pt}{0pt}{popink},
      left=6pt,right=6pt,top=9pt,bottom=9pt,halign=center,valign=center,height=1.75cm,width=\linewidth,nobeforeafter]
    {\popdisplay\fontsize{12.5}{13}\selectfont\color{white}\MakeUppercase{#2}}\par\vspace{3pt}
    {\centering\popsemi\fontsize{7.8}{9.5}\selectfont\color{white}#3\par}
  \end{tcolorbox}}
\newcommand{\popcontenu}{%
  \par\noindent{\popxboldls\fontsize{7.5}{7.5}\selectfont\color{popmuted}CE COURS SIGNÉ CONTIENT}\par\vspace{7pt}
  \noindent\begin{minipage}[t]{0.235\linewidth}\poptile{popblue}{Cours}{complet, illustré, pas à pas}\end{minipage}\hfill
  \begin{minipage}[t]{0.235\linewidth}\poptile{poporange}{Méthodes}{et exercices résolus}\end{minipage}\hfill
  \begin{minipage}[t]{0.235\linewidth}\poptile{poppink}{Exercices}{type examen + corrigés}\end{minipage}\hfill
  \begin{minipage}[t]{0.235\linewidth}\poptile{popgreen}{Fiche bilan}{l'essentiel à retenir}\end{minipage}\par}

% Sommaire
\newtcolorbox{sommaire}{enhanced,colback=white,colframe=popink,boxrule=2.2pt,arc=10pt,
  drop shadow=popink,left=18pt,right=18pt,top=13pt,bottom=13pt,before skip=6pt,after skip=20pt,
  before upper={\poppill{Sommaire}{poppurple}{white}\par\vspace{9pt}\popsemi\fontsize{10.6}{14.5}\selectfont\color{popink}}}

% Signature de fin de cours — poussée tout en bas de la dernière page
\newcommand{\findecours}{%
  \par\vfill\nopagebreak
  \begin{center}\popsquiggle{poporange}\end{center}\vspace{4pt}
  \signatureonbuch
}
\AtBeginDocument{\def\llbracket{\mathopen{[\mkern-3.5mu[}}\def\rrbracket{\mathclose{]\mkern-3.5mu]}}} % intervalles d'entiers (glyphe ⟦ absent de la police)
% Glyphes absents de Fira Math : redéfinis à partir de glyphes disponibles
\AtBeginDocument{%
  \def\setminus{\mathbin{\backslash}}\let\smallsetminus\setminus
  \def\vdots{\mathord{\vbox{\baselineskip3.2pt\lineskiplimit0pt\kern4pt\hbox{.}\hbox{.}\hbox{.}}}}%
  \def\ddots{\mathinner{\mkern1mu\raise6.5pt\hbox{.}\mkern2mu\raise3.6pt\hbox{.}\mkern2mu\raise0.7pt\hbox{.}\mkern1mu}}%
  \def\triangle{\mathord{\scalebox{0.9}{$\symup{\Delta}$}}}%
  \def\nabla{\mathord{\raisebox{\depth}{\scalebox{1}[-1]{$\symup{\Delta}$}}}}%
  \def\checkmark{\popcheck{popdark}}%
  \def\star{\mathbin{\ast}}\def\bigstar{\mathord{\ast}}%
  \def\diamond{\mathbin{\scalebox{0.6}{$\Diamond$}}}%
}

%% FIGFIX-BEGIN v3 — correctifs globaux des figures (docs/onbuch-generation/tools/figfix)
\usepackage{adjustbox}\usepackage{etoolbox}
% toute figure TikZ/pgfplots plus large que la ligne est réduite à la largeur disponible
\BeforeBeginEnvironment{tikzpicture}{\begin{adjustbox}{max width=\linewidth}}
\AfterEndEnvironment{tikzpicture}{\end{adjustbox}}
% légendes pgfplots lisibles par-dessus les courbes
\pgfplotsset{every axis/.append style={legend style={fill=white,fill opacity=0.92,text opacity=1,draw=black!15,inner sep=3pt}}}
% étiquettes d'arêtes (syntaxe edge["..."]) avec fond blanc
\tikzset{every edge quotes/.append style={fill=white,inner sep=1.2pt}}
% bulles de cartes mentales : texte à la ligne dans la bulle, bulle un peu plus grande
\tikzset{every concept/.append style={text width=2.0cm,align=center,minimum size=2.3cm,inner sep=2pt}}
%% FIGFIX-END
\begin{document}

\pagegarde{Niveaux d'énergie de l'atome d'hydrogène}{Physique}{Tle E}{Module 4 : Onde, matière et transformations nucléaires}{P17}{3 h}{Cette leçon explique pourquoi l'énergie de l'électron dans l'atome d'hydrogène est quantifiée et comment les transitions entre niveaux d'énergie rendent compte des spectres d'émission et d'absorption. Elle fournit les outils (E = hν, En = -E0/n²) pour calculer les longueurs d'onde des raies des séries de Lyman, Balmer et Paschen.
\vspace{4pt}
\begin{multicols}{2}\small
\begin{itemize}
\item À la fin de cette leçon, je sais définir un photon et appliquer la relation E = hν = hc/λ
\item À la fin de cette leçon, je sais énoncer la quantification de l'énergie de l'atome d'hydrogène En = -E0/n² et identifier l'état fondamental, les états excités et l'ionisation
\item À la fin de cette leçon, je sais construire et lire un diagramme énergétique de l'atome d'hydrogène
\item À la fin de cette leçon, je sais expliquer l'origine des raies des spectres d'émission et d'absorption
\item À la fin de cette leçon, je sais calculer l'énergie et la longueur d'onde du photon émis ou absorbé lors d'une transition entre deux niveaux
\item À la fin de cette leçon, je sais rattacher une raie à sa série spectrale (Lyman, Balmer, Paschen) et à son domaine (UV, visible, IR)
\end{itemize}\end{multicols}}

\begin{sommaire}
\begin{multicols}{2}
\begin{enumerate}
\item Le photon et la quantification de l'énergie lumineuse
\item Niveaux d'énergie de l'atome d'hydrogène
\item Transitions entre niveaux : émission et absorption de photons
\item Spectres d'émission et d'absorption
\item Séries spectrales de l'hydrogène : Lyman, Balmer, Paschen
\item Méthodes de résolution et applications type Bac
\item Activité d'intégration
\item Méthodes et exercices résolus
\item Exercices
\item Corrigés des exercices
\item Fiche bilan
\end{enumerate}
\end{multicols}
\end{sommaire}

\begin{objectifs}
\begin{itemize}
\item À la fin de cette leçon, je sais définir un photon et appliquer la relation E = hν = hc/λ
\item À la fin de cette leçon, je sais énoncer la quantification de l'énergie de l'atome d'hydrogène En = -E0/n² et identifier l'état fondamental, les états excités et l'ionisation
\item À la fin de cette leçon, je sais construire et lire un diagramme énergétique de l'atome d'hydrogène
\item À la fin de cette leçon, je sais expliquer l'origine des raies des spectres d'émission et d'absorption
\item À la fin de cette leçon, je sais calculer l'énergie et la longueur d'onde du photon émis ou absorbé lors d'une transition entre deux niveaux
\item À la fin de cette leçon, je sais rattacher une raie à sa série spectrale (Lyman, Balmer, Paschen) et à son domaine (UV, visible, IR)
\item À la fin de cette leçon, je sais convertir correctement les énergies entre joules et électronvolts
\item À la fin de cette leçon, je sais interpréter un spectre pour identifier la présence d'un élément chimique
\end{itemize}
\end{objectifs}

\begin{prerequis}
\begin{itemize}
\item Onde lumineuse : célérité c = 3,00×10⁸ m/s, relation c = λν, domaines du spectre électromagnétique (Première)
\item Énergie potentielle électrique et interaction coulombienne (notion qualitative)
\item Conversions d'unités : nm ↔ m, eV ↔ J (1 eV = 1,60×10⁻¹⁹ J)
\item Puissances de dix et calcul littéral (expressions en 1/n²)
\item Notion de spectre obtenu par dispersion (prisme, réseau) vue en Première
\end{itemize}
\end{prerequis}

\newpage

\coursec{Le photon et la quantification de l'énergie lumineuse}

\courssub{Situation d'accroche : pourquoi chaque lumière a-t-elle sa couleur ?}

Le soir tombe sur Douala. Les enseignes au néon s'allument : rouge-orangé, éclatant, immédiatement reconnaissable. À Yaoundé, les lampadaires à vapeur de sodium baignent les rues d'une lumière jaune caractéristique, bien différente de la lumière blanche des lampes à incandescence. Pourquoi le néon émet-il du rouge-orangé et le sodium du jaune ? Pourquoi pas l'inverse ?

Autre énigme : lorsqu'on décompose la lumière du Soleil à travers un prisme de très bonne qualité, on observe le fameux arc-en-ciel... mais traversé de fines \cle{raies noires}, comme si certaines couleurs précises avaient été « volées » en chemin. Ces raies, découvertes par Fraunhofer au début du XIX\textsuperscript{e} siècle, portent une signature : celle des atomes présents dans l'atmosphère solaire. C'est d'ailleurs ainsi que les astrophysiciens déterminent la composition chimique des étoiles situées à des milliards de kilomètres, sans jamais quitter la Terre : ils lisent la « carte d'identité lumineuse » gravée dans leur spectre.

Toutes ces observations ont un point commun : elles sont \emph{inexplicables} si l'on considère la lumière comme une simple onde continue. Il a fallu une révolution conceptuelle, menée par Planck puis Einstein entre 1900 et 1905, pour comprendre que la lumière échange son énergie avec la matière \textbf{par paquets entiers} : les \cle{photons}.

\textbf{Problématique.} Comment l'énergie lumineuse est-elle transportée et échangée avec la matière ? Pourquoi les atomes n'émettent et n'absorbent-ils que certaines couleurs bien précises ?

\courssub{Rappel : la lumière, une onde électromagnétique}

Depuis les travaux de Maxwell et les expériences de Hertz, on sait que la lumière est une \cle{onde électromagnétique} : un champ électrique $\vec{E}$ et un champ magnétique $\vec{B}$ qui se propagent en vibrant, sans support matériel, y compris dans le vide.

\begin{aretenir}[Célérité et relation fondamentale]
Dans le vide, toute onde électromagnétique se propage à la célérité
\[ c = \num{3,00e8}~\text{m/s}. \]
Sa fréquence $\nu$ (en hertz, Hz) et sa longueur d'onde $\lambda$ (en mètres, m) sont liées par
\[ c = \lambda \, \nu. \]
Le domaine visible s'étend approximativement de $\lambda \approx 400$~nm (violet) à $\lambda \approx 800$~nm (rouge). En deçà se trouve l'\textbf{ultraviolet} (UV), au-delà l'\textbf{infrarouge} (IR).
\end{aretenir}

\begin{attention}[Fréquence et longueur d'onde : ne pas les confondre]
La relation $c = \lambda \nu$ montre que $\lambda$ et $\nu$ varient en sens inverse : une \textbf{grande} longueur d'onde correspond à une \textbf{petite} fréquence, et réciproquement. Le rouge ($\lambda \approx 700$~nm) a une fréquence plus faible que le violet ($\lambda \approx 400$~nm). Cette remarque aura des conséquences énergétiques capitales dans la suite.
\end{attention}

\courssub{Les limites du modèle ondulatoire}

Le modèle ondulatoire explique magnifiquement les interférences, la diffraction, la réfraction. Mais à la fin du XIX\textsuperscript{e} siècle, deux phénomènes lui résistent totalement.

\textbf{1. L'effet photoélectrique.} Lorsqu'on éclaire une plaque de zinc avec de la lumière ultraviolette, des électrons sont arrachés du métal. Or l'expérience montre des faits déroutants :
\begin{itemize}
\item en dessous d'une certaine fréquence $\nu_0$ (fréquence seuil), \emph{aucun} électron n'est éjecté, même avec une lumière très intense ;
\item au-dessus de $\nu_0$, l'extraction est \emph{instantanée}, même avec une lumière très faible ;
\item augmenter l'intensité augmente le \emph{nombre} d'électrons éjectés, mais pas leur énergie individuelle.
\end{itemize}
Si la lumière était une onde continue, une lumière rouge très intense devrait, en accumulant de l'énergie, finir par arracher des électrons. Ce n'est jamais le cas. L'énergie semble arriver par « coups » dont l'efficacité dépend de la fréquence, pas de l'intensité.

\textbf{2. Les spectres de raies.} Un gaz d'hydrogène porté à haute température (ou traversé par une décharge électrique) n'émet pas toutes les couleurs : il émet quelques raies fines et précises (une rouge, une bleu-vert, deux violettes dans le visible...). Une onde continue devrait pouvoir transporter n'importe quelle fréquence. D'où viennent ces fréquences privilégiées ?

\begin{savaistu}[Le prix Nobel d'Einstein... pas celui qu'on croit]
Albert Einstein est célèbre pour la relativité. Pourtant, le prix Nobel de physique qu'il reçut en 1921 récompensait son explication de \textbf{l'effet photoélectrique} (1905), fondée sur l'hypothèse des quanta de lumière. C'est ce travail, plus que $E = mc^2$, qui a convaincu le jury : il ouvrait la porte à la physique quantique, la théorie qui gouverne aujourd'hui l'électronique de ton téléphone, les lasers des chirurgiens et les panneaux solaires de l'électrification rurale camerounaise.
\end{savaistu}

\courssub{L'hypothèse de Planck--Einstein : le photon}

En 1900, Max Planck, pour expliquer le rayonnement des corps chauds, fait une hypothèse audacieuse : les échanges d'énergie entre la matière et le rayonnement ne peuvent se faire que par \textbf{paquets} d'énergie proportionnels à la fréquence. En 1905, Einstein radicalise l'idée : ce n'est pas seulement l'échange qui est discontinu, c'est la lumière elle-même qui est constituée de « grains » d'énergie, qu'on appellera plus tard les \cle{photons}.

\begin{definition}[Le photon]
Le \cle{photon} est le \textbf{quantum} (grain élémentaire, indivisible) d'énergie lumineuse. C'est une particule :
\begin{itemize}
\item de \textbf{masse nulle} ;
\item de \textbf{charge nulle} ;
\item se déplaçant dans le vide à la célérité $c = \num{3,00e8}$~m/s ;
\item transportant une énergie qui dépend de la fréquence $\nu$ de la radiation associée.
\end{itemize}
\end{definition}

\begin{propriete}[Relation de Planck--Einstein — admise]
L'énergie d'un photon associé à une radiation de fréquence $\nu$ et de longueur d'onde $\lambda$ est
\[ E = h\nu = \frac{hc}{\lambda}, \]
où $h = \num{6,63e-34}$~J·s est la \textbf{constante de Planck}.

Cette relation est un \emph{postulat} de la physique quantique : elle ne se démontre pas dans le cadre de la physique classique. Elle est \textbf{admise} au programme, mais elle est confirmée avec une précision remarquable par l'expérience, notamment par l'effet photoélectrique dont elle rend compte quantitativement.
\end{propriete}

Vérifions l'homogénéité de la relation $E = hc/\lambda$ : $[h] = \text{J·s}$, $[c] = \text{m·s}^{-1}$, $[\lambda] = \text{m}$, donc
\[ \left[\frac{hc}{\lambda}\right] = \frac{\text{J·s} \times \text{m·s}^{-1}}{\text{m}} = \text{J}. \]
On obtient bien une énergie.

\begin{attention}[Le sens de la proportionnalité]
Comme $E = hc/\lambda$, l'énergie du photon est \textbf{inversement proportionnelle} à la longueur d'onde : plus $\lambda$ est grande, plus le photon est \emph{peu} énergétique. Ainsi :
\[ E_{\text{UV}} > E_{\text{visible}} > E_{\text{IR}}. \]
C'est pourquoi les UV (photons énergétiques) provoquent des coups de soleil et peuvent ioniser des molécules, tandis que les IR (photons peu énergétiques) ne font que chauffer la peau. Erreur classique au Bac : affirmer que « le rouge est plus énergétique que le violet » — c'est l'inverse !
\end{attention}

\begin{popfigure}
\begin{center}
\begin{tikzpicture}[x=1cm,y=1cm,>=Stealth]
% Axe des longueurs d'onde
\draw[->,thick,popdark] (0,0) -- (13.5,0) node[below left=2pt and -30pt,font=\small] {$\lambda$ croissant};
% Graduations
\foreach \x/\lab in {1/{10 nm},3.2/{400 nm},6.4/{800 nm},12/{1 mm}}{
  \draw[thick,popdark] (\x,0) -- (\x,0.18);
  \node[below,font=\footnotesize,popdark] at (\x,-0.05) {\lab};
}
% Zone UV
\fill[poppurple!35] (0.2,0.25) rectangle (3.2,1.15);
\node[font=\small\bfseries,popdark] at (1.7,0.7) {UV};
% Zone visible : dégradé simulé par bandes colorées
\fill[violet!70] (3.2,0.25) rectangle (4.0,1.15);
\fill[blue!70] (4.0,0.25) rectangle (4.8,1.15);
\fill[popgreen!70] (4.8,0.25) rectangle (5.6,1.15);
\fill[orange!80] (5.6,0.25) rectangle (6.0,1.15);
\fill[red!80] (6.0,0.25) rectangle (6.4,1.15);
\node[font=\small\bfseries,popdark] at (4.8,1.45) {Visible};
% Zone IR
\fill[poporange!30] (6.4,0.25) rectangle (12,1.15);
\node[font=\small\bfseries,popdark] at (9.2,0.7) {IR};
% Énergies
\node[font=\footnotesize,align=center,poppurple!80!black] at (1.7,2.3) {$E \approx 10$ à $100$ eV\\ photons énergétiques};
\draw[->,poppurple!80!black] (1.7,1.95) -- (1.7,1.2);
\node[font=\footnotesize,align=center,popdark] at (4.8,2.3) {$E \approx 1{,}6$ à $3{,}1$ eV};
\draw[->,popdark] (4.8,1.95) -- (4.8,1.2);
\node[font=\footnotesize,align=center,poporange!80!black] at (9.2,2.3) {$E < 1{,}6$ eV\\ photons peu énergétiques};
\draw[->,poporange!80!black] (9.2,1.95) -- (9.2,1.2);
% Flèche énergie décroissante
\draw[->,very thick,popink] (1,3.1) -- (11.5,3.1) node[midway,above,font=\small\bfseries,popink] {énergie du photon $E = hc/\lambda$ décroissante};
\end{tikzpicture}
\end{center}
\legende{Spectre électromagnétique autour du visible : plus la longueur d'onde augmente, plus l'énergie des photons diminue.}
\end{popfigure}

\courssub{Unité pratique : l'électron-volt}

Les énergies des photons sont minuscules à l'échelle du joule. On utilise donc une unité adaptée au monde atomique : l'\cle{électron-volt}.

\begin{definition}[L'électron-volt]
L'électron-volt (eV) est l'énergie acquise par un électron accéléré sous une tension de $1$~V :
\[ 1~\text{eV} = e \times 1~\text{V} = \num{1,60e-19}~\text{J}. \]
\end{definition}

\begin{methode}[Calculer rapidement l'énergie d'un photon en eV]
Deux chemins possibles, au choix :
\begin{enumerate}
\item \textbf{En joules (méthode SI)} : convertir d'abord $\lambda$ en \textbf{mètres} ($1~\text{nm} = 10^{-9}$~m), puis appliquer
\[ E = \frac{hc}{\lambda} \quad \text{avec} \quad h = \num{6,63e-34}~\text{J·s}, \quad c = \num{3,00e8}~\text{m/s}, \]
puis convertir en eV en divisant par $\num{1,60e-19}$.
\item \textbf{Formule directe en eV} : en combinant $E = hc/\lambda$ avec la conversion, on obtient la relation pratique
\[ E(\text{eV}) = \frac{1240}{\lambda(\text{nm})}. \]
Démonstration : $\dfrac{hc}{e} = \dfrac{\num{6,63e-34} \times \num{3,00e8}}{\num{1,60e-19}} \approx \num{1,24e-6}$~J·m/C $= 1240$~eV·nm.
\end{enumerate}
\textbf{Piège rédactionnel} : si tu utilises $E = hc/\lambda$ avec $\lambda$ laissé en nanomètres, le résultat est faux d'un facteur $10^{9}$. Toujours convertir !
\end{methode}

\begin{exemplebox}[Énergie d'un photon rouge : la raie $H_\alpha$ de l'hydrogène]
La raie rouge caractéristique de l'hydrogène a pour longueur d'onde $\lambda = 656$~nm. Calculons l'énergie du photon correspondant.

\textbf{En joules} : $\lambda = 656 \times 10^{-9}$~m, donc
\begin{align*}
E &= \frac{hc}{\lambda} = \frac{\num{6,63e-34} \times \num{3,00e8}}{656 \times 10^{-9}} \\
&= \frac{\num{1,989e-25}}{\num{6,56e-7}} \approx \num{3,03e-19}~\text{J}.
\end{align*}
\textbf{En eV} : $E = \dfrac{\num{3,03e-19}}{\num{1,60e-19}} \approx 1,89$~eV. Vérification par la formule directe : $E = 1240/656 \approx 1,89$~eV. Cohérent !
\end{exemplebox}

\begin{popfigure}
\begin{center}
\begin{tikzpicture}[>=Stealth,scale=1]
% Photon incident : onde sinusoïdale
\draw[popink,very thick,domain=0:3.6,samples=100] plot (\x,{0.35*sin(360*\x/1.2)});
\draw[->,popink,very thick] (3.6,0) -- (4.6,0);
\node[popink,font=\small\bfseries] at (1.8,0.85) {photon d'énergie $E = h\nu$};
% Atome : noyau + orbite
\draw[popblue,thick] (6.3,0) circle (1.1);
\fill[poporange] (6.3,0) circle (0.16);
\node[font=\footnotesize\bfseries,white] at (6.3,0) {n};
\fill[popblue] (6.3,1.1) circle (0.09);
\node[font=\footnotesize,popblue] at (7.6,1.15) {électron};
\node[font=\footnotesize,popdark] at (6.3,-1.55) {atome};
% Flèche d'interaction
\draw[->,popgreen,very thick] (6.3,0.25) -- (6.3,0.95);
\node[font=\footnotesize,popgreen!60!black,align=center] at (8.6,-0.5) {absorption possible si\\ $h\nu = $ écart entre deux niveaux};
\end{tikzpicture}
\end{center}
\legende{Un photon d'énergie $h\nu$ frappe un atome. Il ne peut être absorbé que si son énergie correspond exactement à un écart entre deux niveaux d'énergie de l'atome (voir sections suivantes).}
\end{popfigure}

\courssub{Ordres de grandeur et conséquence fondamentale}

Avec la formule $E(\text{eV}) = 1240/\lambda(\text{nm})$, dressons un tableau de référence à connaître :

\begin{center}
\begin{tabularx}{\linewidth}{|l|X|X|X|}\hline
\rowcolor{popblueL} \textbf{Radiation} & \textbf{$\lambda$ typique} & \textbf{$E$ (eV)} & \textbf{Effet sur la matière} \\ \hline
Ultraviolet & 300 nm & $\approx 4{,}1$ eV & coups de soleil, ionisations \\ \hline
Violet & 400 nm & $\approx 3{,}1$ eV & limite du visible \\ \hline
Rouge & 700 nm & $\approx 1{,}8$ eV & limite du visible \\ \hline
Infrarouge & \num{1000} nm & $\approx 1{,}2$ eV & échauffement (thermique) \\ \hline
\end{tabularx}
\end{center}

\begin{aretenir}[Conséquence capitale : la quantification des échanges]
Un échange d'énergie entre la lumière et la matière se fait \textbf{par paquets entiers} : un atome absorbe ou émet \emph{un photon à la fois}, jamais une fraction de photon. L'énergie échangée est donc \cle{quantifiée} : elle ne peut prendre que des valeurs multiples entiers de $h\nu$. C'est ce verrou qui expliquera, dans les sections suivantes, pourquoi les spectres atomiques sont constitués de raies discrètes et non de bandes continues.
\end{aretenir}

\begin{exoresolu}[Photons émis par une lampe à vapeur de sodium]
Les lampadaires à vapeur de sodium de Yaoundé émettent une lumière jaune de longueur d'onde $\lambda = 589$~nm. Une lampe délivre une puissance lumineuse $P = 20$~W.
\begin{enumerate}
\item Calculer l'énergie d'un photon émis, en joules puis en eV.
\item Combien de photons la lampe émet-elle par seconde ?
\end{enumerate}
\tcblower
\textbf{1. Énergie d'un photon.}
\[ E = \frac{hc}{\lambda} = \frac{\num{6,63e-34} \times \num{3,00e8}}{589 \times 10^{-9}} \approx \num{3,38e-19}~\text{J}. \]
En eV : $E = \dfrac{\num{3,38e-19}}{\num{1,60e-19}} \approx 2,11$~eV (vérification : $1240/589 \approx 2,11$~eV).

\textbf{2. Nombre de photons par seconde.} La puissance est l'énergie délivrée par unité de temps : en une seconde, la lampe fournit $E_{\text{tot}} = 20$~J. Comme chaque photon emporte un paquet entier d'énergie $E$ :
\[ N = \frac{E_{\text{tot}}}{E} = \frac{20}{\num{3,38e-19}} \approx \num{5,9e19}~\text{photons par seconde}. \]
Ce nombre colossal explique pourquoi, à notre échelle, la lumière semble continue : les « grains » sont si nombreux et si petits que la granularité nous échappe, comme les grains de sable d'une plage de Kribi vue de loin.
\end{exoresolu}

\begin{exercice}[À toi de jouer]
Un laser utilisé dans un cabinet de dermatologie à Douala émet dans l'ultraviolet à $\lambda = 308$~nm.
\begin{enumerate}
\item Calculer l'énergie d'un photon en joules et en eV.
\item Comparer avec l'énergie d'un photon rouge ($\lambda = 656$~nm). Conclure sur la dangerosité relative des UV.
\end{enumerate}
\end{exercice}

\begin{corrige}[Exercice : laser UV]
\textbf{1.} $E = \dfrac{hc}{\lambda} = \dfrac{\num{6,63e-34} \times \num{3,00e8}}{308 \times 10^{-9}} \approx \num{6,46e-19}$~J, soit $E = 1240/308 \approx 4,03$~eV.

\textbf{2.} Le photon UV ($4,03$~eV) est environ $2,1$ fois plus énergétique que le photon rouge ($1,89$~eV). Chaque photon UV transporte assez d'énergie pour rompre des liaisons moléculaires ou ioniser des atomes : c'est pourquoi les UV sont dangereux pour la peau alors que la lumière rouge ne l'est pas, à intensité égale.
\end{corrige}

\begin{aretenir}[L'essentiel de la section]
\begin{itemize}
\item La lumière est une onde électromagnétique : $c = \lambda\nu$, avec $c = \num{3,00e8}$~m/s.
\item L'effet photoélectrique et les spectres de raies imposent une description corpusculaire : la lumière est constituée de \cle{photons}, particules de masse nulle se déplaçant à $c$.
\item Énergie d'un photon : $E = h\nu = \dfrac{hc}{\lambda}$, avec $h = \num{6,63e-34}$~J·s ; relation pratique $E(\text{eV}) = 1240/\lambda(\text{nm})$.
\item Plus $\lambda$ est grande, moins le photon est énergétique : $E_{\text{UV}} > E_{\text{visible}} > E_{\text{IR}}$.
\item Les échanges lumière--matière se font par paquets entiers : l'énergie est \textbf{quantifiée}.
\end{itemize}
\end{aretenir}

\coursec{Niveaux d'énergie de l'atome d'hydrogène}

Dans la section précédente, nous avons vu que la lumière n'est pas une onde continue du point de vue énergétique : elle est constituée de grains d'énergie appelés \cle{photons}, chacun transportant l'énergie $E = h\nu = \dfrac{hc}{\lambda}$. Mais une question fondamentale demeure : \emph{pourquoi} les atomes émettent-ils ou absorbent-ils des photons de certaines longueurs d'onde seulement, et pas d'autres ? Pourquoi l'hydrogène produit-il un spectre de raies discontinu, et non un arc-en-ciel complet ? La réponse, qui a révolutionné la physique au début du \textsc{xx}\textsuperscript{e} siècle, tient en un mot : l'énergie de l'électron dans l'atome est \cle{quantifiée}. C'est ce que nous allons découvrir dans cette section, en nous concentrant sur l'atome le plus simple de la nature : l'hydrogène.

\courssub{Structure de l'atome d'hydrogène : un système lié}

\textbf{L'atome d'hydrogène, le plus simple des atomes}\par

L'atome d'hydrogène est constitué de deux particules seulement :
\begin{itemize}
\item un \cle{noyau} réduit à un unique \cle{proton}, de charge $+e$ (avec $e = \num{1,60e-19}$ C) et de masse $m_p \approx \num{1,67e-27}$ kg ;
\item un \cle{électron}, de charge $-e$ et de masse $m_e \approx \num{9,11e-31}$ kg, soit environ 1836 fois plus léger que le proton.
\end{itemize}

Ces deux particules, de charges opposées, s'attirent par l'\cle{interaction électrostatique} (force de Coulomb) :
\[
F = \dfrac{1}{4\pi\varepsilon_0}\,\dfrac{e^2}{r^2}
\]
où $r$ est la distance entre le proton et l'électron. Cette force attractive joue exactement le rôle que joue la gravitation pour le système Terre--Lune : elle maintient l'électron \emph{lié} au noyau. On dit que l'électron et le proton forment un \cle{système lié} : pour séparer l'électron du noyau, il faut fournir de l'énergie.

\begin{attention}[Un modèle qui a ses limites]
L'image de l'électron « tournant » autour du noyau comme une planète autour du Soleil (modèle de Rutherford, 1911) est utile pour l'intuition, mais elle est \emph{fondamentalement fausse} en physique classique : d'après l'électromagnétisme, une charge accélérée rayonne de l'énergie ; l'électron devrait donc perdre de l'énergie en spirale et s'écraser sur le noyau en une fraction de microseconde ! L'atome classique est instable... alors que l'atome d'hydrogène existe bel et bien. C'est cette contradiction qui a conduit Niels Bohr (1913), puis la mécanique quantique, à introduire l'idée de \emph{niveaux d'énergie}.
\end{attention}

\courssub{Le constat expérimental : l'énergie est quantifiée}

\textbf{Une observation qui dérange}\par

Lorsqu'on analyse la lumière émise par une lampe à vapeur d'hydrogène (expérience réalisable au laboratoire du lycée avec une lampe spectrale et un prisme ou un réseau), on n'observe pas un spectre continu, mais quelques \emph{raies colorées} isolées : une raie rouge, une raie bleu-vert, deux raies violettes... Chaque raie correspond à un photon d'énergie bien précise $E = \dfrac{hc}{\lambda}$.

Or, si l'électron pouvait posséder \emph{n'importe quelle} valeur d'énergie, il pourrait perdre n'importe quelle quantité d'énergie, et émettre des photons de \emph{toutes} les longueurs d'onde : le spectre serait continu. Le fait qu'il ne l'est pas impose une conclusion révolutionnaire :

\begin{aretenir}[Quantification de l'énergie]
Dans un atome, l'électron ne peut posséder que \cle{certaines valeurs d'énergie bien déterminées}, appelées \cle{niveaux d'énergie}. On dit que l'énergie de l'électron est \cle{quantifiée}. Toute valeur intermédiaire entre deux niveaux est \emph{interdite}.
\end{aretenir}

\begin{savaistu}[Une analogie pour bien comprendre]
Imagine un escalier : tu peux te tenir sur la première marche, la deuxième, la troisième... mais jamais « entre deux marches ». L'énergie de l'électron dans l'atome fonctionne de la même façon : les niveaux d'énergie sont les « marches » de l'escalier atomique. Une rampe continue, elle, correspondrait au monde classique où toutes les énergies sont permises. C'est Max Planck (1900) qui a introduit le premier l'idée de « quantum » d'énergie, Einstein qui l'a appliquée à la lumière (1905), et Bohr qui l'a appliquée à l'atome (1913) — trois révolutions en treize ans !
\end{savaistu}

\courssub{La formule des niveaux d'énergie}

\textbf{Énoncé de la formule}\par

Pour l'atome d'hydrogène, l'expérience et la théorie quantique conduisent à un résultat d'une remarquable simplicité : les niveaux d'énergie ne dépendent que d'un seul nombre entier $n$, appelé \cle{nombre quantique principal}.

\begin{definition}[Niveaux d'énergie de l'atome d'hydrogène]
Un \cle{niveau d'énergie} est une valeur d'énergie permise pour l'électron de l'atome. L'ensemble des niveaux d'énergie de l'atome d'hydrogène est donné par :
\[
\boxed{E_n = -\dfrac{E_0}{n^2}} \qquad \text{avec} \quad E_0 = \SI{13,6}{eV} \quad \text{et} \quad n \in \mathbb{N}^* \; (n = 1, 2, 3, \ldots)
\]
L'entier $n$ est appelé \cle{nombre quantique principal}. L'énergie s'exprime en électronvolts (eV), avec $\SI{1}{eV} = \num{1,60e-19}$ J.
\end{definition}

\begin{propriete}[Origine théorique de la formule (admise)]
La formule $E_n = -\dfrac{E_0}{n^2}$ est une conséquence rigoureuse de la \cle{mécanique quantique} : elle s'obtient en résolvant l'équation de Schrödinger pour l'électron soumis au potentiel coulombien du proton. Elle peut aussi se retrouver dans le modèle plus ancien de Bohr (voir le complément en fin de section). Au programme de Terminale, cette formule est \emph{admise} ; sa justification est confirmée par l'accord parfait avec les spectres expérimentaux (raies de Lyman, Balmer, Paschen...), que nous étudierons dans les sections suivantes.
\end{propriete}

\textbf{La convention du zéro d'énergie}\par

L'énergie d'un système n'est définie qu'à une constante près : il faut donc choisir une \emph{référence}. En physique atomique, la convention universelle est la suivante :

\begin{aretenir}[Convention du zéro d'énergie]
Le zéro de l'énergie est choisi pour l'électron \cle{libre et au repos}, c'est-à-dire infiniment éloigné du noyau et sans vitesse : $E = 0$. Par conséquent, tous les \cle{états liés} (électron attaché au noyau) ont une énergie \emph{négative} : il faut \emph{fournir} de l'énergie à l'atome pour arracher l'électron et l'amener au niveau $E = 0$.
\end{aretenir}

Cette convention est analogue à celle de l'énergie potentielle de pesanteur : un objet au fond d'un puits a une énergie potentielle négative si l'on choisit le zéro au niveau du sol ; il faut lui fournir de l'énergie pour l'extraire du puits. L'électron lié est « au fond du puits » d'attraction électrostatique du proton.

\courssub{État fondamental, états excités, ionisation}

\textbf{L'état fondamental : $n = 1$}\par

Pour $n = 1$, la formule donne :
\[
E_1 = -\dfrac{E_0}{1^2} = -\SI{13,6}{eV}
\]
C'est la plus basse énergie possible (la marche la plus profonde du puits). L'atome y est dans sa configuration la plus stable.

\begin{definition}[État fondamental et états excités]
\begin{itemize}
\item L'\cle{état fondamental} ($n = 1$) est l'état de plus basse énergie de l'atome : c'est son état le plus stable. Un atome laissé à lui-même, sans apport d'énergie, se trouve dans cet état.
\item Tout état correspondant à $n \geq 2$ est appelé \cle{état excité} : l'atome possède alors plus d'énergie que dans son état fondamental. Ces états sont instables : l'atome y séjourne typiquement $10^{-8}$ s avant de « retomber » vers un niveau inférieur.
\item L'\cle{ionisation} correspond au passage de l'électron d'un état lié ($E_n < 0$) à l'état libre ($E \geq 0$) : l'électron est arraché à l'atome, qui devient un ion $\ce{H+}$ (un simple proton).
\end{itemize}
\end{definition}

\textbf{Les états excités : $n = 2, 3, 4, \ldots$}\par

Appliquons la formule pour les premiers niveaux :
\begin{align*}
E_2 &= -\dfrac{13,6}{2^2} = -\dfrac{13,6}{4} = -\SI{3,40}{eV} \\
E_3 &= -\dfrac{13,6}{3^2} = -\dfrac{13,6}{9} = -\SI{1,51}{eV} \\
E_4 &= -\dfrac{13,6}{4^2} = -\dfrac{13,6}{16} = -\SI{0,85}{eV} \\
E_5 &= -\dfrac{13,6}{25} = -\SI{0,54}{eV} \qquad E_6 = -\dfrac{13,6}{36} = -\SI{0,38}{eV}
\end{align*}

On constate un point essentiel : les niveaux se \emph{resserrent} de plus en plus lorsque $n$ augmente, et $E_n \longrightarrow 0$ quand $n \longrightarrow +\infty$. L'écart entre deux niveaux successifs vaut :
\[
E_{n+1} - E_n = 13,6\left(\dfrac{1}{n^2} - \dfrac{1}{(n+1)^2}\right) \text{ eV}
\]
qui diminue rapidement : $E_2 - E_1 = \SI{10,2}{eV}$, $E_3 - E_2 = \SI{1,89}{eV}$, $E_4 - E_3 = \SI{0,66}{eV}$... Le diagramme énergétique ci-dessous visualise cette structure en « marches d'escalier resserrées vers le haut ».

\begin{popfigure}
\begin{center}
\begin{tikzpicture}[scale=1]
% Axe vertical
\draw[->, thick, popdark] (0,-7.5) -- (0,0.8) node[above] {$E$ (eV)};
% Graduations alignées sur les niveaux principaux
\foreach \y/\lab in {-6.8/-13,6, -1.7/-3,40, 0/0} {
  \draw[popdark] (-0.12,\y) -- (0.12,\y);
  \node[left, popdark] at (-0.2,\y) {$\lab$};
}
% Niveaux : échelle y = E/2 (E en eV -> cm)
% n=1 : -13.6 -> -6.8 ; n=2 : -3.4 -> -1.7 ; n=3 : -1.51 -> -0.755 ; n=4 : -0.85 -> -0.425 ; n=5 : -0.544 -> -0.272 ; n=6 : -0.378 -> -0.189 ; ionisation 0
\draw[very thick, popblue] (1,-6.8) -- (5.5,-6.8) node[right, popdark] {$E_1 = -13,6$ eV};
\node[left, popblue] at (1,-6.8) {$n=1$};
\draw[very thick, popblue] (1,-1.7) -- (5.5,-1.7) node[right, popdark] {$E_2 = -3,40$ eV};
\node[left, popblue] at (1,-1.7) {$n=2$};
\draw[very thick, popblue] (1,-0.755) -- (5.5,-0.755) node[right, popdark] {$E_3 = -1,51$ eV};
\node[left, popblue] at (1,-0.755) {$n=3$};
\draw[very thick, popblue] (1,-0.425) -- (5.5,-0.425) node[right, popdark] {$E_4 = -0,85$ eV};
\node[left, popblue] at (1,-0.425) {$n=4$};
\draw[very thick, popblue] (1,-0.272) -- (5.5,-0.272) node[right, popdark] {$E_5 = -0,54$ eV};
\node[left, popblue] at (1,-0.272) {$n=5$};
\draw[very thick, popblue] (1,-0.189) -- (5.5,-0.189) node[right, popdark] {$E_6 = -0,38$ eV};
\node[left, popblue] at (1,-0.189) {$n=6$};
% Niveau d'ionisation
\draw[very thick, dashed, poporange] (1,0) -- (5.5,0) node[right, popdark] {$E = 0$ (ionisation)};
\node[left, poporange] at (1,0) {$n=\infty$};
% Flèche ionisation
\draw[->, very thick, popgreen] (6.3,-6.8) -- (6.3,-0.1) node[midway, right, align=left] {énergie\\d'ionisation\\$E_i = 13,6$ eV};
% Zone états libres
\fill[poporange!15] (1,0.05) rectangle (5.5,0.6);
\node[poporange!80!black] at (3.25,0.32) {\small états libres ($E \geq 0$)};
\end{tikzpicture}
\end{center}
\legende{Diagramme énergétique de l'atome d'hydrogène : les niveaux $E_n = -\dfrac{13,6}{n^2}$ (eV) se resserrent quand $n$ augmente et tendent vers 0. Au-dessus de $E = 0$, l'électron est libre : l'atome est ionisé.}
\end{popfigure}

\textbf{L'ionisation : arracher l'électron}\par

Ioniser l'atome, c'est fournir assez d'énergie à l'électron pour le faire passer de son niveau $E_n$ au niveau $E = 0$ (électron libre au repos). L'\cle{énergie d'ionisation} depuis le niveau $n$ est donc :
\[
E_i = 0 - E_n = \dfrac{E_0}{n^2}
\]
Depuis l'état fondamental ($n = 1$), il faut fournir :
\[
E_i = + \SI{13,6}{eV} = 13,6 \times \num{1,60e-19} = \SI{2,18e-18}{J}
\]
C'est une énergie minuscule à notre échelle, mais considérable à l'échelle d'un atome ! L'atome ionisé, noté $\ce{H+}$, n'est autre qu'un proton isolé.

\begin{attention}[Les deux pièges classiques du Bac]
\begin{itemize}
\item \textbf{Oublier le signe moins.} Les énergies des états liés sont \emph{négatives} : $E_2 = -\SI{3,40}{eV}$, et non $+\SI{3,40}{eV}$. Une énergie de niveau écrite sans le signe moins est systématiquement sanctionnée.
\item \textbf{Croire que l'énergie d'ionisation est négative.} C'est l'inverse : ioniser, c'est \emph{fournir} de l'énergie à l'atome, donc l'énergie d'ionisation est toujours \emph{positive} : $E_i = 0 - E_n = -E_n > 0$. Depuis le fondamental : $E_i = +\SI{13,6}{eV}$.
\end{itemize}
Retiens la phrase-clé : « l'électron lié est endetté énergétiquement ; l'ionisation rembourse sa dette. »
\end{attention}

\begin{popfigure}
\begin{center}
\begin{tikzpicture}[scale=0.9]
% --- Atome fondamental ---
\begin{scope}[shift={(0,0)}]
\fill[poporange] (0,0) circle (0.14);
\node[below=0.35cm, popdark] at (0,0) {\small proton};
\draw[popblue, thick] (0,0) circle (0.75);
\fill[popblue] (0.75,0) circle (0.09);
\draw[->, popblue, thick] (0.75,0) ++(90:0.09) -- ++(0,0.45);
\node[popdark] at (0,-1.5) {\small \textbf{État fondamental}};
\node[popdark] at (0,-1.95) {\small $n=1$, $E_1 = -13,6$ eV};
\end{scope}
% --- Atome excité ---
\begin{scope}[shift={(5,0)}]
\fill[poporange] (0,0) circle (0.14);
\draw[popblue!50, thick] (0,0) circle (0.75);
\draw[popgreen, thick] (0,0) circle (1.35);
\fill[popgreen] (0,1.35) circle (0.09);
\draw[->, popgreen, thick] (0,1.35) ++(0:0.09) -- ++(0.45,0);
\node[popdark] at (0,-1.5) {\small \textbf{État excité}};
\node[popdark] at (0,-1.95) {\small $n \geq 2$ (ex. $E_2 = -3,40$ eV)};
\end{scope}
% --- Atome ionisé ---
\begin{scope}[shift={(10,0)}]
\fill[poporange] (0,0) circle (0.14);
\node[popdark] at (0.1,-0.45) {\small $\ce{H+}$};
\draw[->, very thick, poppink] (1.2,0.6) -- (2.4,1.1);
\fill[poppink] (2.5,1.15) circle (0.09) node[right, popdark] {\small électron libre};
\node[popdark] at (1.2,-1.5) {\small \textbf{Atome ionisé}};
\node[popdark] at (1.2,-1.95) {\small $E \geq 0$ : électron arraché};
\end{scope}
\end{tikzpicture}
\end{center}
\legende{Les trois situations de l'atome d'hydrogène : à gauche, l'électron occupe le niveau fondamental (orbite la plus proche du noyau) ; au centre, l'atome a absorbé de l'énergie et l'électron occupe un niveau excité ; à droite, l'électron a reçu au moins l'énergie d'ionisation et quitte l'atome, qui devient l'ion $\ce{H+}$.}
\end{popfigure}

\courssub{Méthode de calcul et exemples chiffrés}

\begin{methode}[Calculer un niveau d'énergie]
Pour déterminer l'énergie du niveau $n$ de l'atome d'hydrogène :
\begin{enumerate}
\item Calculer $n^2$.
\item Diviser $13,6$ par $n^2$.
\item Affecter le résultat du signe \textbf{MOINS} (état lié) : $E_n = -\dfrac{13,6}{n^2}$ eV.
\item Si une conversion est demandée : multiplier par $\num{1,60e-19}$ pour obtenir des joules.
\end{enumerate}
Exemple : $E_3 = -\dfrac{13,6}{9} = -\SI{1,51}{eV}$, soit $-1,51 \times \num{1,60e-19} = \SI{-2,42e-19}{J}$.
\end{methode}

\begin{exemplebox}[Énergie d'ionisation depuis l'état $n = 2$]
Quelle énergie minimale faut-il fournir à un atome d'hydrogène, pris dans l'état excité $n = 2$, pour l'ioniser ?

L'énergie d'ionisation est la différence entre l'état final (électron libre, $E = 0$) et l'état initial ($E_2$) :
\[
E_i = 0 - E_2 = 0 - (-\SI{3,40}{eV}) = +\SI{3,40}{eV}
\]
Il faut donc fournir seulement $\SI{3,40}{eV}$, soit quatre fois moins que depuis l'état fondamental ($\SI{13,6}{eV}$). C'est logique : l'électron du niveau $n = 2$ est déjà « plus haut » dans le puits d'énergie, il est plus facile à extraire. En joules : $E_i = 3,40 \times \num{1,60e-19} = \SI{5,44e-19}{J}$.
\end{exemplebox}

\begin{exoresolu}[Niveaux d'énergie et ionisation]
\textbf{Énoncé.} On donne $E_0 = \SI{13,6}{eV}$ et $\SI{1}{eV} = \num{1,60e-19}$ J.
\begin{enumerate}
\item Calculer les énergies $E_1$, $E_2$ et $E_4$ de l'atome d'hydrogène, en eV puis en joules.
\item Comparer $E_2 - E_1$ et $E_4 - E_2$. Conclure.
\item Un atome d'hydrogène est dans l'état $n = 4$. Quelle énergie minimale faut-il lui fournir pour l'ioniser ?
\end{enumerate}
\tcblower
\textbf{Solution détaillée.}
\begin{enumerate}
\item Appliquons la formule $E_n = -\dfrac{E_0}{n^2}$ :
\begin{align*}
E_1 &= -\dfrac{13,6}{1} = -\SI{13,6}{eV} = -13,6 \times \num{1,60e-19} = \SI{-2,18e-18}{J} \\
E_2 &= -\dfrac{13,6}{4} = -\SI{3,40}{eV} = \SI{-5,44e-19}{J} \\
E_4 &= -\dfrac{13,6}{16} = -\SI{0,85}{eV} = \SI{-1,36e-19}{J}
\end{align*}
\item Calculons les écarts :
\[
E_2 - E_1 = -3,40 - (-13,6) = \SI{10,2}{eV} \qquad ; \qquad E_4 - E_2 = -0,85 - (-3,40) = \SI{2,55}{eV}
\]
On constate que $E_4 - E_2 < E_2 - E_1$ : les niveaux d'énergie \emph{se resserrent} lorsque $n$ augmente. C'est une conséquence directe de la décroissance en $\dfrac{1}{n^2}$.
\item L'énergie minimale d'ionisation depuis $n = 4$ est :
\[
E_i = 0 - E_4 = +\SI{0,85}{eV} = \SI{1,36e-19}{J}
\]
Cette énergie est \emph{positive} : elle doit être \emph{fournie} à l'atome (par un photon, une collision, un chauffage...). Toute énergie supérieure à $\SI{0,85}{eV}$ ionise aussi l'atome ; l'excédent devient l'énergie cinétique de l'électron libéré.
\end{enumerate}
\end{exoresolu}

\courssub{Complément culturel (hors programme strict) : d'où vient la formule ?}

\begin{savaistu}[Le modèle de Bohr (1913)]
Niels Bohr a retrouvé la formule $E_n = -\dfrac{E_0}{n^2}$ en postulant que le \cle{moment cinétique} de l'électron est quantifié : $m_e v r = n\hbar$, où $\hbar = \dfrac{h}{2\pi}$. En combinant ce postulat avec la force de Coulomb (qui joue le rôle de force centripète : $\dfrac{1}{4\pi\varepsilon_0}\dfrac{e^2}{r^2} = \dfrac{m_e v^2}{r}$), on obtient des rayons quantifiés $r_n = n^2 a_0$ (avec $a_0 = \SI{0,053}{nm}$, le \emph{rayon de Bohr}) et des énergies $E_n = -\dfrac{E_0}{n^2}$. Cette démonstration, très formatrice, peut être proposée en exercice guidé : elle illustre magnifiquement comment un postulat de quantification conduit à des résultats mesurables. Le modèle de Bohr, bien que dépassé par la mécanique quantique (équation de Schrödinger, 1926), reste une étape historique décisive — Bohr reçut le prix Nobel en 1922.
\end{savaistu}

\textbf{Ce qu'il faut retenir de cette section}\par

\begin{aretenir}[L'essentiel]
\begin{itemize}
\item L'énergie de l'électron dans l'atome d'hydrogène est \cle{quantifiée} : $E_n = -\dfrac{E_0}{n^2}$ avec $E_0 = \SI{13,6}{eV}$ et $n \geq 1$ entier.
\item Le zéro d'énergie correspond à l'électron libre au repos ; les états liés ont des énergies \emph{négatives}.
\item $n = 1$ : état \cle{fondamental} ($E_1 = -\SI{13,6}{eV}$), le plus stable ; $n \geq 2$ : états \cle{excités}.
\item Les niveaux se resserrent quand $n$ croît et tendent vers 0.
\item L'\cle{ionisation} ($\ce{H} \rightarrow \ce{H+} + e^-$) exige un apport d'énergie $E_i = -E_n > 0$ : $\SI{13,6}{eV}$ depuis le fondamental.
\end{itemize}
\end{aretenir}

Nous savons désormais que l'électron ne peut occuper que des « marches » énergétiques bien précises. Reste à comprendre comment il \emph{passe} d'une marche à l'autre : c'est précisément l'objet de la section suivante, où nous verrons que chaque transition s'accompagne de l'émission ou de l'absorption d'un photon dont l'énergie $h\nu$ est exactement égale à l'écart entre deux niveaux — la clé qui explique enfin les spectres de raies.

\coursec{Transitions entre niveaux : émission et absorption de photons}

Dans la section précédente, nous avons découvert que l'énergie de l'atome d'hydrogène ne peut prendre que des valeurs bien précises, données par $E_n = -\dfrac{E_0}{n^2}$ avec $E_0 = \SI{13,6}{eV}$. L'énergie de l'atome est \cle{quantifiée} : elle ne peut varier que par \emph{sauts} d'un niveau à un autre. Mais comment l'atome passe-t-il concrètement d'un niveau à un autre ? C'est toute la question des \cle{transitions}, que nous allons maintenant étudier. Tu vas voir que chaque saut s'accompagne de l'émission ou de l'absorption d'un photon, et que c'est ce mécanisme qui explique l'origine de la lumière émise par les lampes spectrales, les néons... et les étoiles.

\courssub{Le principe de conservation de l'énergie lors d'une transition}

\textbf{Intuition : des marches d'escalier, pas une rampe.}

Imagine un enfant qui joue sur un escalier : il peut se tenir sur la première marche, la deuxième, la troisième... mais jamais \emph{entre} deux marches. Pour monter, il doit fournir exactement l'énergie correspondant à la hauteur d'une ou plusieurs marches ; pour descendre, il libère exactement cette énergie. L'atome d'hydrogène fonctionne de la même manière : ses niveaux d'énergie sont les « marches », et l'énergie échangée lors d'un saut est exactement égale à la différence d'énergie entre les deux niveaux.

\begin{definition}[Transition entre niveaux]
Une \cle{transition} est le passage de l'atome d'un niveau d'énergie $E_{n_i}$ à un autre niveau d'énergie $E_{n_p}$. L'atome ne peut jamais se trouver dans un état intermédiaire : la transition est \emph{instantanée} et s'accompagne de l'échange d'un photon avec l'extérieur, de sorte que l'énergie totale du système \{atome + rayonnement\} soit conservée.
\end{definition}

Le principe physique fondamental qui gouverne toutes les transitions est la \cle{conservation de l'énergie}. Considérons un atome qui passe du niveau $n_p$ au niveau $n_i$ en émettant un photon d'énergie $h\nu$. Le bilan énergétique s'écrit :
\[ \underbrace{E_{n_p}}_{\text{énergie initiale de l'atome}} = \underbrace{E_{n_i}}_{\text{énergie finale de l'atome}} + \underbrace{h\nu}_{\text{énergie emportée par le photon}} \]

\begin{experience}[Démonstration guidée : expression de l'énergie du photon échangé]
Retrouvons pas à pas l'expression de l'énergie du photon échangé lors d'une transition entre les niveaux $n_i$ et $n_p$, avec $n_p > n_i$.
\begin{enumerate}
\item \textbf{Écris la conservation de l'énergie.} L'atome passe de $E_{n_p}$ à $E_{n_i}$ en émettant un photon : $E_{n_p} = E_{n_i} + h\nu$.
\item \textbf{Isole l'énergie du photon} : $h\nu = E_{n_p} - E_{n_i}$.
\item \textbf{Remplace par l'expression des niveaux} : $E_{n_p} = -\dfrac{E_0}{n_p^2}$ et $E_{n_i} = -\dfrac{E_0}{n_i^2}$.
\item \textbf{Calcule la différence} :
\[ h\nu = -\frac{E_0}{n_p^2} - \left(-\frac{E_0}{n_i^2}\right) = \frac{E_0}{n_i^2} - \frac{E_0}{n_p^2} = E_0\left(\frac{1}{n_i^2} - \frac{1}{n_p^2}\right) \]
\item \textbf{Vérifie le signe.} Comme $n_p > n_i$, on a $\dfrac{1}{n_i^2} > \dfrac{1}{n_p^2}$, donc $h\nu > 0$ : le résultat est cohérent, un photon transporte toujours une énergie positive.
\end{enumerate}
Cette démarche — partir d'un bilan de conservation — est un grand classique du Baccalauréat : sache la refaire sans hésiter.
\end{experience}

\begin{propriete}[Énergie du photon échangé lors d'une transition]
Lors d'une transition entre le niveau $n_i$ et le niveau $n_p$ de l'atome d'hydrogène, l'énergie du photon échangé est :
\[ |\Delta E| = h\nu = \frac{hc}{\lambda} = E_0\left(\frac{1}{n_i^2} - \frac{1}{n_p^2}\right) \quad \text{avec} \quad E_0 = \SI{13,6}{eV} \]
\begin{itemize}
\item Si $n_p > n_i$ (l'atome \emph{descend}) : il y a \cle{émission} d'un photon.
\item Si $n_p < n_i$ (l'atome \emph{monte}) : il y a \cle{absorption} d'un photon.
\end{itemize}
La longueur d'onde associée est $\lambda = \dfrac{hc}{|\Delta E|}$. (La démonstration à partir de la conservation de l'énergie est exigible.)
\end{propriete}

\courssub{L'émission spontanée : l'atome redescend en émettant un photon}

Un atome porté sur un niveau excité $n_p > 1$ est instable : au bout d'une durée très courte (de l'ordre de $10^{-8}$ s), il retombe spontanément vers un niveau d'énergie inférieure $n_i$. L'énergie excédentaire est emportée par un \cle{photon émis}, dont l'énergie est rigoureusement fixée :
\[ h\nu = E_{n_p} - E_{n_i} = E_0\left(\frac{1}{n_i^2} - \frac{1}{n_p^2}\right) \]
C'est le mécanisme qui produit la lumière dans les lampes à décharge (lampes à vapeur d'hydrogène, de sodium, de mercure) : le courant électrique excite les atomes, qui se désexcitent ensuite en émettant des photons.

\begin{exemplebox}[La raie rouge H$_\alpha$ de l'hydrogène]
Calculons la longueur d'onde du photon émis lors de la transition $n_p = 3 \rightarrow n_i = 2$.
\begin{align*}
\Delta E &= 13{,}6 \times \left(\frac{1}{2^2} - \frac{1}{3^2}\right) = 13{,}6 \times \left(\frac{1}{4} - \frac{1}{9}\right) = 13{,}6 \times \frac{5}{36} \approx \SI{1,89}{eV} \\
\Delta E &= 1{,}89 \times 1{,}60\times 10^{-19} \approx \SI{3,02e-19}{J} \\
\lambda &= \frac{hc}{\Delta E} = \frac{6{,}63\times 10^{-34} \times 3{,}00\times 10^{8}}{3{,}02\times 10^{-19}} \approx \SI{6,56e-7}{m} = \SI{656}{nm}
\end{align*}
Cette raie rouge, appelée raie \cle{H$_\alpha$}, est la raie la plus intense du spectre visible de l'hydrogène. C'est elle qui donne leur couleur rougeâtre aux nébuleuses d'hydrogène photographiées par les télescopes !
\end{exemplebox}

\courssub{L'absorption : un photon sur mesure, rien d'autre}

Voici maintenant le point le plus subtil — et le plus souvent mal compris — de la leçon. Pour qu'un atome absorbe un photon et monte du niveau $n_i$ au niveau $n_p$, il ne suffit pas que le photon soit « énergétique » : il faut que son énergie soit \emph{exactement} égale à la différence entre les deux niveaux :
\[ h\nu = E_{n_p} - E_{n_i} \quad \text{(égalité stricte)} \]

\begin{aretenir}[La règle d'or de l'absorption]
Un atome ne peut absorber un photon que si l'énergie de celui-ci correspond \textbf{exactement} à l'écart entre deux de ses niveaux d'énergie. Un photon « trop » ou « pas assez » énergétique traverse l'atome sans interagir (sauf cas de la photo-ionisation, voir plus loin). L'absorption est \cle{sélective}.
\end{aretenir}

\begin{savaistu}[Pourquoi l'atome ne peut pas absorber « n'importe quelle » lumière]
C'est précisément cette absorption sélective qui explique les spectres d'absorption : quand de la lumière blanche (qui contient toutes les longueurs d'onde) traverse un gaz d'hydrogène froid, seuls les photons dont l'énergie coïncide avec un écart entre niveaux sont absorbés. Les autres passent sans encombre. C'est ainsi que les astrophysiciens identifient les éléments présents dans l'atmosphère des étoiles : chaque raie sombre du spectre stellaire est la signature d'une transition précise d'un élément précis. L'hélium a d'ailleurs été découvert dans le Soleil (du grec \emph{helios}) en 1868, par ses raies d'absorption, avant d'être isolé sur Terre !
\end{savaistu}

\begin{popfigure}
\begin{center}
\begin{tikzpicture}[scale=1.0]
% Axe énergie
\draw[->,thick,popdark] (-0.4,0) -- (-0.4,5.6) node[above] {$E$ (eV)};
% Niveaux : E1=-13.6 -> y=0 ; E2=-3.4 -> y=3.4 ; E3=-1.51 -> y=4.36 ; E4=-0.85 -> y=4.83 ; 0 -> y=5.3
\draw[very thick,popblue] (0,0) -- (3.2,0) node[right] {$n=1$ \quad $E_1=-13{,}6$};
\draw[very thick,popblue] (0,3.4) -- (3.2,3.4) node[right] {$n=2$ \quad $E_2=-3{,}40$};
\draw[very thick,popblue] (0,4.36) -- (3.2,4.36) node[right] {$n=3$ \quad $E_3=-1{,}51$};
\draw[very thick,popblue] (0,4.83) -- (3.2,4.83) node[right] {$n=4$ \quad $E_4=-0{,}85$};
\draw[dashed,popmuted] (0,5.3) -- (3.2,5.3) node[right] {$E=0$ (ionisation)};
% Absorption (bleu) 1 -> 3
\draw[->,very thick,popblue] (0.9,0) -- (0.9,4.36) node[midway,left] {\small absorption $h\nu$};
% Émission (rouge/orange) 3 -> 2
\draw[->,very thick,poporange] (1.8,4.36) -- (1.8,3.4) node[midway,right] {\small émission H$_\alpha$};
% Émission 4 -> 1
\draw[->,very thick,poporange] (2.6,4.83) -- (2.6,0) node[midway,right] {\small émission};
\end{tikzpicture}
\end{center}
\legende{Diagramme énergétique de l'atome d'hydrogène : flèches montantes (bleu) = absorption d'un photon ; flèches descendantes (orange) = émission d'un photon. La longueur de chaque flèche représente l'énergie $h\nu$ du photon échangé.}
\end{popfigure}

\courssub{La désexcitation en cascade : plusieurs sauts, plusieurs photons}

Un atome excité sur un niveau élevé ne redescend pas forcément directement au niveau fondamental. Il peut effectuer le retour \emph{en plusieurs étapes}, en passant par des niveaux intermédiaires : c'est la \cle{désexcitation en cascade}. À chaque saut correspond l'émission d'un photon, dont l'énergie est fixée par les deux niveaux concernés.

Par exemple, un atome porté au niveau $n = 4$ peut :
\begin{itemize}
\item redescendre directement au fondamental : $4 \rightarrow 1$, un seul photon très énergétique ;
\item ou passer par le niveau 2 : $4 \rightarrow 2$ puis $2 \rightarrow 1$, soit deux photons ;
\item ou encore : $4 \rightarrow 3 \rightarrow 2 \rightarrow 1$, soit trois photons.
\end{itemize}

\begin{attention}[Conservation de l'énergie dans une cascade]
La somme des énergies des photons émis au cours d'une cascade est toujours égale à l'énergie du photon qui aurait été émis lors de la transition directe : $h\nu_{4\rightarrow 2} + h\nu_{2\rightarrow 1} = h\nu_{4\rightarrow 1}$. Vérifie-le avec les valeurs numériques : c'est un excellent réflexe de contrôle.
\end{attention}

\begin{popfigure}
\begin{center}
\begin{tikzpicture}[scale=1.0]
\draw[->,thick,popdark] (-0.4,0) -- (-0.4,5.6) node[above] {$E$};
\draw[very thick,popblue] (0,0) -- (5.4,0) node[right] {$n=1$};
\draw[very thick,popblue] (0,3.4) -- (5.4,3.4) node[right] {$n=2$};
\draw[very thick,popblue] (0,4.36) -- (5.4,4.36) node[right] {$n=3$};
\draw[very thick,popblue] (0,4.83) -- (5.4,4.83) node[right] {$n=4$};
% Cascade 4 -> 2 -> 1
\draw[->,very thick,poporange] (1.2,4.83) -- (1.2,3.4);
\draw[->,very thick,poporange] (1.9,3.4) -- (1.9,0);
% Photons émis (flèches ondulées)
\draw[->,thick,popgold,decorate,decoration={snake,amplitude=0.6mm}] (1.2,4.1) -- (2.9,4.35) node[right,popdark] {\small photon 1 : $h\nu_1 = E_4 - E_2$};
\draw[->,thick,popgold,decorate,decoration={snake,amplitude=0.6mm}] (1.9,1.7) -- (3.6,1.95) node[right,popdark] {\small photon 2 : $h\nu_2 = E_2 - E_1$};
% Atome excité
\fill[poporange] (0.6,4.83) circle (2pt);
\fill[popblue] (2.5,0) circle (2pt);
\end{tikzpicture}
\end{center}
\legende{Cascade de désexcitation $n=4 \rightarrow n=2 \rightarrow n=1$ : deux sauts successifs, donc deux photons émis, de longueurs d'onde différentes. On a $h\nu_1 + h\nu_2 = E_4 - E_1$.}
\end{popfigure}

\courssub{La photo-ionisation : quand le photon est trop énergétique}

Que se passe-t-il si un photon incident possède une énergie $h\nu$ \emph{supérieure} à l'énergie d'ionisation $E_i = \SI{13,6}{eV}$ (pour un atome dans son état fondamental) ? La règle d'absorption sélective ne s'applique plus, car l'électron arraché n'est plus lié : son énergie n'est plus quantifiée. Le photon est alors absorbé, l'électron est \cle{arraché} à l'atome, et l'excédent d'énergie se retrouve sous forme d'énergie cinétique de l'électron libre :
\[ h\nu = E_i + E_c \quad \text{soit} \quad E_c = h\nu - E_i \]
C'est la \cle{photo-ionisation}. C'est le seul cas où un atome peut absorber un photon dont l'énergie n'est pas exactement calée sur un écart entre deux niveaux — car l'état final (électron libre) appartient à un continuum d'énergie.

\begin{exemplebox}[Photo-ionisation de l'hydrogène]
Un atome d'hydrogène dans son état fondamental absorbe un photon d'énergie $h\nu = \SI{15,0}{eV}$. L'électron est arraché avec une énergie cinétique :
\[ E_c = 15{,}0 - 13{,}6 = \SI{1,4}{eV} = 1{,}4 \times 1{,}60\times 10^{-19} \approx \SI{2,2e-19}{J} \]
En revanche, un photon de $\SI{12,0}{eV}$ ne peut être absorbé : il ne correspond à aucune transition ($E_2 - E_1 = \SI{10,2}{eV}$, $E_3 - E_1 = \SI{12,1}{eV}$) et son énergie est inférieure au seuil d'ionisation. Il traverse l'atome sans effet.
\end{exemplebox}

\courssub{Unicité des raies : la carte d'identité de l'atome}

Puisque les niveaux d'énergie sont quantifiés et parfaitement déterminés, chaque transition correspond à \emph{une seule} valeur de $\Delta E$, donc à \emph{une seule} longueur d'onde :
\[ \lambda = \frac{hc}{|\Delta E|} \]
La lumière émise ou absorbée par un élément est donc constituée de \cle{raies fines}, infiniment caractéristiques : deux éléments différents n'ont jamais le même spectre, car leurs niveaux d'énergie sont différents.

\begin{aretenir}[Le spectre, carte d'identité de l'élément]
Le spectre d'émission ou d'absorption d'un élément chimique est unique : il constitue sa \cle{carte d'identité}. L'analyse spectrale permet d'identifier la composition d'un gaz, d'une flamme, de l'atmosphère d'une étoile ou d'un échantillon, sans contact direct. C'est le principe de la spectroscopie, utilisée en astrophysique, en médecine et dans le contrôle qualité industriel.
\end{aretenir}

\courssub{Méthode de résolution et applications type Bac}

\begin{methode}[Déterminer la longueur d'onde d'une transition]
\begin{enumerate}
\item \textbf{Calcule} l'énergie échangée en électronvolts : $\Delta E = 13{,}6 \times \left(\dfrac{1}{n_i^2} - \dfrac{1}{n_p^2}\right)$ (en eV).
\item \textbf{Convertis} en joules : multiplie par $1{,}60\times 10^{-19}$ J/eV.
\item \textbf{Applique} la relation de Planck-Einstein : $\lambda = \dfrac{hc}{\Delta E}$ avec $h = 6{,}63\times 10^{-34}$ J·s et $c = 3{,}00\times 10^{8}$ m/s.
\item \textbf{Convertis} en nanomètres ($1\ \text{nm} = 10^{-9}$ m) et conclus sur le domaine spectral (UV, visible, IR).
\end{enumerate}
Astuce de calcul rapide : $\lambda\ (\text{nm}) \approx \dfrac{1240}{\Delta E\ (\text{eV})}$.
\end{methode}

\begin{attention}[Erreur fréquente : inverser $n_i$ et $n_p$]
Dans la formule $\Delta E = E_0\left(\dfrac{1}{n_i^2} - \dfrac{1}{n_p^2}\right)$, $n_i$ est le niveau \emph{le plus bas} (le plus proche du noyau) et $n_p$ le niveau \emph{le plus haut}. Si tu les inverses, tu obtiens une énergie négative — ce qui est absurde pour un photon. \textbf{Réflexe de survie au Bac} : vérifie toujours que $\Delta E$ est \emph{positif} pour une émission, et que la longueur d'onde obtenue tombe dans un domaine spectral plausible. Autre piège : oublier de convertir les eV en joules avant d'appliquer $\lambda = hc/\Delta E$.
\end{attention}

\begin{exoresolu}[Cascade d'émission d'un atome d'hydrogène excité]
Un atome d'hydrogène, initialement dans son état fondamental, absorbe un photon et passe au niveau $n = 3$. Il se désexcite ensuite en cascade : $n = 3 \rightarrow n = 2$, puis $n = 2 \rightarrow n = 1$.
\begin{enumerate}
\item Calculer l'énergie du photon absorbé.
\item Calculer les longueurs d'onde $\lambda_1$ et $\lambda_2$ des deux photons émis. Préciser leur domaine spectral.
\end{enumerate}
Données : $E_0 = \SI{13,6}{eV}$ ; $h = 6{,}63\times 10^{-34}$ J·s ; $c = 3{,}00\times 10^{8}$ m/s ; $1\ \text{eV} = 1{,}60\times 10^{-19}$ J.
\tcblower
\textbf{1. Énergie du photon absorbé.} L'absorption se fait du niveau $n_i = 1$ vers le niveau $n_p = 3$ :
\[ \Delta E_{\text{abs}} = 13{,}6 \times \left(\frac{1}{1^2} - \frac{1}{3^2}\right) = 13{,}6 \times \frac{8}{9} \approx \SI{12,1}{eV} \]
Le photon absorbé doit avoir \emph{exactement} cette énergie.

\textbf{2. Photons émis lors de la cascade.}

Transition $3 \rightarrow 2$ :
\[ \Delta E_1 = 13{,}6 \times \left(\frac{1}{4} - \frac{1}{9}\right) = \SI{1,89}{eV} \quad \Rightarrow \quad \lambda_1 = \frac{1240}{1{,}89} \approx \SI{656}{nm} \]
C'est la raie rouge H$_\alpha$ (domaine visible).

Transition $2 \rightarrow 1$ :
\[ \Delta E_2 = 13{,}6 \times \left(\frac{1}{1} - \frac{1}{4}\right) = 13{,}6 \times \frac{3}{4} = \SI{10,2}{eV} \quad \Rightarrow \quad \lambda_2 = \frac{1240}{10{,}2} \approx \SI{122}{nm} \]
Ce photon appartient à l'\emph{ultraviolet} (série de Lyman).

\textbf{Vérification} : $\Delta E_1 + \Delta E_2 = 1{,}89 + 10{,}2 \approx \SI{12,1}{eV} = \Delta E_{\text{abs}}$ : l'énergie absorbée est bien intégralement restituée par la cascade. La conservation de l'énergie est satisfaite.
\end{exoresolu}

\begin{exercice}[À toi de jouer]
Un atome d'hydrogène dans son état fondamental reçoit un photon d'énergie $\SI{12,75}{eV}$.
\begin{enumerate}
\item Montrer que ce photon peut être absorbé et déterminer le niveau atteint.
\item L'atome se désexcite directement vers le niveau $n = 2$. Calculer la longueur d'onde du photon émis.
\item Un photon d'énergie $\SI{11,0}{eV}$ peut-il être absorbé par un atome d'hydrogène au repos dans son état fondamental ? Justifier.
\end{enumerate}
\end{exercice}

\begin{corrige}[Exercice : Absorption et désexcitation de l'hydrogène]
\textbf{1. Absorption du photon de $\SI{12,75}{eV}$.}
L'énergie nécessaire pour passer de l'état fondamental ($n_i=1$) au niveau $n_p$ est :
\[ \Delta E = 13{,}6 \left(1 - \frac{1}{n_p^2}\right) \]
On cherche $n_p$ tel que $\Delta E = 12{,}75$ eV.
\[ 1 - \frac{1}{n_p^2} = \frac{12{,}75}{13{,}6} = 0{,}9375 \quad \Rightarrow \quad \frac{1}{n_p^2} = 0{,}0625 = \frac{1}{16} \quad \Rightarrow \quad n_p = 4 \]
Le photon est absorbé et l'atome atteint le niveau $n=4$.

\textbf{2. Désexcitation directe $4 \rightarrow 2$.}
\[ \Delta E = 13{,}6 \left(\frac{1}{2^2} - \frac{1}{4^2}\right) = 13{,}6 \left(\frac{1}{4} - \frac{1}{16}\right) = 13{,}6 \times \frac{3}{16} = \SI{2,55}{eV} \]
\[ \lambda = \frac{1240}{2{,}55} \approx \SI{486}{nm} \]
Ce photon est dans le domaine \emph{visible} (bleu-vert, raie H$_\beta$ de la série de Balmer).

\textbf{3. Photon de $\SI{11,0}{eV}$.}
Les transitions possibles depuis $n=1$ sont :
$1 \rightarrow 2$ : $\Delta E = \SI{10,2}{eV}$
$1 \rightarrow 3$ : $\Delta E = \SI{12,1}{eV}$
$1 \rightarrow 4$ : $\Delta E = \SI{12,75}{eV}$
$1 \rightarrow \infty$ (ionisation) : $\Delta E = \SI{13,6}{eV}$
L'énergie $\SI{11,0}{eV}$ ne correspond à aucune de ces valeurs et est inférieure au seuil d'ionisation. \textbf{Non}, ce photon ne peut pas être absorbé.
\end{corrige}

En résumé, les transitions de l'atome d'hydrogène obéissent à une logique implacable dictée par la conservation de l'énergie : l'atome ne gagne ou ne perd de l'énergie que par sauts, et chaque saut impose l'énergie — donc la couleur — du photon échangé. Dans la prochaine section, nous verrons comment l'ensemble de ces transitions se manifeste expérimentalement sous la forme des spectres d'émission et d'absorption.

\coursec{Spectres d'émission et d'absorption}

Dans la section précédente, nous avons établi qu'un atome ne peut échanger de l'énergie avec la lumière que par \emph{paquets} bien précis : lors d'une transition entre deux niveaux $E_n$ et $E_p$, il absorbe ou émet un photon d'énergie $h\nu = |E_n - E_p|$, donc de longueur d'onde $\lambda = hc/|E_n - E_p|$ parfaitement déterminée. Mais comment \emph{observe-t-on} concrètement ces échanges ? Comment un physicien « voit-il » les niveaux d'énergie d'un atome ? La réponse tient en un mot : les \cle{spectres}. En dispersant la lumière avec un prisme ou un réseau, on obtient une empreinte lumineuse qui révèle, raie par raie, la structure énergétique de la matière. C'est l'une des plus belles confirmations expérimentales de la physique quantique, et un outil d'analyse d'une puissance extraordinaire, du laboratoire de lycée jusqu'aux télescopes qui sondent les étoiles.

\courssub{Les trois types de spectres}

\textbf{Observation préalable.} Éclaire un prisme avec la lumière d'une lampe à filament : tu obtiens un dégradé ininterrompu de couleurs, du violet au rouge. Recommence avec la lumière d'un tube à décharge contenant de l'hydrogène (comme ceux utilisés en TP au lycée) : catastrophe pour ton intuition ! Au lieu d'un arc-en-ciel continu, tu ne vois que quelques \emph{raies colorées fines}, isolées sur fond noir. La nature de la source change radicalement l'aspect du spectre. Classons les trois situations fondamentales.

\begin{definition}[Spectre continu]
Un \cle{spectre continu} est un spectre contenant \textbf{toutes} les longueurs d'onde d'un domaine, sans interruption. Il est émis par un \textbf{corps chaud dense} : solide porté à incandescence (filament de tungstène d'une lampe classique, braise), liquide ou gaz sous forte pression (la surface du Soleil). Plus le corps est chaud, plus son spectre s'enrichit vers le violet.
\end{definition}

\begin{definition}[Spectre d'émission de raies]
Un \cle{spectre d'émission de raies} est l'ensemble des radiations émises par un atome (ou un ion) excité qui se désexcite. Il est constitué de \textbf{raies colorées fines sur fond noir}, chaque raie correspondant à une longueur d'onde précise. On l'obtient en dispersant la lumière émise par un gaz \textbf{chaud ou excité} (décharge électrique dans un tube, flamme, arc électrique) à \textbf{faible pression}.
\end{definition}

\begin{definition}[Spectre d'absorption]
Un \cle{spectre d'absorption} est l'ensemble des radiations absorbées par un atome éclairé en lumière blanche. On l'obtient en faisant traverser un gaz \textbf{froid} (faible pression) par la lumière d'une source de spectre continu : on observe alors un \textbf{fond continu traversé de raies noires}, aux longueurs d'onde que le gaz a prélevées.
\end{definition}

\begin{experience}[Les trois spectres au banc d'optique]
\textbf{Matériel :} source de lumière blanche (lampe à filament), tube à décharge à hydrogène, cuve de gaz froid, prisme (ou réseau), fente, écran.

\textbf{Protocole et observations :}
\begin{enumerate}
\item \emph{Source blanche seule} $\rightarrow$ prisme : bande continue arc-en-ciel sur l'écran.
\item \emph{Tube à hydrogène seul} $\rightarrow$ prisme : quatre raies visibles isolées sur fond noir (rouge, turquoise, bleu-violet, violet).
\item \emph{Source blanche} $\rightarrow$ \emph{cuve de gaz froid} $\rightarrow$ prisme : fond continu, mais des raies noires apparaissent \emph{exactement} aux positions des raies colorées de l'étape 2.
\end{enumerate}

\textbf{Interprétation :} le gaz froid absorbe les mêmes photons qu'il émettrait s'il était excité. Les raies noires d'absorption et les raies brillantes d'émission sont les deux visages d'une même réalité : les transitions entre niveaux quantifiés.
\end{experience}

\begin{popfigure}
\begin{tikzpicture}[x=1cm,y=1cm]
% --- Bande 1 : spectre continu ---
\shade[left color=violet!80!black, right color=red!90!black] (0,4.2) rectangle (9,5.0);
\node[left, font=\small\bfseries, text=popdark] at (-0.1,4.6) {Spectre continu};
\node[left, font=\scriptsize, text=popmuted] at (-0.1,4.25) {(corps chaud dense)};
% --- Bande 2 : spectre d'émission de l'hydrogène ---
\fill[black] (0,2.6) rectangle (9,3.4);
\node[left, font=\small\bfseries, text=popdark] at (-0.1,3.0) {Spectre d'émission};
\node[left, font=\scriptsize, text=popmuted] at (-0.1,2.65) {(H gazeux excité)};
% raies Balmer : 656, 486, 434, 410 nm ; echelle : x = (700-lambda)/(700-400)*9
\fill[red!85!black]    (1.32,2.6) rectangle (1.42,3.4);
\fill[cyan!90!black]   (6.42,2.6) rectangle (6.52,3.4);
\fill[blue!80!black]   (7.98,2.6) rectangle (8.08,3.4);
\fill[violet!90!black] (8.70,2.6) rectangle (8.80,3.4);
% --- Bande 3 : spectre d'absorption ---
\shade[left color=violet!80!black, right color=red!90!black] (0,1.0) rectangle (9,1.8);
\fill[black] (1.32,1.0) rectangle (1.42,1.8);
\fill[black] (6.42,1.0) rectangle (6.52,1.8);
\fill[black] (7.98,1.0) rectangle (8.08,1.8);
\fill[black] (8.70,1.0) rectangle (8.80,1.8);
\node[left, font=\small\bfseries, text=popdark] at (-0.1,1.4) {Spectre d'absorption};
\node[left, font=\scriptsize, text=popmuted] at (-0.1,1.05) {(H froid en lumière blanche)};
% --- lignes pointillées d'alignement ---
\foreach \x in {1.37,6.47,8.03,8.75} {
  \draw[dashed, popmuted, thin] (\x,1.0) -- (\x,5.0);
}
% --- légendes des raies ---
\node[above, font=\scriptsize, text=popdark] at (1.37,5.05) {$656$ nm};
\node[above, font=\scriptsize, text=popdark] at (6.47,5.05) {$486$ nm};
\node[above, font=\scriptsize, text=popdark] at (8.03,5.05) {$434$ nm};
\node[above, font=\scriptsize, text=popdark] at (8.75,5.05) {$410$ nm};
% --- axe des longueurs d'onde ---
\draw[->, popdark] (0,0.55) -- (9.3,0.55) node[right, font=\scriptsize] {$\lambda$};
\node[below, font=\scriptsize, text=popmuted] at (0.3,0.55) {violet ($400$ nm)};
\node[below, font=\scriptsize, text=popmuted] at (8.7,0.55) {rouge ($700$ nm)};
\end{tikzpicture}
\legende{Les trois types de spectres, alignés en longueur d'onde. Les raies noires d'absorption se situent exactement aux positions des raies brillantes d'émission : mêmes longueurs d'onde, donc mêmes transitions.}
\end{popfigure}

\begin{popfigure}
\begin{tikzpicture}[font=\small, >=Stealth]
% source
\draw[fill=popgoldL, draw=popdark] (0,0) circle (0.35);
\node[font=\scriptsize, text=popdark] at (0,-0.75) {Source de};
\node[font=\scriptsize, text=popdark] at (0,-1.05) {lumière blanche};
\foreach \a in {20,0,-20} {\draw[->, poporange, thick] (0.4,0) -- ++(\a:0.9);}
% fente
\draw[fill=popdark] (1.55,-0.5) rectangle (1.7,-0.08);
\draw[fill=popdark] (1.55,0.08) rectangle (1.7,0.5);
\node[font=\scriptsize, text=popdark] at (1.62,-0.75) {fente};
% cuve de gaz
\draw[fill=popblueL!50, draw=popblue, thick] (2.4,-0.55) rectangle (4.2,0.55);
\node[font=\scriptsize, text=popblue] at (3.3,0) {gaz froid};
\node[font=\scriptsize, text=popblue] at (3.3,-0.85) {(absorbe certaines $\lambda$)};
\draw[->, poporange, thick] (1.75,0) -- (2.4,0);
\draw[->, poporange, thick] (4.2,0) -- (5.1,0);
% prisme
\draw[fill=poppurpleL!60, draw=poppurple, thick] (5.4,-0.6) -- (6.1,0.6) -- (6.8,-0.6) -- cycle;
\node[font=\scriptsize, text=poppurple] at (6.1,-0.85) {prisme};
% faisceau dispersé
\draw[->, red!85!black, thick] (6.5,0.1) -- (8.6,0.55);
\draw[->, popgreen, thick] (6.5,0.05) -- (8.6,0.1);
\draw[->, violet!85!black, thick] (6.5,0) -- (8.6,-0.35);
% écran
\draw[fill=popcream, draw=popdark, thick] (8.7,-1.0) rectangle (8.9,1.0);
\node[font=\scriptsize, text=popdark, rotate=90] at (9.15,0) {écran};
% raies noires sur l'écran
\fill[black] (8.7,0.30) rectangle (8.9,0.38);
\fill[black] (8.7,-0.05) rectangle (8.9,0.03);
\fill[black] (8.7,-0.40) rectangle (8.9,-0.32);
\end{tikzpicture}
\legende{Dispositif d'observation d'un spectre d'absorption : la lumière blanche traverse le gaz froid, qui prélève certaines longueurs d'onde ; le prisme disperse le faisceau transmis et l'écran révèle un fond continu strié de raies noires.}
\end{popfigure}

\courssub{Interprétation énergétique des raies}

Pourquoi des raies, et non un continuum ? Parce que l'énergie d'un atome est \cle{quantifiée}. Reprenons le bilan énergétique de la section précédente.

\textbf{Émission.} Un atome excité (par une décharge électrique ou une flamme) se désexcite d'un niveau $E_n$ vers un niveau inférieur $E_p$ ($n > p$). L'énergie libérée est emportée par un photon unique :
\[
h\nu = E_n - E_p \quad \Longrightarrow \quad \boxed{\lambda_{\text{émission}} = \dfrac{hc}{E_n - E_p}}
\]
Comme les niveaux $E_n$ sont discrets, seules certaines valeurs de $\lambda$ sont possibles : le spectre est fait de \emph{raies}, pas d'un continuum. Chaque raie brillante est la signature lumineuse d'une \textbf{désexcitation}.

\textbf{Absorption.} Un atome froid, au repos dans son état fondamental (ou un état bas), ne peut absorber un photon que si celui-ci transporte \emph{exactement} l'énergie d'une transition permise :
\[
h\nu = E_p - E_n \quad (p > n) \quad \Longrightarrow \quad \boxed{\lambda_{\text{absorption}} = \dfrac{hc}{E_p - E_n}}
\]
Dans la lumière blanche qui traverse le gaz, les photons de ces longueurs d'onde précises sont « capturés » : ils manquent à l'arrivée, d'où les \emph{raies noires}. Chaque raie noire est la signature d'une \textbf{excitation}.

\begin{propriete}[Coïncidence des raies d'émission et d'absorption]
Les raies d'absorption d'un gaz ont \textbf{les mêmes longueurs d'onde} que ses raies d'émission. C'est une conséquence directe de la quantification des niveaux d'énergie : émission et absorption mettent en jeu les \emph{mêmes} écarts $|E_n - E_p|$, donc les mêmes photons $h\nu = hc/\lambda$.
\end{propriete}

\begin{attention}[Pièges fréquents]
\begin{itemize}
\item Un spectre de raies n'est \textbf{pas} un spectre continu : ne confonds pas la bande arc-en-ciel du filament avec les raies isolées du gaz excité.
\item Toutes les raies ne sont \textbf{pas} dans le visible ! L'hydrogène n'a que quatre raies visibles (série de Balmer) ; la série de Lyman est entièrement dans l'\emph{ultraviolet} et la série de Paschen dans l'\emph{infrarouge}. Un spectre complet s'étend bien au-delà de ce que l'œil perçoit.
\item Un photon d'énergie « presque » égale à $E_n - E_p$ n'est \textbf{pas} absorbé : la résonance est exacte. C'est tout ou rien.
\item Le gaz absorbant doit être \textbf{froid} (peu dense, atomes dans les états bas) ; un gaz chaud et dense émet un spectre continu.
\end{itemize}
\end{attention}

\courssub{Le spectre, carte d'identité de l'élément}

Chaque élément chimique possède son propre diagramme énergétique : l'hydrogène a ses niveaux en $-13{,}6/n^2$ (eV), le sodium, le néon, le mercure ont les leurs, tous différents. Par conséquent :

\begin{aretenir}[Caractère identifiant du spectre]
Deux éléments différents n'ont \textbf{jamais} le même spectre de raies. Le spectre d'émission (ou d'absorption) d'un atome est sa \cle{carte d'identité} lumineuse : il permet d'identifier un élément sans aucune réaction chimique, même à distance, même en quantité infime.
\end{aretenir}

\begin{methode}[Identifier un élément dans un spectre inconnu]
\begin{enumerate}
\item Repérer les raies du spectre inconnu et mesurer leurs longueurs d'onde $\lambda_1, \lambda_2, \ldots$ (par lecture sur l'échelle graduée du spectroscope).
\item Comparer ces valeurs à un \textbf{tableau de raies de référence} des éléments connus.
\item Conclure : l'élément dont \emph{toutes} les raies caractéristiques coïncident avec celles du spectre inconnu est présent dans l'échantillon. Si plusieurs familles de raies coexistent, l'échantillon est un mélange.
\item Vérifier la cohérence : pour l'hydrogène, contrôler que $\lambda = hc/|E_n - E_p|$ redonne les valeurs mesurées.
\end{enumerate}
\end{methode}

\begin{exemplebox}[La raie jaune du sodium et les lampadaires]
Les lampes à vapeur de sodium qui éclairent en jaune-orangé certaines rues de Douala et de Yaoundé doivent leur couleur à une raie intense de l'atome Na : le doublet jaune de longueur d'onde moyenne $\lambda \approx 589$ nm. L'énergie du photon correspondant est :
\[
E = \dfrac{hc}{\lambda} = \dfrac{6{,}63\times 10^{-34}\times 3{,}00\times 10^{8}}{589\times 10^{-9}} \approx 3{,}38\times 10^{-19}\ \text{J} \approx 2{,}11\ \text{eV}.
\]
C'est précisément l'écart entre le premier état excité et l'état fondamental de l'atome de sodium. Même principe pour les enseignes lumineuses au néon (rouge-orangé) et les tubes fluorescents (vapeur de mercure + poudre fluorescente).
\end{exemplebox}

\begin{exoresolu}[Identification d'un gaz inconnu]
\textbf{Énoncé.} On disperse la lumière émise par un tube à décharge contenant un gaz inconnu. On observe quatre raies visibles de longueurs d'onde : $656$ nm, $486$ nm, $434$ nm et $410$ nm. Un tableau de référence donne pour l'hydrogène les raies $656{,}3$ nm, $486{,}1$ nm, $434{,}0$ nm, $410{,}2$ nm ; pour le néon, des raies à $585$ nm, $640$ nm, $703$ nm\ldots
\begin{enumerate}
\item Identifier le gaz.
\item Vérifier que la raie à $656$ nm correspond à la transition $n = 3 \rightarrow n = 2$ de l'hydrogène. On donne $E_n = -13{,}6/n^2$ (eV), $h = 6{,}63\times 10^{-34}$ J.s, $c = 3{,}00\times 10^{8}$ m/s, $1$ eV $= 1{,}60\times 10^{-19}$ J.
\end{enumerate}
\tcblower
\textbf{Solution.}
\begin{enumerate}
\item Les quatre raies observées coïncident exactement avec les raies de référence de l'hydrogène, et aucune raie du néon n'apparaît : le gaz est donc du \textbf{dihydrogène} (dont les molécules sont dissociées en atomes par la décharge).
\item Énergie de la transition $3 \rightarrow 2$ :
\begin{align*}
\Delta E &= E_3 - E_2 = -13{,}6\left(\dfrac{1}{3^2} - \dfrac{1}{2^2}\right) = 13{,}6\times \dfrac{5}{36} \approx 1{,}89\ \text{eV}\\
&= 1{,}89\times 1{,}60\times 10^{-19} \approx 3{,}02\times 10^{-19}\ \text{J}.
\end{align*}
Longueur d'onde du photon émis :
\[
\lambda = \dfrac{hc}{\Delta E} = \dfrac{6{,}63\times 10^{-34}\times 3{,}00\times 10^{8}}{3{,}02\times 10^{-19}} \approx 6{,}58\times 10^{-7}\ \text{m} = 658\ \text{nm}.
\]
À la précision des données près, on retrouve la raie rouge observée à $656$ nm : la vérification est concluante.
\end{enumerate}
\end{exoresolu}

\courssub{Les raies de Fraunhofer : lire la composition du Soleil}

En 1814, l'opticien allemand Joseph von Fraunhofer observe le spectre solaire avec un spectroscope soigné : le magnifique continuum arc-en-ciel est traversé de \emph{centaines de raies noires fines}. Que se passe-t-il ? La surface du Soleil (la photosphère, à environ $5\,800$ K) est un corps dense qui émet un spectre \emph{continu}. Mais cette lumière traverse ensuite la chromosphère, l'atmosphère solaire plus \emph{froide} : les atomes qu'elle contient (hydrogène, hélium, sodium, calcium, fer\ldots) absorbent leurs raies caractéristiques. Les \cle{raies de Fraunhofer} sont donc le spectre d'absorption de l'atmosphère solaire, imprimé sur le continuum de la photosphère.

En comparant les positions de ces raies noires aux spectres d'émission mesurés en laboratoire, on détermine la \textbf{composition chimique} d'un astre situé à 150 millions de kilomètres, sans y envoyer le moindre échantillon. C'est la naissance de l'astrophysique spectroscopique : la même méthode révèle aujourd'hui la composition des étoiles, des nébuleuses et des atmosphères d'exoplanètes.

\begin{savaistu}[L'hélium : un élément découvert dans le Soleil avant la Terre]
Lors de l'éclipse totale de 1868, l'astronome Jules Janssen observe une raie jaune inconnue ($\lambda \approx 588$ nm) dans le spectre de la couronne solaire, qui ne correspond à \emph{aucun} élément terrestre connu. Norman Lockyer propose qu'il s'agit d'un nouvel élément et le baptise \emph{hélium}, du grec \emph{helios} : le Soleil. Il faudra attendre 1895 pour que William Ramsay isole ce gaz sur Terre, dans un minerai d'uranium. L'hélium reste à ce jour le seul élément découvert hors de notre planète : la preuve éclatante que le spectre est une carte d'identité universelle de la matière.
\end{savaistu}

\begin{aretenir}[L'essentiel de la section]
\begin{itemize}
\item Corps chaud \textbf{dense} $\rightarrow$ spectre \textbf{continu}.
\item Gaz excité à faible pression $\rightarrow$ spectre d'\textbf{émission} de raies (raies colorées sur fond noir) : désexcitations.
\item Gaz froid éclairé en lumière blanche $\rightarrow$ spectre d'\textbf{absorption} (raies noires sur fond continu) : excitations.
\item Raies d'émission et d'absorption aux \textbf{mêmes} longueurs d'onde : $\lambda = hc/|E_n - E_p|$.
\item Chaque élément possède un spectre unique : outil d'identification en chimie et en astrophysique (raies de Fraunhofer).
\end{itemize}
\end{aretenir}

Dans la section suivante, nous appliquerons cette lecture des spectres à l'atome le plus simple de tous, l'hydrogène, en organisant ses raies en \emph{séries} — Lyman, Balmer, Paschen — dont la régularité presque musicale a guidé Bohr vers son modèle quantique.

\coursec{Séries spectrales de l'hydrogène : Lyman, Balmer, Paschen}

Dans la section précédente, nous avons appris à interpréter les spectres d'émission et d'absorption : chaque raie correspond à une transition entre deux niveaux d'énergie, et sa longueur d'onde est fixée par la relation $\Delta E = \dfrac{hc}{\lambda}$. Mais le spectre de l'hydrogène compte des dizaines de raies ! Faut-il les apprendre toutes par cœur ? Heureusement non : les physiciens ont remarqué que ces raies se rangent en \cle{familles} parfaitement organisées, appelées \cle{séries spectrales}. C'est cette organisation remarquable que nous allons découvrir — et qui constitue l'une des plus belles réussites du modèle de Bohr.

\courssub{Notion de série spectrale}

Reprenons le diagramme énergétique de l'atome d'hydrogène. Une transition d'émission fait passer l'électron d'un niveau supérieur $n_p$ vers un niveau inférieur $n_i$ (avec $n_p > n_i$), en libérant un photon d'énergie :
\[ \Delta E = E_{n_p} - E_{n_i} = E_0\left(\frac{1}{n_i^2} - \frac{1}{n_p^2}\right) \qquad \text{avec } E_0 = 13{,}6\ \text{eV}. \]

L'idée directrice est la suivante : au lieu de classer les raies par leur longueur d'onde, classons-les \emph{par leur niveau d'arrivée}. Toutes les transitions qui se terminent sur le même niveau $n_i$ forment une famille cohérente : leurs énergies sont du même ordre de grandeur, donc leurs raies se situent dans le même domaine du spectre électromagnétique.

\begin{definition}[Série spectrale]
Une \cle{série spectrale} de l'hydrogène est l'ensemble des transitions d'émission aboutissant au \textbf{même niveau final} $n_i$, depuis tous les niveaux supérieurs $n_p > n_i$. Les trois séries au programme sont :
\begin{itemize}
\item la \cle{série de Lyman} : transitions $n_p \to 1$ ($n_p \geq 2$), toutes situées dans l'\textbf{ultraviolet} ;
\item la \cle{série de Balmer} : transitions $n_p \to 2$ ($n_p \geq 3$), dont quatre raies sont dans le \textbf{visible} ;
\item la \cle{série de Paschen} : transitions $n_p \to 3$ ($n_p \geq 4$), toutes situées dans l'\textbf{infrarouge}.
\end{itemize}
\end{definition}

\begin{attention}[Émission ou absorption ?]
Les séries sont définies ici pour l'\textbf{émission} (l'atome descend et émet un photon). En absorption, le même niveau $n_i$ joue le rôle de niveau de \emph{départ} : l'atome absorbe exactement les mêmes longueurs d'onde. Mais attention : à température ordinaire, l'atome d'hydrogène est dans son état fondamental $n = 1$ : seule la série de Lyman est observable en absorption. C'est pourquoi les spectres d'absorption des gaz froids montrent surtout des raies UV.
\end{attention}

\courssub{La série de Lyman : le domaine de l'ultraviolet}

La série de Lyman regroupe toutes les transitions qui ramènent l'électron au niveau fondamental $n_i = 1$. Comme le niveau $n = 1$ est le plus profond ($E_1 = -13{,}6\ \text{eV}$), ces transitions sont les \textbf{plus énergétiques} de toutes : les photons émis ont de courtes longueurs d'onde, dans l'ultraviolet.

\begin{exemplebox}[Première raie de Lyman : la raie L$\alpha$]
Transition $n_p = 2 \to n_i = 1$ :
\begin{align*}
\Delta E &= 13{,}6\left(\frac{1}{1^2} - \frac{1}{2^2}\right) = 13{,}6 \times \frac{3}{4} = 10{,}2\ \text{eV} \\
\lambda &= \frac{hc}{\Delta E} = \frac{6{,}63\times 10^{-34}\times 3{,}00\times 10^{8}}{10{,}2\times 1{,}60\times 10^{-19}} \approx 1{,}22\times 10^{-7}\ \text{m} = 122\ \text{nm}.
\end{align*}
Cette raie, notée L$\alpha$, est dans l'ultraviolet ($\lambda < 400\ \text{nm}$). C'est d'ailleurs la raie la plus intense émise par le Soleil dans l'UV ; elle est absorbée par la haute atmosphère terrestre.
\end{exemplebox}

La transition la plus énergétique possible est la chute depuis l'infini ($n_p \to \infty$, énergie nulle) vers $n = 1$ : elle libère exactement l'énergie d'ionisation, $\Delta E = 13{,}6\ \text{eV}$, soit $\lambda = \dfrac{hc}{13{,}6\ \text{eV}} \approx 91{,}2\ \text{nm}$. C'est la \cle{limite de la série de Lyman} : aucune raie de cette série ne peut avoir une longueur d'onde inférieure à 91,2 nm.

\courssub{La série de Balmer : les quatre raies visibles}

La série de Balmer regroupe les transitions aboutissant au niveau $n_i = 2$. C'est historiquement la première série découverte, car c'est la seule qui possède des raies dans le domaine visible — celles que l'on observe au lycée avec une lampe à hydrogène et un réseau de diffraction.

\begin{exemplebox}[Les quatre raies visibles de Balmer]
En appliquant $\Delta E = 13{,}6\left(\dfrac{1}{4} - \dfrac{1}{n_p^2}\right)$ eV puis $\lambda = \dfrac{hc}{\Delta E}$, on obtient :
\begin{center}
\begin{tabularx}{\linewidth}{|c|c|X|c|}\hline
\rowcolor{popblueL} \textbf{Raie} & \textbf{Transition} & \textbf{$\Delta E$ (eV)} & \textbf{$\lambda$ (nm)} \\ \hline
H$\alpha$ & $3 \to 2$ & $13{,}6\left(\tfrac14-\tfrac19\right)=1{,}89$ & 656 (rouge) \\ \hline
H$\beta$ & $4 \to 2$ & $13{,}6\left(\tfrac14-\tfrac1{16}\right)=2{,}55$ & 486 (bleu-vert) \\ \hline
H$\gamma$ & $5 \to 2$ & $13{,}6\left(\tfrac14-\tfrac1{25}\right)=2{,}86$ & 434 (violet) \\ \hline
H$\delta$ & $6 \to 2$ & $13{,}6\left(\tfrac14-\tfrac1{36}\right)=3{,}02$ & 410 (violet extrême) \\ \hline
\end{tabularx}
\end{center}
La raie H$\alpha$, rouge, est la plus célèbre : c'est elle qui donne leur couleur rougeâtre aux nébuleuses d'hydrogène photographiées par les astronomes.
\end{exemplebox}

\begin{attention}[Piège classique : « toute la série de Balmer est visible »]
\textbf{Faux !} Seules \textbf{quatre} raies de Balmer sont visibles (H$\alpha$ à H$\delta$). Dès $n_p = 7 \to 2$, on trouve $\lambda \approx 397\ \text{nm}$, déjà dans l'ultraviolet proche. Toutes les raies suivantes ($n_p \geq 7$) sont invisibles à l'œil. Retenir : le visible s'arrête vers $400\ \text{nm}$, et la série de Balmer continue bien au-delà.
\end{attention}

\begin{exemplebox}[Limite de la série de Balmer]
La transition de plus grande énergie de la série est $\infty \to 2$ :
\[ \Delta E_{\max} = \frac{13{,}6}{2^2} = 3{,}40\ \text{eV} \quad\Longrightarrow\quad \lambda_{\min} = \frac{hc}{\Delta E_{\max}} \approx 365\ \text{nm}. \]
Toute raie de Balmer vérifie donc l'encadrement : $\boxed{365\ \text{nm} \leq \lambda \leq 656\ \text{nm}}$. Cet encadrement est un excellent réflexe de vérification au Baccalauréat : si tu trouves une raie de Balmer à 500 nm ou à 300 nm, il y a une erreur de calcul !
\end{exemplebox}

\courssub{La série de Paschen et les séries infrarouges}

La série de Paschen regroupe les transitions aboutissant à $n_i = 3$. Le niveau $n = 3$ étant proche du zéro d'énergie ($E_3 = -1{,}51\ \text{eV}$), les transitions y sont peu énergétiques : \textbf{toutes les raies de Paschen sont dans l'infrarouge}.

\begin{exemplebox}[Première raie de Paschen]
Transition $4 \to 3$ :
\[ \Delta E = 13{,}6\left(\frac{1}{9} - \frac{1}{16}\right) = 0{,}66\ \text{eV} \quad\Longrightarrow\quad \lambda = \frac{hc}{\Delta E} \approx 1875\ \text{nm}. \]
On est bien dans l'infrarouge ($\lambda > 800\ \text{nm}$). La limite de série ($\infty \to 3$) vaut $\lambda_{\min} = \dfrac{hc}{13{,}6/9\ \text{eV}} \approx 820\ \text{nm}$, également infrarouge.
\end{exemplebox}

On peut poursuivre la généralisation : la \cle{série de Brackett} ($n_i = 4$) et la \cle{série de Pfund} ($n_i = 5$) existent aussi ; toutes leurs raies sont dans l'infrarouge lointain. Elles ne sont pas exigibles au programme, mais leur existence confirme la solidité du modèle : plus le niveau d'arrivée $n_i$ est élevé, moins les transitions sont énergétiques, et plus les raies s'enfoncent dans l'infrarouge.

\begin{propriete}[Raie extrême et limite d'une série]
Pour une série de niveau final $n_i$ :
\begin{itemize}
\item la raie de \textbf{plus grande longueur d'onde} (la moins énergétique) correspond à la transition $n_p = n_i + 1 \to n_i$ ;
\item la \cle{limite de série}, c'est-à-dire la \textbf{plus petite longueur d'onde} (la plus énergétique), correspond à la transition $n_p = \infty \to n_i$, d'énergie $\Delta E = \dfrac{E_0}{n_i^2}$.
\end{itemize}
Entre ces deux bornes, les raies se resserrent de plus en plus lorsque $n_p$ augmente, car les niveaux d'énergie se rapprochent.
\end{propriete}

\courssub{Formule générale de Ritz--Rydberg}

En combinant $\Delta E = E_0\left(\dfrac{1}{n_i^2} - \dfrac{1}{n_p^2}\right)$ et $\Delta E = \dfrac{hc}{\lambda}$, on obtient directement la loi générale donnant toutes les raies de l'hydrogène :

\begin{propriete}[Formule de Ritz--Rydberg]
La longueur d'onde de toute raie de l'atome d'hydrogène vérifie :
\[ \frac{1}{\lambda} = \frac{E_0}{hc}\left(\frac{1}{n_i^2} - \frac{1}{n_p^2}\right) = R_H\left(\frac{1}{n_i^2} - \frac{1}{n_p^2}\right), \]
où $R_H = \dfrac{E_0}{hc} \approx 1{,}097\times 10^{7}\ \text{m}^{-1}$ est la \cle{constante de Rydberg}. Cette formule est admise comme conséquence directe des niveaux de Bohr ; sa démonstration n'est pas exigible, mais savoir la retrouver à partir de $\Delta E = hc/\lambda$ est un excellent réflexe.
\end{propriete}

Vérifions l'homogénéité : $E_0$ est en joules, $hc$ en $\text{J}\cdot\text{m}$, donc $\dfrac{E_0}{hc}$ est bien en $\text{m}^{-1}$, comme $\dfrac{1}{\lambda}$. La formule est cohérente.

\begin{savaistu}[Balmer avant Bohr : la science en marche]
En 1885, Johann Balmer, un modeste professeur de mathématiques suisse, découvre par tâtonnement que les longueurs d'onde visibles de l'hydrogène obéissent à la formule empirique $\lambda = \lambda_0 \dfrac{n^2}{n^2 - 4}$. Il ne comprend pas \emph{pourquoi} — personne ne le comprend. Il faudra attendre \textbf{28 ans}, et le modèle de Bohr (1913), pour que cette formule tombe naturellement de la quantification des niveaux d'énergie. Bohr a d'ailleurs raconté : « Dès que j'ai vu la formule de Balmer, tout m'est apparu clair. » Belle leçon de démarche scientifique : l'expérience et les régularités empiriques précèdent souvent la théorie qui les explique.
\end{savaistu}

\courssub{Diagramme énergétique complet et classement des séries}

\begin{popfigure}
\begin{tikzpicture}[scale=1.0, every node/.style={font=\small}]
% Niveaux d'énergie
\draw[thick, popdark] (0,0) -- (8.5,0) node[right] {$n=1$ \quad $E_1=-13{,}6$ eV};
\draw[thick, popdark] (0,2.55) -- (8.5,2.55) node[right] {$n=2$ \quad $-3{,}40$ eV};
\draw[thick, popdark] (0,3.02) -- (8.5,3.02) node[right] {$n=3$ \quad $-1{,}51$ eV};
\draw[thick, popdark] (0,3.19) -- (8.5,3.19) node[right] {$n=4$ \quad $-0{,}85$ eV};
\draw[thick, popdark] (0,3.28) -- (8.5,3.28) node[right] {$n=5$};
\draw[thick, popdark] (0,3.33) -- (8.5,3.33) node[right] {$n=6$};
\draw[thick, popdark, dashed] (0,3.4) -- (8.5,3.4) node[right] {$n=\infty$ \quad $E=0$ (ionisation)};
% Lyman (violet)
\draw[-{Stealth}, very thick, poppurple] (1.2,2.55) -- (1.2,0.05);
\draw[-{Stealth}, very thick, poppurple] (1.6,3.02) -- (1.6,0.05);
\draw[-{Stealth}, very thick, poppurple] (2.0,3.19) -- (2.0,0.05);
\draw[-{Stealth}, very thick, poppurple, dashed] (2.4,3.4) -- (2.4,0.05);
\node[poppurple, font=\small\bfseries] at (1.8,-0.45) {Lyman (UV)};
% Balmer (vert)
\draw[-{Stealth}, very thick, popgreen] (3.6,3.02) -- (3.6,2.6);
\draw[-{Stealth}, very thick, popgreen] (4.0,3.19) -- (4.0,2.6);
\draw[-{Stealth}, very thick, popgreen] (4.4,3.28) -- (4.4,2.6);
\draw[-{Stealth}, very thick, popgreen, dashed] (4.8,3.4) -- (4.8,2.6);
\node[popgreen, font=\small\bfseries] at (4.2,2.15) {Balmer (visible)};
% Paschen (rouge)
\draw[-{Stealth}, very thick, poppink] (6.2,3.19) -- (6.2,3.07);
\draw[-{Stealth}, very thick, poppink] (6.6,3.28) -- (6.6,3.07);
\draw[-{Stealth}, very thick, poppink, dashed] (7.0,3.4) -- (7.0,3.07);
\node[poppink, font=\small\bfseries] at (6.6,2.8) {Paschen (IR)};
\end{tikzpicture}
\legende{Diagramme énergétique de l'atome d'hydrogène (échelle verticale non linéaire pour la lisibilité). Les flèches violettes forment la série de Lyman, les vertes la série de Balmer, les rouges la série de Paschen. Les flèches en pointillés représentent les limites de séries ($n_p \to \infty$).}
\end{popfigure}

Ce diagramme permet de lire d'un coup d'œil un classement énergétique essentiel :

\begin{aretenir}[Classement énergétique des transitions]
Pour un même niveau de départ $n_p$, la transition vers $n = 1$ (Lyman) est \textbf{plus énergétique} que celle vers $n = 2$ (Balmer), elle-même plus énergétique que celle vers $n = 3$ (Paschen) :
\[ \Delta E_{\text{Lyman}} > \Delta E_{\text{Balmer}} > \Delta E_{\text{Paschen}} \quad\Longleftrightarrow\quad \lambda_{\text{Lyman}} < \lambda_{\text{Balmer}} < \lambda_{\text{Paschen}}. \]
En effet, plus le niveau d'arrivée est profond, plus la « chute » de l'électron est grande, et plus le photon emporte d'énergie.
\end{aretenir}

\begin{popfigure}
\begin{tikzpicture}
\begin{axis}[
  width=\linewidth, height=4.6cm,
  axis lines=middle,
  xlabel={$\lambda$ (nm)}, ylabel={},
  xmin=350, xmax=750, ymin=0, ymax=1,
  ytick=\empty,
  xtick={365,400,410,434,486,656,700},
  x tick label style={font=\scriptsize},
  xlabel style={font=\small},
  clip=false
]
% Domaine visible
\fill[popgreenL, opacity=0.6] (axis cs:400,0) rectangle (axis cs:700,0.9);
\node[font=\scriptsize, popdark] at (axis cs:550,0.78) {domaine visible};
\node[font=\scriptsize, popmuted] at (axis cs:378,0.78) {UV};
\node[font=\scriptsize, popmuted] at (axis cs:726,0.78) {IR};
% Limite de série
\draw[very thick, poporange, dashed] (axis cs:365,0) -- (axis cs:365,0.55);
\node[font=\scriptsize, poporange, anchor=south] at (axis cs:365,0.55) {limite 365 nm};
% Raies de Balmer
\draw[very thick, poppink] (axis cs:656,0) -- (axis cs:656,0.45);
\node[font=\scriptsize, poppink, anchor=south] at (axis cs:656,0.45) {H$\alpha$};
\draw[very thick, popblue] (axis cs:486,0) -- (axis cs:486,0.45);
\node[font=\scriptsize, popblue, anchor=south] at (axis cs:486,0.45) {H$\beta$};
\draw[very thick, poppurple] (axis cs:434,0) -- (axis cs:434,0.45);
\node[font=\scriptsize, poppurple, anchor=south] at (axis cs:434,0.45) {H$\gamma$};
\draw[very thick, poppurple!70!popblue] (axis cs:410,0) -- (axis cs:410,0.3);
\node[font=\scriptsize, poppurple!70!popblue, anchor=south west] at (axis cs:404,0.3) {H$\delta$};
\end{axis}
\end{tikzpicture}
\legende{Positions des raies de la série de Balmer sur l'axe des longueurs d'onde. Seules H$\alpha$, H$\beta$, H$\gamma$ et H$\delta$ tombent dans le domaine visible (zone verte) ; la limite de série (365 nm) est dans l'ultraviolet.}
\end{popfigure}

\courssub{Méthode : identifier la série d'une raie inconnue}

\begin{methode}[Identifier la série d'une raie de longueur d'onde $\lambda$]
\begin{enumerate}
\item Calculer l'énergie du photon : $\Delta E = \dfrac{hc}{\lambda}$, convertie en électronvolts (diviser par $1{,}60\times 10^{-19}$).
\item Tester successivement $n_i = 1$, puis $n_i = 2$, puis $n_i = 3$ : chercher un entier $n_p > n_i$ tel que
\[ \Delta E = 13{,}6\left(\frac{1}{n_i^2} - \frac{1}{n_p^2}\right)\ \text{eV}. \]
Astuce : isoler $\dfrac{1}{n_p^2} = \dfrac{1}{n_i^2} - \dfrac{\Delta E}{13{,}6}$ ; si le résultat est l'inverse du carré d'un entier, c'est gagné.
\item Conclure : $n_i = 1$ donne Lyman, $n_i = 2$ donne Balmer, $n_i = 3$ donne Paschen.
\item Vérifier la cohérence avec le domaine spectral : Lyman $\to$ UV, Balmer $\to$ visible (4 raies) ou UV proche, Paschen $\to$ IR.
\end{enumerate}
\end{methode}

\begin{exoresolu}[À quelle série appartient la raie de 486 nm ?]
Énoncé : Une lampe à hydrogène émet une raie bleu-vert de longueur d'onde $\lambda = 486\ \text{nm}$. Déterminer la transition correspondante et la série à laquelle elle appartient. On donne $h = 6{,}63\times 10^{-34}\ \text{J}\cdot\text{s}$, $c = 3{,}00\times 10^{8}\ \text{m/s}$, $1\ \text{eV} = 1{,}60\times 10^{-19}\ \text{J}$.
\tcblower
\textbf{Étape 1 : énergie du photon.}
\[ \Delta E = \frac{hc}{\lambda} = \frac{6{,}63\times 10^{-34}\times 3{,}00\times 10^{8}}{486\times 10^{-9}} = 4{,}09\times 10^{-19}\ \text{J} = \frac{4{,}09\times 10^{-19}}{1{,}60\times 10^{-19}} \approx 2{,}56\ \text{eV}. \]
\textbf{Étape 2 : test des séries.}
\begin{itemize}
\item Lyman ($n_i = 1$) : la raie la \emph{moins} énergétique vaut déjà 10,2 eV $\gg$ 2,56 eV. Impossible.
\item Balmer ($n_i = 2$) : $\dfrac{1}{n_p^2} = \dfrac{1}{4} - \dfrac{2{,}56}{13{,}6} = 0{,}250 - 0{,}188 = 0{,}062 \approx \dfrac{1}{16}$, donc $n_p = 4$. Parfait !
\end{itemize}
\textbf{Étape 3 : conclusion.} La raie correspond à la transition $4 \to 2$ : c'est la raie \textbf{H$\beta$} de la \textbf{série de Balmer}. La couleur bleu-vert et la position dans le visible confirment le résultat.
\end{exoresolu}

\begin{exercice}[Série et limite]
\begin{enumerate}
\item Calculer la longueur d'onde de la raie $5 \to 3$ de l'hydrogène. À quelle série appartient-elle et dans quel domaine spectral se situe-t-elle ?
\item Déterminer la longueur d'onde de la limite de la série de Paschen.
\item Une raie de l'hydrogène a pour longueur d'onde $\lambda = 95\ \text{nm}$. Identifier la transition correspondante.
\end{enumerate}
\end{exercice}

\begin{corrige}[Exercice : Série et limite]
\textbf{1.} $\Delta E = 13{,}6\left(\dfrac{1}{9}-\dfrac{1}{25}\right) = 13{,}6\times 0{,}0711 \approx 0{,}97\ \text{eV}$, d'où $\lambda = \dfrac{hc}{\Delta E} \approx 1280\ \text{nm}$. Transition vers $n_i = 3$ : série de \textbf{Paschen}, domaine \textbf{infrarouge}. \par
\textbf{2.} Limite de Paschen : $\Delta E = \dfrac{13{,}6}{9} = 1{,}51\ \text{eV}$, soit $\lambda_{\min} = \dfrac{hc}{\Delta E} \approx 820\ \text{nm}$ (infrarouge proche). \par
\textbf{3.} $\Delta E = \dfrac{hc}{\lambda} = \dfrac{6{,}63\times 10^{-34}\times 3{,}00\times 10^{8}}{95\times 10^{-9}} \approx 2{,}09\times 10^{-18}\ \text{J} \approx 13{,}1\ \text{eV}$. C'est proche de 13,6 eV : série de Lyman. On teste $n_i=1$ : pour $n_p=5$, $\Delta E = 13{,}6\left(1-\tfrac{1}{25}\right) = 13{,}06\ \text{eV}$, ce qui donne $\lambda \approx 95\ \text{nm}$. Il s'agit bien de la transition $5 \to 1$ (série de Lyman, ultraviolet).
\end{corrige}

\begin{aretenir}[L'essentiel sur les séries spectrales]
\begin{itemize}
\item Une série regroupe toutes les transitions vers un même niveau final : \textbf{Lyman} ($\to 1$, UV), \textbf{Balmer} ($\to 2$, 4 raies visibles), \textbf{Paschen} ($\to 3$, IR).
\item Formule générale : $\dfrac{1}{\lambda} = \dfrac{E_0}{hc}\left(\dfrac{1}{n_i^2} - \dfrac{1}{n_p^2}\right)$.
\item Raie la plus longue d'une série : $n_p = n_i + 1$ ; limite de série : $n_p = \infty$.
\item Encadrements utiles : Lyman $\lambda \geq 91{,}2\ \text{nm}$ ; Balmer $365\ \text{nm} \leq \lambda \leq 656\ \text{nm}$ ; Paschen $\lambda \geq 820\ \text{nm}$.
\end{itemize}
\end{aretenir}

\coursec{Méthodes de résolution et applications type Bac}

Dans la section précédente, tu as appris à classer les raies de l'hydrogène en séries : Lyman (retour vers $n=1$, ultraviolet), Balmer (retour vers $n=2$, visible), Paschen (retour vers $n=3$, infrarouge). Tu sais désormais \emph{comprendre} le diagramme énergétique. Reste à transformer cette compréhension en \cle{points au Bac}. C'est l'objet de cette dernière section : te donner des \cle{méthodes} systématiques, applicables les yeux fermés, pour chaque type de question que l'on rencontre dans les sujets camerounais. La bonne nouvelle : ces questions sont parmi les plus mécaniques de tout le programme de Terminale C.

\courssub{Les quatre formules clés et la gestion des unités}

Toutes les questions du Bac sur l'atome d'hydrogène se ramènent à quatre relations. Mémorise-les avec leurs unités, c'est la moitié du travail.

\begin{methode}[Fiche récapitulative des quatre formules clés]
\begin{enumerate}
\item \textbf{Énergie du niveau $n$} : $E_n = -\dfrac{13{,}6}{n^2}$ (en eV), avec $n \in \mathbb{N}^*$.
\item \textbf{Énergie d'une transition} du niveau $n_i$ vers le niveau $n_p$ (avec $n_i > n_p$ pour une émission) :
\[ |\Delta E| = 13{,}6\left(\frac{1}{n_p^2} - \frac{1}{n_i^2}\right) \quad (\text{en eV}). \]
\item \textbf{Longueur d'onde associée} : $\lambda(\text{nm}) = \dfrac{1240}{\Delta E(\text{eV})}$, car $hc = 1{,}24\times 10^{-6}\ \text{eV}\cdot\text{m} = 1240\ \text{eV}\cdot\text{nm}$.
\item \textbf{Énergie d'ionisation} depuis le niveau $n$ : $E_i(n) = \dfrac{13{,}6}{n^2}$ (en eV).
\end{enumerate}
\end{methode}

\begin{attention}[Le piège n°1 : mélanger les unités]
L'erreur la plus fréquente au Bac consiste à écrire $\lambda = \dfrac{hc}{\Delta E}$ avec $hc$ en $\text{J}\cdot\text{m}$ et $\Delta E$ en eV, ou bien $\lambda$ en nm avec $hc$ en $\text{J}\cdot\text{m}$. Le résultat est alors faux de plusieurs ordres de grandeur. \textbf{Règle d'or : choisis UN système cohérent et convertis AVANT le calcul.}
\begin{itemize}
\item Système « électron-volt » (recommandé, plus rapide) : $\Delta E$ en eV, $hc = 1240\ \text{eV}\cdot\text{nm}$, $\lambda$ obtenue directement en nm.
\item Système international : $\Delta E$ en joules ($1\ \text{eV} = 1{,}60\times 10^{-19}\ \text{J}$), $h = 6{,}63\times 10^{-34}\ \text{J}\cdot\text{s}$, $c = 3{,}00\times 10^{8}\ \text{m}\cdot\text{s}^{-1}$, $\lambda$ obtenue en m puis convertie en nm ($1\ \text{nm} = 10^{-9}\ \text{m}$).
\end{itemize}
Vérifie toujours la cohérence : une raie visible doit donner $\lambda$ entre 400 et 800 nm environ.
\end{attention}

\courssub{Méthode type 1 : calculer l'énergie d'un niveau donné}

\textbf{Quand l'utiliser ?} L'énoncé demande « l'énergie du niveau $n$ », « l'énergie de l'état fondamental », « l'énergie de l'état excité $n = 3$ »...

\begin{methode}[Démarche]
\begin{enumerate}
\item Identifier la valeur de $n$ dans l'énoncé.
\item Appliquer directement $E_n = -\dfrac{13{,}6}{n^2}$ (eV).
\item Ne pas oublier le \textbf{signe moins} : les états liés ont une énergie négative.
\item Convertir en joules si demandé : $E(\text{J}) = E(\text{eV}) \times 1{,}60\times 10^{-19}$.
\end{enumerate}
\end{methode}

\begin{exemplebox}[Énergie du niveau $n = 3$]
\[ E_3 = -\frac{13{,}6}{3^2} = -\frac{13{,}6}{9} = -1{,}51\ \text{eV}. \]
En joules : $E_3 = -1{,}51 \times 1{,}60\times 10^{-19} = -2{,}42\times 10^{-19}\ \text{J}$.
\end{exemplebox}

\courssub{Méthode type 2 : calculer la longueur d'onde d'une transition donnée}

\textbf{Quand l'utiliser ?} L'énoncé donne les deux niveaux (par exemple « l'atome passe de $n = 4$ à $n = 2$ ») et demande la longueur d'onde du photon émis ou absorbé.

\begin{methode}[Démarche]
\begin{enumerate}
\item Repérer $n_i$ (niveau de départ) et $n_p$ (niveau d'arrivée). Pour une \cle{émission}, $n_i > n_p$ ; pour une \cle{absorption}, $n_i < n_p$.
\item Calculer $|\Delta E| = 13{,}6\left(\dfrac{1}{n_p^2} - \dfrac{1}{n_i^2}\right)$ en eV.
\item En déduire $\lambda = \dfrac{1240}{|\Delta E|}$ en nm.
\item Conclure sur le domaine spectral (UV : $\lambda < 400$ nm ; visible : $400 \lesssim \lambda \lesssim 800$ nm ; IR : $\lambda > 800$ nm) et sur la série concernée.
\end{enumerate}
\end{methode}

\begin{exoresolu}[La raie H$_\beta$ : transition $n = 4 \to n = 2$]
Un atome d'hydrogène, initialement au niveau $n = 4$, se désexcite vers le niveau $n = 2$. Calculer l'énergie du photon émis, sa longueur d'onde, et préciser le domaine spectral et la série.
\tcblower
\textbf{Étape 1 — Énergie de la transition.}
\begin{align*}
|\Delta E| &= 13{,}6\left(\frac{1}{n_p^2} - \frac{1}{n_i^2}\right) = 13{,}6\left(\frac{1}{2^2} - \frac{1}{4^2}\right) \\
&= 13{,}6\left(\frac{1}{4} - \frac{1}{16}\right) = 13{,}6 \times \frac{3}{16} = 2{,}55\ \text{eV}.
\end{align*}
\textbf{Étape 2 — Longueur d'onde.}
\[ \lambda = \frac{1240}{2{,}55} = 486\ \text{nm}. \]
\textbf{Étape 3 — Conclusion.} $\lambda = 486$ nm appartient au domaine visible (raie bleu-vert) : c'est la raie \cle{H$_\beta$} de la série de \cle{Balmer} (arrivée sur $n = 2$). Vérification de cohérence : le résultat est bien entre 400 et 800 nm, conforme à une raie visible.
\end{exoresolu}

\courssub{Méthode type 3 : identifier une transition à partir d'une raie observée}

\textbf{Quand l'utiliser ?} C'est le problème « inverse » : l'énoncé donne $\lambda$ (ou l'énergie du photon) et demande de retrouver les niveaux concernés.

\begin{methode}[Démarche]
\begin{enumerate}
\item Calculer l'énergie du photon : $\Delta E = \dfrac{1240}{\lambda(\text{nm})}$ en eV.
\item Exploiter les indices de l'énoncé : si la raie est visible, c'est Balmer donc $n_p = 2$ ; si UV, Lyman donc $n_p = 1$ ; si IR, Paschen donc $n_p = 3$.
\item Résoudre $\Delta E = 13{,}6\left(\dfrac{1}{n_p^2} - \dfrac{1}{n_i^2}\right)$ pour trouver $n_i$ entier.
\item \textbf{Vérifier impérativement} que $n_i$ est un entier supérieur à $n_p$ ; sinon, revoir l'hypothèse sur $n_p$.
\end{enumerate}
\end{methode}

\begin{exoresolu}[Identifier la transition d'une raie rouge à 656 nm]
On observe dans le spectre d'émission de l'hydrogène une raie rouge de longueur d'onde $\lambda = 656$ nm. Déterminer les niveaux entre lesquels s'effectue la transition.
\tcblower
\textbf{Étape 1 — Énergie du photon.} $\Delta E = \dfrac{1240}{656} = 1{,}89\ \text{eV}$.

\textbf{Étape 2 — Hypothèse.} La raie est visible : série de Balmer, donc $n_p = 2$.

\textbf{Étape 3 — Résolution.}
\[ 1{,}89 = 13{,}6\left(\frac{1}{4} - \frac{1}{n_i^2}\right) \implies \frac{1}{n_i^2} = \frac{1}{4} - \frac{1{,}89}{13{,}6} = 0{,}250 - 0{,}139 = 0{,}111 \approx \frac{1}{9}. \]
Donc $n_i^2 = 9$, soit $n_i = 3$ : entier, l'hypothèse est validée.

\textbf{Conclusion :} transition $n = 3 \to n = 2$, c'est la raie \cle{H$_\alpha$} (rouge) de la série de Balmer.
\end{exoresolu}

\courssub{Méthode type 4 : énergie d'ionisation depuis un état excité}

\textbf{Quand l'utiliser ?} L'énoncé demande « l'énergie minimale pour ioniser l'atome depuis l'état $n$ » ou « la longueur d'onde maximale d'un photon capable d'ioniser l'atome depuis $n$ ».

Ioniser, c'est arracher l'électron et l'amener du niveau $n$ au niveau $n \to \infty$ où $E_\infty = 0$. Le bilan énergétique donne immédiatement :
\[ E_i(n) = E_\infty - E_n = 0 - \left(-\frac{13{,}6}{n^2}\right) = \frac{13{,}6}{n^2}\ \text{(eV)}. \]

\begin{exemplebox}[Ionisation depuis $n = 2$]
$E_i(2) = \dfrac{13{,}6}{4} = 3{,}40\ \text{eV}$. Un photon ionisant doit vérifier $\lambda \leq \dfrac{1240}{3{,}40} = 365\ \text{nm}$ (ultraviolet). Remarque : c'est exactement la \emph{limite de la série de Balmer}.
\end{exemplebox}

\begin{attention}[Énergie d'ionisation ≠ énergie du niveau]
L'énergie d'ionisation est \textbf{positive} (il faut \emph{fournir} de l'énergie) : $E_i(n) = +13{,}6/n^2$, alors que l'énergie du niveau est négative : $E_n = -13{,}6/n^2$. Depuis le fondamental, $E_i(1) = 13{,}6\ \text{eV}$ : c'est la valeur maximale, car l'électron y est le plus lié.
\end{attention}

\courssub{Méthode type 5 : cascade de désexcitation et dénombrement des raies}

\textbf{Quand l'utiliser ?} L'énoncé dit : « un ensemble d'atomes est porté au niveau $n$ ; combien de raies peut-on observer lors de la désexcitation ? »

Un atome au niveau $n$ peut redescendre directement vers le fondamental, ou par étapes (cascade) : $n \to n-1 \to \dots \to 1$. Chaque « saut » possible entre deux niveaux correspond à une raie distincte. Le nombre de raies est donc le nombre de \cle{couples} de niveaux $(n_i, n_p)$ avec $n_i > n_p$ parmi les niveaux $1, 2, \dots, n$.

\begin{propriete}[Nombre de raies d'une désexcitation depuis le niveau $n$]
Le nombre de raies distinctes émises lors de la désexcitation d'une population d'atomes d'hydrogène excités au niveau $n$ est :
\[ N = \frac{n(n-1)}{2}. \]
\emph{Dénombrement :} on choisit 2 niveaux distincts parmi $n$ (le couple définit une transition unique, le plus haut étant le niveau de départ) : $N = \binom{n}{2} = \dfrac{n(n-1)}{2}$. Par exemple, depuis $n = 4$ : $N = \dfrac{4 \times 3}{2} = 6$ raies.
\end{propriete}

\begin{popfigure}
\begin{tikzpicture}[scale=1.0, font=\small]
% Niveaux d'énergie
\draw[thick, popdark] (0,0) -- (6.5,0) node[right] {$n=1$ \quad $E_1 = -13{,}6$ eV};
\draw[thick, popdark] (0,1.7) -- (6.5,1.7) node[right] {$n=2$ \quad $E_2 = -3{,}40$ eV};
\draw[thick, popdark] (0,2.6) -- (6.5,2.6) node[right] {$n=3$ \quad $E_3 = -1{,}51$ eV};
\draw[thick, popdark] (0,3.15) -- (6.5,3.15) node[right] {$n=4$ \quad $E_4 = -0{,}85$ eV};
% Transitions depuis 4
\draw[-{Stealth}, very thick, popblue] (0.8,3.15) -- (0.8,1.75) node[midway, left, popblue] {\scriptsize $4\to2$ : 486 nm};
\draw[-{Stealth}, very thick, poppurple] (1.6,3.15) -- (1.6,0.05) node[midway, left, poppurple] {\scriptsize $4\to1$ : 97 nm};
\draw[-{Stealth}, very thick, poporange] (2.4,3.15) -- (2.4,2.65) node[midway, right, poporange] {\scriptsize $4\to3$ : 1875 nm};
% Transitions depuis 3
\draw[-{Stealth}, very thick, popgreen] (3.2,2.6) -- (3.2,1.75) node[midway, right, popgreen] {\scriptsize $3\to2$ : 656 nm};
\draw[-{Stealth}, very thick, poppink] (4.0,2.6) -- (4.0,0.05) node[pos=0.6, right, poppink] {\scriptsize $3\to1$ : 103 nm};
% Transition depuis 2
\draw[-{Stealth}, very thick, popgold] (4.8,1.7) -- (4.8,0.05) node[midway, right, popgold] {\scriptsize $2\to1$ : 122 nm};
\end{tikzpicture}
\legende{Les six raies possibles d'une désexcitation depuis $n = 4$ : trois appartiennent à Lyman (UV), deux à Balmer (visible : 486 nm et 656 nm), une à Paschen (IR).}
\end{popfigure}

\courssub{L'organigramme de décision : quelle formule pour quelle question ?}

Face à un énoncé de Bac, la première chose à faire est d'identifier \emph{ce que l'on te demande}. L'arbre de décision ci-dessous résume toute la stratégie.

\begin{popfigure}
\begin{tikzpicture}[
  font=\small,
  decision/.style={rectangle, rounded corners, draw=popblue, fill=popblueL, text width=4.2cm, align=center, minimum height=1cm},
  formule/.style={rectangle, draw=poporange, fill=poporangeL, text width=4.6cm, align=center, minimum height=0.9cm},
  node distance=0.55cm and 1.2cm]
\node[decision] (q) {\textbf{Que demande l'énoncé ?}};
\node[decision, below left=0.7cm and 2.2cm of q] (niv) {Énergie d'un niveau $n$ ?};
\node[decision, below=0.7cm of q] (trans) {Longueur d'onde d'une transition $n_i \to n_p$ ?};
\node[decision, below right=0.7cm and 2.2cm of q] (ion) {Ionisation depuis $n$ ? Nombre de raies ?};
\node[formule, below=0.5cm of niv] (f1) {$E_n = -\dfrac{13{,}6}{n^2}$ (eV)};
\node[formule, below=0.5cm of trans] (f2) {$|\Delta E| = 13{,}6\left(\dfrac{1}{n_p^2}-\dfrac{1}{n_i^2}\right)$ puis $\lambda = \dfrac{1240}{|\Delta E|}$};
\node[formule, below=0.5cm of ion] (f3) {$E_i = \dfrac{13{,}6}{n^2}$ \quad ; \quad $N = \dfrac{n(n-1)}{2}$};
\draw[-{Stealth}, thick, popdark] (q) -- (niv);
\draw[-{Stealth}, thick, popdark] (q) -- (trans);
\draw[-{Stealth}, thick, popdark] (q) -- (ion);
\draw[-{Stealth}, thick, popdark] (niv) -- (f1);
\draw[-{Stealth}, thick, popdark] (trans) -- (f2);
\draw[-{Stealth}, thick, popdark] (ion) -- (f3);
\end{tikzpicture}
\legende{Arbre de décision : identifier la question, appliquer la formule correspondante, convertir les unités avant tout calcul.}
\end{popfigure}

\courssub{La rédaction attendue au Bac}

Un correcteur attend une copie structurée. Voici le squelette de rédaction qui garantit tous les points :

\begin{methode}[Structure de rédaction type Bac]
\begin{enumerate}
\item \textbf{Schéma} : tracer le diagramme énergétique (même simplifié : deux ou trois niveaux horizontaux), y placer la flèche de transition dans le bon sens (vers le bas : émission ; vers le haut : absorption).
\item \textbf{Expression littérale} : écrire la relation utilisée avec des symboles, par exemple $|\Delta E| = E_0\left(\dfrac{1}{n_p^2} - \dfrac{1}{n_i^2}\right)$, avant tout chiffre.
\item \textbf{Application numérique} : remplacer, calculer, donner le résultat avec son \textbf{unité} et un nombre raisonnable de chiffres significatifs (2 ou 3).
\item \textbf{Conclusion physique} : nommer la raie, la série, le domaine spectral (UV, visible, IR), et dire si le photon est émis ou absorbé.
\end{enumerate}
\end{methode}

\begin{exoresolu}[Sujet type Bac : synthèse complète]
Des atomes d'hydrogène sont excités au niveau $n = 3$.
\begin{enumerate}
\item Calculer l'énergie du niveau $n = 3$.
\item Déterminer le nombre de raies observables lors de la désexcitation et calculer la longueur d'onde de la raie de plus grande énergie.
\item Quelle énergie faut-il fournir pour ioniser un atome se trouvant au niveau $n = 3$ ?
\end{enumerate}
\tcblower
\textbf{1.} $E_3 = -\dfrac{13{,}6}{9} = -1{,}51\ \text{eV}$.

\textbf{2.} Nombre de raies : $N = \dfrac{3 \times 2}{2} = 3$ raies ($3\to2$, $3\to1$, $2\to1$). La raie de plus grande énergie correspond au plus grand écart : la transition $3 \to 1$.
\[ |\Delta E| = 13{,}6\left(1 - \frac{1}{9}\right) = 13{,}6 \times \frac{8}{9} = 12{,}1\ \text{eV} \implies \lambda = \frac{1240}{12{,}1} = 103\ \text{nm}. \]
C'est une raie ultraviolette de la série de Lyman.

\textbf{3.} $E_i(3) = \dfrac{13{,}6}{9} = 1{,}51\ \text{eV}$, soit $2{,}42\times 10^{-19}\ \text{J}$. Tout photon d'énergie supérieure ou égale à cette valeur ($\lambda \leq 821$ nm) peut ioniser l'atome depuis $n = 3$.
\end{exoresolu}

\begin{exercice}[Pour t'entraîner]
\begin{enumerate}
\item Calculer la longueur d'onde de la transition $n = 5 \to n = 2$. À quelle série appartient-elle ? Est-elle visible ?
\item Une raie ultraviolette de longueur d'onde $\lambda = 95$ nm est observée. Identifier la transition.
\item Des atomes sont excités au niveau $n = 5$. Combien de raies peut-on observer ? Quelle est la plus grande longueur d'onde émise ?
\end{enumerate}
\end{exercice}

\begin{corrige}[Exercice n°1]
\textbf{1.} $|\Delta E| = 13{,}6\left(\dfrac{1}{4} - \dfrac{1}{25}\right) = 13{,}6 \times 0{,}21 = 2{,}86\ \text{eV}$ ; $\lambda = \dfrac{1240}{2{,}86} = 434$ nm : raie H$_\gamma$ de Balmer, visible (violet).

\textbf{2.} $\Delta E = \dfrac{1240}{95} = 13{,}05\ \text{eV}$. UV donc $n_p = 1$ : $\dfrac{1}{n_i^2} = 1 - \dfrac{13{,}05}{13{,}6} = 0{,}040 \approx \dfrac{1}{25}$, donc $n_i = 5$. Transition $5 \to 1$ (Lyman).

\textbf{3.} $N = \dfrac{5 \times 4}{2} = 10$ raies. La plus grande longueur d'onde correspond au plus petit écart : $5 \to 4$ : $|\Delta E| = 13{,}6\left(\dfrac{1}{16} - \dfrac{1}{25}\right) = 0{,}306\ \text{eV}$, $\lambda = \dfrac{1240}{0{,}306} = 4050$ nm (infrarouge, série de Brackett).
\end{corrige}

\begin{savaistu}[Des points « mécaniques » à ne pas laisser passer]
Au Baccalauréat C au Cameroun, les questions sur les niveaux d'énergie de l'hydrogène rapportent typiquement \textbf{3 à 5 points} et figurent parmi les plus « mécaniques » de l'épreuve : les mêmes quatre formules, les mêmes conversions, les mêmes rédactions reviennent d'année en année. Un élève qui s'est entraîné sur une dizaine d'exercices de ce type est quasiment assuré de les obtenir. C'est un investissement exceptionnel : quelques heures d'entraînement pour des points garantis le jour J. Alors, refais les exemples de cette section sans regarder les corrigés — c'est ainsi qu'on transforme une méthode en réflexe.
\end{savaistu}

\newpage

\coursec{Activité d'intégration}

\coursec{Niveaux d'énergie de l'atome d'hydrogène}

\courssub{Le photon et la quantification de l'énergie lumineuse}

\begin{definition}[Photon et relation de Planck-Einstein]
La lumière se propage sous forme de \cle{photons}, particules d'énergie $E$ liée à la fréquence $\nu$ (ou à la longueur d'onde $\lambda$) par la \textbf{relation de Planck-Einstein} :
\[
E = h\nu = \frac{hc}{\lambda}
\]
où $h = \num{6,63e-34}\ \mathrm{J\,s}$ (constante de Planck) et $c = \num{3,00e8}\ \mathrm{m\,s^{-1}}$ (célérité de la lumière dans le vide).
L'énergie d'un photon est \textbf{quantifiée} : elle ne peut prendre que des valeurs discrètes proportionnelles à la fréquence.
\end{definition}

\begin{propriete}[Formule pratique en électronique-volt et nanomètre]
En utilisant $1\ \mathrm{eV} = \num{1,602e-19}\ \mathrm{J}$, on obtient $hc \approx \num{1240}\ \mathrm{eV\,nm}$. La relation devient :
\[
\boxed{\lambda\ (\mathrm{nm}) = \frac{1240}{\Delta E\ (\mathrm{eV})}}
\]
Cette formule est très commode pour la spectroscopie atomique.
\end{propriete}

\courssub{Niveaux d'énergie de l'atome d'hydrogène}

\begin{propriete}[Modèle de Bohr — Énergies des états liés]
Dans l'atome d'hydrogène, l'électron ne peut occuper que des \cle{états stationnaires} d'énergie quantifiée :
\[
\boxed{E_n = -\frac{E_0}{n^2} \qquad (n \in \mathbb{N}^*)}
\]
avec $E_0 = 13,6\ \mathrm{eV}$ (énergie d'ionisation).
\begin{itemize}
\item $n = 1$ : \cle{état fondamental} ($E_1 = -13,6\ \mathrm{eV}$), le plus stable.
\item $n \geq 2$ : \cle{états excités} (énergies moins négatives).
\item $n \to \infty$ : \cle{limite d'ionisation} ($E_\infty = 0$), l'électron est libre au repos.
\end{itemize}
\emph{Analyse dimensionnelle :} $E_0$ en eV, $n$ sans dimension $\Rightarrow E_n$ en eV. \checkmark
\end{propriete}

\begin{aretenir}[Diagramme énergétique]
\begin{itemize}
\item Les niveaux sont \textbf{négatifs} (états liés) et se \textbf{resserrent} quand $n$ augmente ($\Delta E_{n\to n+1} \propto 1/n^3$).
\item L'échelle d'énergie est orientée vers le haut : $E=0$ en haut (ionisation), $E_1$ en bas (fondamental).
\item Une flèche \textbf{vers le bas} ($n_i > n_f$) = \textbf{émission} d'un photon d'énergie $\Delta E = E_{n_i} - E_{n_f} > 0$.
\item Une flèche \textbf{vers le haut} ($n_i < n_f$) = \textbf{absorption} d'un photon d'énergie $\Delta E = E_{n_f} - E_{n_i} > 0$.
\end{itemize}
\end{aretenir}

\courssub{Transitions, spectres et séries spectrales}

\begin{definition}[Spectres d'émission et d'absorption]
\begin{itemize}
\item \textbf{Spectre d'émission :} atome excité $\to$ désexcitation $\to$ émission de photons de longueurs d'onde \textbf{caractéristiques} (raies brillantes sur fond noir).
\item \textbf{Spectre d'absorption :} lumière blanche traverse un gaz froid $\to$ atomes absorbent les photons de certaines $\lambda$ $\to$ \textbf{raies sombres} sur fond continu aux mêmes positions que les raies d'émission.
\end{itemize}
\end{definition}

\begin{propriete}[Séries spectrales de l'hydrogène]
Les transitions se regroupent en \cle{séries} selon le niveau d'arrivée $n_f$ :
\begin{center}
\begin{tabularx}{\linewidth}{|l|X|X|X|}
\hline
\rowcolor{popblueL} \textbf{Série} & \textbf{Niveau d'arrivée $n_f$} & \textbf{Domaine spectral principal} & \textbf{Transitions emblématiques} \\
\hline
Lyman & $1$ & Ultraviolet (UV) $< 400\ \mathrm{nm}$ & $2\to1\ (122\ \mathrm{nm}),\ 3\to1\ (103\ \mathrm{nm})$ \\
\hline
Balmer & $2$ & Visible ($400$--$800\ \mathrm{nm}$) + proche UV & $\mathrm{H}_\alpha\ 3\to2\ (656\ \mathrm{nm}),\ \mathrm{H}_\beta\ 4\to2\ (486\ \mathrm{nm}),\ \mathrm{H}_\gamma\ 5\to2\ (434\ \mathrm{nm})$ \\
\hline
Paschen & $3$ & Infrarouge (IR) $> 800\ \mathrm{nm}$ & $4\to3\ (1875\ \mathrm{nm}),\ 5\to3\ (1282\ \mathrm{nm})$ \\
\hline
\end{tabularx}
\end{center}
\end{propriete}

\courssub{Activité d'intégration : La carte d'identité lumineuse de l'hydrogène}

\courssub{Objectifs de l'activité}

À l'issue de cette activité d'intégration, tu seras capable de :
\begin{itemize}
\item construire le \cle{diagramme énergétique} complet de l'atome d'hydrogène à partir de la relation $E_n = -\dfrac{E_0}{n^2}$ ;
\item calculer les \cle{longueurs d'onde} des transitions entre niveaux à l'aide de la relation de Planck-Einstein $E = h\nu = \dfrac{hc}{\lambda}$ ;
\item rattacher chaque raie à sa \cle{série spectrale} (Lyman, Balmer, Paschen) et à son domaine (ultraviolet, visible, infrarouge) ;
\item identifier un \cle{spectre inconnu} en comparant les positions des raies observées aux valeurs calculées ;
\item interpréter qualitativement le décalage des raies spectrales d'un astre lointain (effet Doppler).
\end{itemize}

Cette activité mobilise l'ensemble des savoir-faire de la leçon dans une situation authentique d'astrophysicien : l'analyse spectrale est l'outil principal d'identification de la composition chimique des étoiles.

\courssub{Situation}

Le lycée scientifique de Ngaoundéré organise sa semaine des sciences. Le club d'astronomie, encadré par le professeur de physique, prépare une exposition intitulée \emph{« Lire la lumière des étoiles »}. Le professeur explique à ses élèves que chaque élément chimique possède une \cle{carte d'identité lumineuse} : un ensemble de raies colorées dont les positions sont uniques. L'hydrogène, élément le plus abondant de l'Univers (il constitue environ $75\,\%$ de la matière ordinaire), est le premier élément dont le spectre a été expliqué par la physique quantique, grâce au modèle de Bohr (1913).

Le professeur confie à chaque groupe de 3 ou 4 élèves une mission : \textbf{construire la carte d'identité lumineuse complète de l'hydrogène}, puis l'utiliser pour authentifier un \emph{spectre mystère} photographié par un télescope amateur lors de l'observation d'une nébuleuse, et enfin discuter un cas réel d'astrophysique.

\textbf{Données fournies à chaque groupe :}
\begin{itemize}
\item constante de Planck : $h = \num{6,626e-34}\ \mathrm{J\,s}$ ;
\item célérité de la lumière dans le vide : $c = \num{3,00e8}\ \mathrm{m\,s^{-1}}$ ;
\item charge élémentaire : $e = \num{1,602e-19}\ \mathrm{C}$, donc $1\ \mathrm{eV} = \num{1,602e-19}\ \mathrm{J}$ ;
\item énergie d'ionisation de l'hydrogène : $E_0 = 13,6\ \mathrm{eV}$ ;
\item énergies des niveaux : $E_n = -\dfrac{E_0}{n^2}$ avec $n \in \mathbb{N}^*$ ;
\item limites du spectre visible : de $400\ \mathrm{nm}$ (violet) à $800\ \mathrm{nm}$ (rouge) ; en dessous de $400\ \mathrm{nm}$ : ultraviolet (UV) ; au-dessus de $800\ \mathrm{nm}$ : infrarouge (IR).
\end{itemize}

\textbf{Le spectre mystère} est une bande photographique présentant trois raies brillantes nettes aux longueurs d'onde suivantes :
\[
\lambda_a = 656\ \mathrm{nm} \qquad \lambda_b = 486\ \mathrm{nm} \qquad \lambda_c = 434\ \mathrm{nm}.
\]

\textbf{Matériel par groupe :} papier millimétré, règle, calculatrice, feutres de couleur, grande bande de papier pour l'axe spectral.

\courssub{Consignes}

\textbf{Travail en groupe (1 h 30), restitution orale de 5 min par groupe avec le diagramme affiché au tableau.}

\begin{enumerate}
\item \textbf{Calcul des niveaux d'énergie.} Calculer, en électronvolts, les énergies $E_1$, $E_2$, $E_3$, $E_4$ et $E_5$ des cinq premiers niveaux de l'atome d'hydrogène. Présenter les résultats dans un tableau avec trois chiffres significatifs.
\item \textbf{Diagramme énergétique.} Sur papier millimétré, tracer le diagramme énergétique à l'échelle (échelle conseillée : $1\ \mathrm{cm}$ pour $1\ \mathrm{eV}$). Représenter les niveaux par des segments horizontaux, l'état fondamental, les états excités, et le niveau d'ionisation $E = 0$. Que remarque-t-on sur l'espacement des niveaux quand $n$ augmente ?
\item \textbf{Calcul des longueurs d'onde.} Pour chacune des transitions suivantes : $2 \to 1$, $3 \to 2$, $4 \to 2$, $5 \to 2$ et $4 \to 3$, calculer la variation d'énergie $|\Delta E|$ en eV puis en joules, et en déduire la longueur d'onde $\lambda$ du photon émis, exprimée en nanomètres. On pourra vérifier au préalable que $\lambda\ (\mathrm{nm}) = \dfrac{1240}{\Delta E\ (\mathrm{eV})}$.
\item \textbf{Axe spectral.} Sur la grande bande de papier, graduer un axe de $100\ \mathrm{nm}$ à $2000\ \mathrm{nm}$, y placer les cinq raies calculées, et colorier celles qui appartiennent au domaine visible.
\item \textbf{Classement par séries.} Rattacher chaque raie à sa série (Lyman : arrivée en $n = 1$ ; Balmer : arrivée en $n = 2$ ; Paschen : arrivée en $n = 3$) et à son domaine (UV, visible ou IR). Compléter le diagramme énergétique par des flèches verticales représentant ces transitions.
\item \textbf{Le spectre mystère.} Comparer les trois raies du spectre mystère aux valeurs calculées. Prouver qu'il s'agit bien de l'hydrogène et identifier la transition correspondant à chaque raie.
\item \textbf{Question d'ouverture.} Un astrophysicien camerounais, stagiaire à l'observatoire, observe une étoile lointaine dont le spectre montre les raies de Balmer, mais toutes décalées vers les grandes longueurs d'onde (par exemple la raie rouge apparaît à $670\ \mathrm{nm}$ au lieu de $656\ \mathrm{nm}$). Que peut-on en déduire qualitativement sur le mouvement de cette étoile ? Ce phénomène porte un nom : lequel ?
\end{enumerate}

\courssub{Ressources à mobiliser}

\begin{aretenir}[Boîte à outils de l'activité]
\begin{itemize}
\item \textbf{Quantification des niveaux :} $E_n = -\dfrac{E_0}{n^2}$ avec $E_0 = 13,6\ \mathrm{eV}$. L'état fondamental est $n = 1$ ; les états $n \geq 2$ sont excités ; $E_\infty = 0$ correspond à l'ionisation.
\item \textbf{Photon émis ou absorbé :} lors d'une transition $n_i \to n_f$, $|\Delta E| = |E_{n_f} - E_{n_i}| = h\nu = \dfrac{hc}{\lambda}$.
\item \textbf{Formule pratique :} $\lambda\ (\mathrm{nm}) = \dfrac{1240}{\Delta E\ (\mathrm{eV})}$, car $hc = \num{6,626e-34} \times \num{3,00e8} = \num{1,988e-25}\ \mathrm{J\,m} = 1240\ \mathrm{eV\,nm}$.
\item \textbf{Séries :} Lyman (arrivée $n = 1$, UV), Balmer (arrivée $n = 2$, visible et proche UV), Paschen (arrivée $n = 3$, IR).
\item \textbf{Conversion :} $1\ \mathrm{eV} = \num{1,602e-19}\ \mathrm{J}$ ; $1\ \mathrm{nm} = 10^{-9}\ \mathrm{m}$.
\end{itemize}
\end{aretenir}

\courssub{Démarche et corrigé commenté}

\begin{exoresolu}[La carte d'identité lumineuse de l'hydrogène — corrigé complet]
\textbf{Étape 1 : énergies des niveaux.} On applique $E_n = -\dfrac{13,6}{n^2}$ (en eV) :
\begin{align*}
E_1 &= -13,6\ \mathrm{eV} \\
E_2 &= -\frac{13,6}{4} = -3,40\ \mathrm{eV} \\
E_3 &= -\frac{13,6}{9} = -1,51\ \mathrm{eV} \\
E_4 &= -\frac{13,6}{16} = -0,850\ \mathrm{eV} \\
E_5 &= -\frac{13,6}{25} = -0,544\ \mathrm{eV} \\
E_\infty &= 0\ \mathrm{eV}
\end{align*}

\textbf{Étape 2 : diagramme énergétique.}
\begin{popfigure}
\begin{tikzpicture}[scale=0.7, every node/.style={font=\small}]
% Axes
\draw[->, thick] (0,0) -- (0,15) node[above] {$E$ (eV)};
\draw[->, thick] (0,0) -- (10,0) node[right] {$n$};
% Levels
\foreach \n/\E/\label in {1/-13.6/$n=1$ (fondamental), 2/-3.40/$n=2$, 3/-1.51/$n=3$, 4/-0.850/$n=4$, 5/-0.544/$n=5$} {
    \draw[thick, popdark] (1,-\E) -- (9,-\E);
    \node[left] at (0.8,-\E) {\label\ : \E\ eV};
}
% Ionization limit
\draw[dashed, thick, poporange] (1,0) -- (9,0) node[right] {$E_\infty = 0$ (ionisation)};
% Transitions - Emission (downward arrows)
% Lyman 2->1
\draw[popblue, thick, ->] (5,-3.40) -- (5,-13.6) node[midway, right] {Lyman $\alpha$};
% Balmer 3->2, 4->2, 5->2
\draw[popgreen, thick, ->] (5.5,-1.51) -- (5.5,-3.40) node[midway, right] {Balmer $\mathrm{H}_\alpha$};
\draw[popgreen, thick, ->] (6,-0.850) -- (6,-3.40) node[midway, right] {Balmer $\mathrm{H}_\beta$};
\draw[popgreen, thick, ->] (6.5,-0.544) -- (6.5,-3.40) node[midway, right] {Balmer $\mathrm{H}_\gamma$};
% Paschen 4->3
\draw[poporange, thick, ->] (7,-0.850) -- (7,-1.51) node[midway, right] {Paschen};
% Legend
\begin{scope}[shift={(0,15.5)}]
\node[popblue] {\textbullet\ Lyman (UV)};
\node[popgreen, xshift=3cm] {\textbullet\ Balmer (Visible)};
\node[poporange, xshift=6cm] {\textbullet\ Paschen (IR)};
\end{scope}
\end{tikzpicture}
\legende{Diagramme énergétique de l'atome d'hydrogène avec les transitions étudiées. Les niveaux se resserrent vers $E=0$.}
\end{popfigure}
Les niveaux se \emph{resserrent} quand $n$ augmente : l'écart $E_2 - E_1 = 10,2\ \mathrm{eV}$ est bien plus grand que $E_5 - E_4 = 0,306\ \mathrm{eV}$. Les niveaux s'accumulent vers $E = 0$ (limite d'ionisation). C'est une conséquence directe de la décroissance en $1/n^2$.

\textbf{Étape 3 : longueurs d'onde.} On utilise $\lambda = \dfrac{1240}{\Delta E}$ avec $\Delta E$ en eV.
\begin{itemize}
\item Transition $2 \to 1$ : $\Delta E = 13,6 - 3,40 = 10,2\ \mathrm{eV}$, soit $\lambda = \dfrac{1240}{10,2} = 122\ \mathrm{nm}$.
\item Transition $3 \to 2$ : $\Delta E = 3,40 - 1,51 = 1,89\ \mathrm{eV}$, soit $\lambda = \dfrac{1240}{1,89} = 656\ \mathrm{nm}$.
\item Transition $4 \to 2$ : $\Delta E = 3,40 - 0,850 = 2,55\ \mathrm{eV}$, soit $\lambda = \dfrac{1240}{2,55} = 486\ \mathrm{nm}$.
\item Transition $5 \to 2$ : $\Delta E = 3,40 - 0,544 = 2,86\ \mathrm{eV}$, soit $\lambda = \dfrac{1240}{2,86} = 434\ \mathrm{nm}$.
\item Transition $4 \to 3$ : $\Delta E = 1,51 - 0,850 = 0,661\ \mathrm{eV}$, soit $\lambda = \dfrac{1240}{0,661} = 1875\ \mathrm{nm}$.
\end{itemize}

Vérification en unités SI pour la transition $3 \to 2$ : $\Delta E = 1,89 \times \num{1,602e-19} = \num{3,03e-19}\ \mathrm{J}$, d'où
\[
\lambda = \frac{hc}{\Delta E} = \frac{\num{6,626e-34} \times \num{3,00e8}}{\num{3,03e-19}} = \num{6,56e-7}\ \mathrm{m} = 656\ \mathrm{nm}. \checkmark
\]

\textbf{Étape 4 : axe spectral.}
\begin{popfigure}
\begin{tikzpicture}[scale=0.8]
% Axis
\draw[->, thick] (0,0) -- (19,0) node[right] {$\lambda$ (nm)};
% Zones
\fill[popblue!15] (0,0.2) rectangle (1.22,0.8) node[midway, above, font=\tiny] {UV};
\fill[popgreen!15] (4,0.2) rectangle (8,0.8) node[midway, above, font=\tiny] {Visible};
\fill[poporange!15] (12,0.2) rectangle (19,0.8) node[midway, above, font=\tiny] {IR};
% Ticks and labels
\foreach \x/\label in {0/0, 1.22/122, 4/400, 6.56/656, 4.86/486, 4.34/434, 8/800, 12/1000, 18.75/1875, 19/2000} {
    \draw (\x, -0.1) -- (\x, 0.1);
    \node[below, font=\tiny, rotate=45, anchor=north east] at (\x, -0.15) {\label};
}
% Spectral lines
\draw[popblue, very thick] (1.22, 0) -- (1.22, 1.2) node[above, font=\small] {Lyman $\alpha$};
\draw[popgreen, very thick] (6.56, 0) -- (6.56, 1.2) node[above, font=\small] {$\mathrm{H}_\alpha$};
\draw[popgreen, very thick] (4.86, 0) -- (4.86, 1.2) node[above, font=\small] {$\mathrm{H}_\beta$};
\draw[popgreen, very thick] (4.34, 0) -- (4.34, 1.2) node[above, font=\small] {$\mathrm{H}_\gamma$};
\draw[poporange, very thick] (18.75, 0) -- (18.75, 1.2) node[above, font=\small] {Paschen $\alpha$};
% Visible color bar
\shade[left color=violet, right color=red] (4, -0.5) rectangle (8, -0.3);
\node[below, font=\tiny] at (6, -0.6) {400 nm \quad 500 \quad 600 \quad 700 \quad 800 nm};
\end{tikzpicture}
\legende{Axe spectral des raies de l'hydrogène calculées. Les raies de Balmer ($\mathrm{H}_\alpha, \mathrm{H}_\beta, \mathrm{H}_\gamma$) sont dans le visible.}
\end{popfigure}
Les raies à $656$, $486$ et $434\ \mathrm{nm}$ tombent dans le visible ($400$ à $800\ \mathrm{nm}$) : on les colorie respectivement en rouge, en bleu-vert et en violet. La raie à $122\ \mathrm{nm}$ est dans l'UV, celle à $1875\ \mathrm{nm}$ dans l'IR : invisibles à l'œil.

\textbf{Étape 5 : classement par séries.}
\begin{center}
\begin{tabularx}{\linewidth}{|l|X|X|X|X|}
\hline
\rowcolor{popblueL} \textbf{Transition} & \textbf{$\Delta E$ (eV)} & \textbf{$\lambda$ (nm)} & \textbf{Série} & \textbf{Domaine} \\
\hline
$2 \to 1$ & $10,2$ & $122$ & Lyman & UV \\
\hline
$3 \to 2$ & $1,89$ & $656$ & Balmer ($\mathrm{H}_\alpha$) & visible (rouge) \\
\hline
$4 \to 2$ & $2,55$ & $486$ & Balmer ($\mathrm{H}_\beta$) & visible (bleu-vert) \\
\hline
$5 \to 2$ & $2,86$ & $434$ & Balmer ($\mathrm{H}_\gamma$) & visible (violet) \\
\hline
$4 \to 3$ & $0,661$ & $1875$ & Paschen & IR \\
\hline
\end{tabularx}
\end{center}

\textbf{Étape 6 : identification du spectre mystère.} Les trois raies observées coïncident exactement avec nos calculs :
\[
656\ \mathrm{nm} \leftrightarrow 3 \to 2 \ (\mathrm{H}_\alpha), \quad 486\ \mathrm{nm} \leftrightarrow 4 \to 2 \ (\mathrm{H}_\beta), \quad 434\ \mathrm{nm} \leftrightarrow 5 \to 2 \ (\mathrm{H}_\gamma).
\]
Ce sont les trois raies visibles les plus intenses de la série de Balmer de l'hydrogène. Aucun autre élément ne présente des raies à ces positions exactes : le spectre mystère est bien celui de l'\textbf{hydrogène}. La nébuleuse observée contient donc de l'hydrogène excité, dont les atomes se désexcitent en émettant ces photons caractéristiques.

\textbf{Étape 7 : question d'ouverture.} Si toutes les raies de Balmer sont décalées vers les grandes longueurs d'onde (vers le rouge), c'est le \cle{décalage vers le rouge} ou \emph{redshift}, manifestation de l'\textbf{effet Doppler-Fizeau} : l'étoile \textbf{s'éloigne de l'observateur}. Plus le décalage relatif $\dfrac{\Delta \lambda}{\lambda}$ est grand, plus la vitesse d'éloignement est grande. C'est précisément ce type d'observation (Hubble, 1929) qui a établi l'expansion de l'Univers : les galaxies lointaines s'éloignent d'autant plus vite qu'elles sont distantes.
\end{exoresolu}

\begin{attention}[Pièges fréquents rencontrés pendant l'activité]
\begin{itemize}
\item Ne pas oublier le signe moins dans $E_n$ : les énergies des états liés sont négatives, et c'est la \emph{différence} $|\Delta E|$ qui donne l'énergie du photon, toujours positive.
\item Convertir systématiquement les eV en joules avant d'appliquer $\lambda = \dfrac{hc}{\Delta E}$ en unités SI, ou utiliser la formule pratique $\lambda\ (\mathrm{nm}) = 1240/\Delta E\ (\mathrm{eV})$ sans mélanger les deux.
\item Une transition $n_i \to n_f$ avec $n_i > n_f$ correspond à une \emph{émission} ; l'inverse à une \emph{absorption}. Ne pas inverser.
\item Le photon absorbé doit avoir une énergie \emph{exactement égale} à l'écart entre deux niveaux : un photon de $10,0\ \mathrm{eV}$ ne peut pas provoquer la transition $1 \to 2$ ($10,2\ \mathrm{eV}$).
\item Attention aux chiffres significatifs : $E_0 = 13,6\ \mathrm{eV}$ (3 chiffres) impose 3 chiffres significatifs sur les résultats ($E_3 = -1,51\ \mathrm{eV}$, pas $-1,5$).
\end{itemize}
\end{attention}

\courssub{Bilan de l'activité}

Cette activité a permis de mettre en évidence plusieurs idées fondamentales de la leçon :

\begin{itemize}
\item \textbf{La quantification est une nécessité expérimentale.} Le fait que le spectre de l'hydrogène soit composé de raies discrètes, et non d'un continuum, s'explique parfaitement par l'existence de niveaux d'énergie quantifiés en $-E_0/n^2$ : seules certaines énergies de photons sont possibles.
\item \textbf{Le spectre est une carte d'identité.} Les positions des raies calculées ($656$, $486$, $434\ \mathrm{nm}$) coïncident avec les raies réellement observées : l'analyse spectrale permet d'identifier un élément chimique à distance, sans jamais le toucher. C'est ainsi qu'on connaît la composition des étoiles, des nébuleuses et de l'atmosphère des planètes.
\item \textbf{L'organisation en séries.} Toutes les raies d'une même série partagent le même niveau d'arrivée : Lyman ($n = 1$, UV), Balmer ($n = 2$, visible), Paschen ($n = 3$, IR). Le diagramme énergétique rend cette classification visuelle et évidente.
\item \textbf{Une porte ouverte vers l'astrophysique moderne.} Le décalage des raies (effet Doppler-Fizeau) transforme le spectre en \emph{radar} : on mesure la vitesse des astres, on découvre l'expansion de l'Univers. La même physique quantique qui décrit un atome minuscule permet ainsi de comprendre le cosmos tout entier — une belle illustration de la puissance unificatrice de la physique, du laboratoire de Ngaoundéré jusqu'aux galaxies lointaines.
\end{itemize}

\begin{savaistu}[L'hélium, l'élément découvert dans le Soleil]
En 1868, l'astronome Jules Janssen observe une raie jaune inconnue ($587,6\ \mathrm{nm}$) dans le spectre du Soleil, qui ne correspond à aucun élément terrestre. On la baptise \emph{hélium}, du grec \emph{hélios} (Soleil). Ce n'est qu'en 1895 que l'hélium est identifié sur Terre ! C'est la seule fois de l'histoire où un élément chimique a été découvert dans un astre avant de l'être dans un laboratoire — une preuve éclatante que la spectroscopie est un véritable outil d'exploration de l'Univers.
\end{savaistu}

\courssub{Méthodes de résolution et applications type Bac}

\begin{methode}[Calculer la longueur d'onde d'une transition dans l'hydrogène]
\begin{enumerate}
\item Identifier les niveaux initial $n_i$ et final $n_f$.
\item Calculer les énergies $E_{n_i}$ et $E_{n_f}$ avec $E_n = -13,6/n^2$ (en eV).
\item Calculer $\Delta E = |E_{n_f} - E_{n_i}|$ (en eV).
\item Appliquer $\lambda\ (\mathrm{nm}) = 1240 / \Delta E\ (\mathrm{eV})$.
\item Vérifier le domaine spectral (UV, visible, IR) et la série correspondante.
\end{enumerate}
\end{methode}

\begin{exoresolu}[Exemple type Bac : Absorption dans une nébuleuse]
\textbf{Énoncé.} Une nébuleuse froide d'hydrogène est éclairée par une source continue. On observe une raie d'absorption à $\lambda = 102,6\ \mathrm{nm}$. Déterminer la transition correspondante.
\textbf{Solution.}
\begin{enumerate}
\item $\Delta E = 1240 / 102,6 = 12,09\ \mathrm{eV}$.
\item On cherche $n_i, n_f$ tels que $|E_{n_f} - E_{n_i}| = 12,09\ \mathrm{eV}$.
\item Comme $\lambda < 400\ \mathrm{nm}$, c'est l'UV $\Rightarrow$ série de Lyman ($n_f = 1$).
\item $E_1 = -13,6\ \mathrm{eV}$. $\Delta E = E_{n_i} - E_1 = 13,6 - 13,6/n_i^2 = 12,09$.
\item $13,6/n_i^2 = 1,51 \Rightarrow n_i^2 = 9 \Rightarrow n_i = 3$.
\item Transition : \textbf{$3 \to 1$ (Lyman $\beta$)}. \checkmark
\end{enumerate}
\end{exoresolu}

\begin{exercice}[Entraînement : Série de Brackett]
La série de Brackett correspond aux transitions arrivant sur le niveau $n_f = 4$.
\begin{enumerate}
\item Calculer la longueur d'onde de la première raie de cette série (transition $5 \to 4$).
\item Dans quel domaine spectral se situe-t-elle ?
\item Pourquoi cette série est-elle difficile à observer depuis le sol ?
\end{enumerate}
\end{exercice}

\begin{corrige}[Exercice : Série de Brackett]
\begin{enumerate}
\item $E_4 = -0,850\ \mathrm{eV}$, $E_5 = -0,544\ \mathrm{eV}$. $\Delta E = 0,850 - 0,544 = 0,306\ \mathrm{eV}$.
$\lambda = 1240 / 0,306 = 4052\ \mathrm{nm} = 4,05\ \mu\mathrm{m}$.
\item Domaine \textbf{infrarouge moyen} (IR).
\item L'atmosphère terrestre absorbe fortement dans l'infrarouge moyen (vapeur d'eau, $\mathrm{CO}_2$). L'observation nécessite des télescopes spatiaux (comme le JWST) ou des sites d'altitude très sèche (Atacama).
\end{enumerate}
\end{corrige}

\newpage

\coursec{Méthodes et exercices résolus}

\coursec{Niveaux d'énergie de l'atome d'hydrogène}

\courssub{Fiche méthode 1 : Calculer l'énergie d'un photon à partir de sa fréquence ou de sa longueur d'onde}

\begin{methode}[Photon : de $\lambda$ ou $\nu$ à l'énergie]
\begin{enumerate}
\item \textbf{Identifier la donnée} : fréquence $\nu$ (en Hz) ou longueur d'onde $\lambda$ (en m). Si $\lambda$ est donnée en nm, convertir : $1~\text{nm} = 10^{-9}~\text{m}$.
\item \textbf{Choisir la relation} : $E = h\nu$ si on connaît $\nu$ ; $E = \dfrac{hc}{\lambda}$ si on connaît $\lambda$.
\item \textbf{Appliquer en unités SI} : $h = 6{,}63\times 10^{-34}$~J$\cdot$s, $c = 3{,}00\times 10^{8}$~m/s. Le résultat est en joules.
\item \textbf{Convertir en eV si demandé} : $E(\text{eV}) = \dfrac{E(\text{J})}{1{,}60\times 10^{-19}}$.
\item \textbf{Réflexe pratique} : la formule approchée $E(\text{eV}) \approx \dfrac{1240}{\lambda(\text{nm})}$ permet de vérifier rapidement l'ordre de grandeur.
\end{enumerate}
\end{methode}

\begin{exoresolu}[Énergie d'un photon de lumière verte]
Au Cameroun, certains feux de signalisation à LED émettent une lumière verte de longueur d'onde $\lambda = \SI{532}{nm}$.
\begin{enumerate}
\item Calculer la fréquence de cette radiation.
\item Calculer l'énergie d'un photon en joules, puis en électronvolts.
\end{enumerate}
Données : $h = 6{,}63\times 10^{-34}$~J$\cdot$s ; $c = 3{,}00\times 10^{8}$~m/s ; $1~\text{eV} = 1{,}60\times 10^{-19}$~J.
\tcblower
\textbf{1. Fréquence de la radiation.}

La relation entre fréquence et longueur d'onde dans le vide est $\nu = \dfrac{c}{\lambda}$.

Conversion : $\lambda = \SI{532}{nm} = 532\times 10^{-9}~\text{m} = 5{,}32\times 10^{-7}$~m.

\[
\nu = \frac{3{,}00\times 10^{8}}{5{,}32\times 10^{-7}} = 5{,}64\times 10^{14}~\text{Hz}.
\]

\textbf{2. Énergie du photon.}

En joules :
\[
E = h\nu = 6{,}63\times 10^{-34}\times 5{,}64\times 10^{14} = 3{,}74\times 10^{-19}~\text{J}.
\]

Vérification par la formule directe : $E = \dfrac{hc}{\lambda} = \dfrac{6{,}63\times 10^{-34}\times 3{,}00\times 10^{8}}{5{,}32\times 10^{-7}} = 3{,}74\times 10^{-19}$~J. Cohérent.

En électronvolts :
\[
E = \frac{3{,}74\times 10^{-19}}{1{,}60\times 10^{-19}} = 2{,}34~\text{eV}.
\]

Vérification rapide : $E \approx \dfrac{1240}{532} \approx 2{,}33$~eV. L'ordre de grandeur est correct.

\textbf{Conclusion :} chaque photon vert émis par la LED transporte une énergie $E \approx \SI{3,74e-19}{J}$, soit $\SI{2,34}{eV}$. C'est l'ordre de grandeur typique des transitions électroniques dans les atomes : la lumière visible correspond bien à des sauts entre niveaux d'énergie atomiques.
\end{exoresolu}

\courssub{Fiche méthode 2 : Appliquer la relation $E_n = -\dfrac{E_0}{n^2}$}

\begin{methode}[Calculer l'énergie d'un niveau de l'atome d'hydrogène]
\begin{enumerate}
\item \textbf{Identifier le niveau} : $n = 1$ (fondamental), $n = 2, 3, \ldots$ (états excités), $n \to \infty$ (atome ionisé, $E_\infty = 0$).
\item \textbf{Appliquer la formule} : $E_n = -\dfrac{13{,}6}{n^2}$ (en eV). Le signe moins est obligatoire : l'électron est lié au noyau.
\item \textbf{Convertir en joules si demandé} : multiplier par $1{,}60\times 10^{-19}$.
\item \textbf{Interpréter} : plus $n$ est grand, plus $E_n$ est proche de zéro (niveaux qui se resserrent) ; l'énergie d'ionisation depuis le niveau $n$ vaut $E_{\text{ion}} = -E_n = \dfrac{13{,}6}{n^2}$ eV.
\end{enumerate}
\end{methode}

\begin{exoresolu}[Énergies des premiers niveaux et ionisation]
\begin{enumerate}
\item Calculer, en eV puis en J, l'énergie de l'atome d'hydrogène dans son état fondamental ($n=1$) et dans son premier état excité ($n=2$).
\item Quelle énergie minimale faut-il fournir pour ioniser l'atome pris dans son état fondamental ? Dans l'état $n=2$ ?
\item Comparer $E_2 - E_1$ et $E_3 - E_2$. Conclure sur la disposition des niveaux.
\end{enumerate}
\tcblower
\textbf{1. Énergies des niveaux.}

État fondamental :
\[
E_1 = -\frac{13{,}6}{1^2} = -13{,}6~\text{eV} = -13{,}6\times 1{,}60\times 10^{-19} = -2{,}18\times 10^{-18}~\text{J}.
\]

Premier état excité :
\[
E_2 = -\frac{13{,}6}{2^2} = -\frac{13{,}6}{4} = -3{,}40~\text{eV} = -5{,}44\times 10^{-19}~\text{J}.
\]

\textbf{2. Énergies d'ionisation.}

Ioniser, c'est amener l'électron du niveau $n$ au niveau $E_\infty = 0$. L'énergie minimale à fournir est donc $\Delta E = E_\infty - E_n = -E_n$.

Depuis le fondamental : $E_{\text{ion}} = 0 - (-13{,}6) = 13{,}6$~eV.

Depuis $n=2$ : $E_{\text{ion}} = 0 - (-3{,}40) = 3{,}40$~eV.

\textbf{3. Comparaison des écarts.}

$E_3 = -\dfrac{13{,}6}{9} = -1{,}51$~eV.

\[
E_2 - E_1 = -3{,}40 - (-13{,}6) = 10{,}2~\text{eV} \quad ; \quad E_3 - E_2 = -1{,}51 - (-3{,}40) = 1{,}89~\text{eV}.
\]

\textbf{Conclusion :} les écarts entre niveaux successifs diminuent fortement quand $n$ augmente : les niveaux d'énergie se resserrent vers la limite $E_\infty = 0$. C'est pourquoi, sur un diagramme énergétique, les traits horizontaux se rapprochent de plus en plus en montant.
\end{exoresolu}

\courssub{Fiche méthode 3 : Déterminer la longueur d'onde d'une transition}

\begin{methode}[Transition $n_i \to n_f$ : énergie et longueur d'onde du photon]
\begin{enumerate}
\item \textbf{Repérer le sens de la transition} : si $n_i > n_f$, l'atome se désexcite et \emph{émet} un photon ; si $n_i < n_f$, l'atome \emph{absorbe} un photon.
\item \textbf{Calculer l'énergie échangée} (toujours positive pour le photon) :
\[
E_{\text{photon}} = |E_f - E_i| = 13{,}6\left|\frac{1}{n_f^2} - \frac{1}{n_i^2}\right|~\text{(eV)}.
\]
\item \textbf{En déduire la longueur d'onde} : $\lambda = \dfrac{hc}{E_{\text{photon}}}$ avec $E$ converti en joules, ou directement $\lambda(\text{nm}) \approx \dfrac{1240}{E(\text{eV})}$.
\item \textbf{Conclure sur le domaine spectral} : $\lambda < 400$~nm : ultraviolet ; $400 < \lambda < 800$~nm : visible ; $\lambda > 800$~nm : infrarouge.
\end{enumerate}
\end{methode}

\begin{exoresolu}[Transition de $n=3$ vers $n=2$ : la raie rouge de l'hydrogène]
L'électron d'un atome d'hydrogène passe du niveau $n=3$ au niveau $n=2$.
\begin{enumerate}
\item Cette transition s'accompagne-t-elle d'une émission ou d'une absorption ? Justifier.
\item Calculer l'énergie du photon mis en jeu, en eV puis en J.
\item Calculer la longueur d'onde de la radiation correspondante. À quel domaine appartient-elle ?
\end{enumerate}
Données : $h = 6{,}63\times 10^{-34}$~J$\cdot$s ; $c = 3{,}00\times 10^{8}$~m/s ; $1~\text{eV} = 1{,}60\times 10^{-19}$~J.
\tcblower
\textbf{1. Sens de la transition.}

L'électron passe d'un niveau de nombre quantique élevé ($n_i = 3$) vers un niveau plus bas ($n_f = 2$). L'énergie de l'atome diminue : $E_2 < E_3$. L'énergie perdue est emportée par un photon : c'est une \textbf{émission}.

\textbf{2. Énergie du photon.}

\[
E_3 = -\frac{13{,}6}{9} = -1{,}51~\text{eV} \quad ; \quad E_2 = -\frac{13{,}6}{4} = -3{,}40~\text{eV}.
\]

\[
E_{\text{photon}} = E_3 - E_2 = -1{,}51 - (-3{,}40) = 1{,}89~\text{eV}.
\]

En joules : $E_{\text{photon}} = 1{,}89\times 1{,}60\times 10^{-19} = 3{,}02\times 10^{-19}$~J.

\textbf{3. Longueur d'onde.}

\[
\lambda = \frac{hc}{E_{\text{photon}}} = \frac{6{,}63\times 10^{-34}\times 3{,}00\times 10^{8}}{3{,}02\times 10^{-19}} = 6{,}59\times 10^{-7}~\text{m} = 659~\text{nm}.
\]

Vérification : $\lambda \approx \dfrac{1240}{1{,}89} \approx 656$~nm (écart dû aux arrondis). 

\textbf{Conclusion :} la transition $3 \to 2$ produit un photon de longueur d'onde $\lambda \approx \SI{656}{nm}$, située dans le \textbf{visible} (rouge). C'est la célèbre raie rouge H$\alpha$ de la série de Balmer, la plus intense du spectre de l'hydrogène, observable avec une simple lampe à décharge et un prisme.
\end{exoresolu}

\courssub{Fiche méthode 4 : Identifier une transition à partir de la longueur d'onde}

\begin{methode}[Retrouver les niveaux $n_i$ et $n_f$ connaissant $\lambda$]
\begin{enumerate}
\item \textbf{Calculer l'énergie du photon} : $E = \dfrac{hc}{\lambda}$ (en J), convertie en eV.
\item \textbf{Écrire la relation de quantification} : il faut trouver deux entiers $n_i > n_f \geq 1$ tels que
\[
E = 13{,}6\left(\frac{1}{n_f^2} - \frac{1}{n_i^2}\right)~\text{eV}.
\]
\item \textbf{Procéder par essais raisonnés} : tester $n_f = 1$ (Lyman), puis $n_f = 2$ (Balmer), etc. Pour chaque $n_f$, calculer $\dfrac{1}{n_i^2} = \dfrac{1}{n_f^2} - \dfrac{E}{13{,}6}$ et vérifier si le résultat est l'inverse du carré d'un entier.
\item \textbf{Vérifier} en recalculant $\lambda$ à partir du couple $(n_i, n_f)$ trouvé.
\end{enumerate}
\end{methode}

\begin{exoresolu}[Identification d'une raie ultraviolette]
Le spectre d'émission de l'hydrogène présente une raie de longueur d'onde $\lambda = \SI{121,6}{nm}$.
\begin{enumerate}
\item Calculer l'énergie du photon correspondant, en eV.
\item Montrer que cette raie correspond à la transition $n=2 \to n=1$.
\item À quelle série appartient cette raie ? Dans quel domaine spectral se situe-t-elle ?
\end{enumerate}
\tcblower
\textbf{1. Énergie du photon.}

\[
E = \frac{hc}{\lambda} = \frac{6{,}63\times 10^{-34}\times 3{,}00\times 10^{8}}{121{,}6\times 10^{-9}} = 1{,}636\times 10^{-18}~\text{J}.
\]

Conversion : $E = \dfrac{1{,}636\times 10^{-18}}{1{,}60\times 10^{-19}} = 10{,}2$~eV.

\textbf{2. Identification de la transition.}

Cherchons d'abord dans la série de Lyman ($n_f = 1$) :
\[
\frac{1}{n_i^2} = \frac{1}{1^2} - \frac{E}{13{,}6} = 1 - \frac{10{,}2}{13{,}6} = 1 - 0{,}75 = 0{,}25 = \frac{1}{4}.
\]

Donc $n_i^2 = 4$, soit $n_i = 2$ : c'est bien un entier. La transition est $n = 2 \to n = 1$.

Vérification : $E_2 - E_1 = -3{,}40 - (-13{,}6) = 10{,}2$~eV, et $\lambda = \dfrac{1240}{10{,}2} \approx 122$~nm. Cohérent avec l'énoncé.

\textbf{3. Série et domaine.}

La transition aboutit au niveau fondamental $n=1$ : elle appartient à la \textbf{série de Lyman}. Comme $\lambda = \SI{121,6}{nm} < 400$~nm, la raie se situe dans l'\textbf{ultraviolet}.

\textbf{Conclusion :} la raie à $\SI{121,6}{nm}$ est la raie Lyman-$\alpha$ ($2 \to 1$), la raie la plus énergétique de l'hydrogène ; elle est très utilisée en astrophysique pour détecter l'hydrogène dans les galaxies lointaines.
\end{exoresolu}

\courssub{Fiche méthode 5 : Exploiter un diagramme énergétique (émission, absorption, ionisation)}

\begin{methode}[Lire et exploiter un diagramme de niveaux]
\begin{enumerate}
\item \textbf{Orienter le diagramme} : l'énergie croît vers le haut ; le niveau $n=1$ est en bas ($-13{,}6$~eV) ; la limite $E=0$ correspond à l'ionisation.
\item \textbf{Émission} : flèche descendante de $n_i$ vers $n_f < n_i$ ; le photon émis a pour énergie $h\nu = E_i - E_f$.
\item \textbf{Absorption} : flèche montante ; le photon absorbé doit avoir \emph{exactement} l'énergie $E_f - E_i$ : l'absorption est sélective (quantifiée).
\item \textbf{Ionisation} : flèche partant du niveau $n$ et dépassant $E = 0$ ; si le photon incident a une énergie $E > |E_n|$, l'excédent devient l'énergie cinétique de l'électron éjecté : $E_c = E - |E_n|$.
\item \textbf{Compter les raies} : un atome excité au niveau $n$ peut émettre au maximum $\dfrac{n(n-1)}{2}$ raies différentes en cascade vers le fondamental.
\end{enumerate}
\end{methode}

\begin{exoresolu}[Absorption sélective et photo-ionisation]
Un atome d'hydrogène est dans son état fondamental ($E_1 = -13{,}6$~eV). On l'éclaire avec des photons d'énergie : (a) $\SI{10,0}{eV}$ ; (b) $\SI{10,2}{eV}$ ; (c) $\SI{15,0}{eV}$.
\begin{enumerate}
\item Pour chaque photon, dire s'il peut être absorbé et préciser l'état final de l'atome.
\item Dans le cas (c), calculer l'énergie cinétique de l'électron éjecté.
\item L'atome est maintenant excité au niveau $n=3$. Combien de raies différentes peut-il émettre en revenant au fondamental ? Les citer.
\end{enumerate}
On rappelle : $E_2 = -3{,}40$~eV, $E_3 = -1{,}51$~eV, $E_4 = -0{,}85$~eV.
\tcblower
\textbf{1. Absorption ou non.}

Un photon ne peut être absorbé que si son énergie correspond \emph{exactement} à l'écart entre deux niveaux (ou est supérieure à l'énergie d'ionisation).

Écarts possibles depuis le fondamental :
\[
E_2 - E_1 = 10{,}2~\text{eV} \ ; \quad E_3 - E_1 = 12{,}09~\text{eV} \ ; \quad E_4 - E_1 = 12{,}75~\text{eV}.
\]

(a) Photon de $\SI{10,0}{eV}$ : ne correspond à aucun écart et est inférieure à $\SI{13,6}{eV}$ (ionisation impossible). \textbf{Le photon n'est pas absorbé} : il traverse sans interaction.

(b) Photon de $\SI{10,2}{eV}$ : égal à $E_2 - E_1$. \textbf{Le photon est absorbé} et l'atome passe au niveau $n=2$.

(c) Photon de $\SI{15,0}{eV}$ : supérieure à $\SI{13,6}{eV}$. \textbf{Le photon est absorbé et l'atome est ionisé} : l'électron est éjecté.

\textbf{2. Énergie cinétique de l'électron éjecté.}

Conservation de l'énergie : $E_{\text{photon}} = |E_1| + E_c$, d'où
\[
E_c = 15{,}0 - 13{,}6 = 1{,}4~\text{eV} = 2{,}24\times 10^{-19}~\text{J}.
\]

\textbf{3. Raies émises depuis $n=3$.}

Transitions possibles vers des niveaux inférieurs : $3 \to 2$, $3 \to 1$, puis $2 \to 1$ (en cascade après $3 \to 2$). Soit $\dfrac{3\times 2}{2} = 3$ raies :
\begin{itemize}
\item $3 \to 1$ : $E = 12{,}09$~eV, $\lambda \approx \SI{103}{nm}$ (UV, Lyman) ;
\item $3 \to 2$ : $E = 1{,}89$~eV, $\lambda \approx \SI{656}{nm}$ (visible, Balmer) ;
\item $2 \to 1$ : $E = 10{,}2$~eV, $\lambda \approx \SI{122}{nm}$ (UV, Lyman).
\end{itemize}

\textbf{Conclusion :} l'absorption est un phénomène quantifié et sélectif : seuls les photons d'énergie exactement égale à un écart de niveaux sont absorbés ; au-delà de $\SI{13,6}{eV}$, tout photon ionise l'atome et l'excédent devient énergie cinétique de l'électron.
\end{exoresolu}

\courssub{Fiche méthode 6 : Séries spectrales et raies limites}

\begin{methode}[Séries de Lyman, Balmer, Paschen : raies extrêmes]
\begin{enumerate}
\item \textbf{Identifier la série} par son niveau d'arrivée : Lyman $n_f = 1$ (UV), Balmer $n_f = 2$ (visible/UV proche), Paschen $n_f = 3$ (IR).
\item \textbf{Raie de plus grande longueur d'onde} (énergie minimale) : transition depuis le niveau immédiatement supérieur, $n_i = n_f + 1$.
\item \textbf{Raie limite} (plus courte longueur d'onde, énergie maximale) : transition depuis $n_i \to \infty$ :
\[
E_{\text{limite}} = \frac{13{,}6}{n_f^2}~\text{eV} \quad ; \quad \lambda_{\text{limite}} = \frac{hc}{E_{\text{limite}}}.
\]
\item \textbf{Encadrer la série} : toutes les raies vérifient $\lambda_{\text{limite}} < \lambda \leq \lambda_{\max}$.
\end{enumerate}
\end{methode}

\begin{exoresolu}[Bornes de la série de Balmer]
La série de Balmer regroupe les transitions aboutissant au niveau $n=2$.
\begin{enumerate}
\item Calculer la longueur d'onde de la raie de plus grande longueur d'onde de cette série.
\item Calculer la longueur d'onde limite de la série.
\item Le spectre de l'hydrogène présente une raie à $\lambda = \SI{434}{nm}$. Peut-elle appartenir à la série de Balmer ? À quelle transition correspond-elle ?
\end{enumerate}
\tcblower
\textbf{1. Raie de plus grande longueur d'onde.}

Énergie minimale : transition $n_i = 3 \to n_f = 2$ :
\[
E = 13{,}6\left(\frac{1}{2^2} - \frac{1}{3^2}\right) = 13{,}6\left(\frac{1}{4} - \frac{1}{9}\right) = 13{,}6\times\frac{5}{36} = 1{,}89~\text{eV}.
\]
\[
\lambda_{\max} = \frac{1240}{1{,}89} \approx 656~\text{nm} \quad \text{(raie H}\alpha\text{, rouge)}.
\]

\textbf{2. Longueur d'onde limite.}

Transition depuis $n_i \to \infty$ :
\[
E_{\text{limite}} = \frac{13{,}6}{2^2} = 3{,}40~\text{eV} \quad ; \quad \lambda_{\text{limite}} = \frac{1240}{3{,}40} \approx 365~\text{nm}.
\]

Toutes les raies de Balmer vérifient donc $365~\text{nm} < \lambda \leq 656~\text{nm}$.

\textbf{3. Raie à 434 nm.}

$365 < 434 < 656$ : la raie \textbf{peut} appartenir à la série de Balmer. Cherchons $n_i$ :
\[
E = \frac{1240}{434} = 2{,}86~\text{eV} \quad ; \quad \frac{1}{n_i^2} = \frac{1}{4} - \frac{2{,}86}{13{,}6} = 0{,}250 - 0{,}210 = 0{,}040 \approx \frac{1}{25}.
\]

Donc $n_i = 5$ : il s'agit de la transition $5 \to 2$ (raie H$\gamma$, bleu-violet).

\textbf{Conclusion :} la série de Balmer s'étend de $\SI{365}{nm}$ (limite, UV proche) à $\SI{656}{nm}$ (H$\alpha$) ; la raie à $\SI{434}{nm}$ correspond à la transition $5 \to 2$.
\end{exoresolu}

\begin{attention}[Les 5 erreurs qui coûtent des points au Bac]
\begin{enumerate}
\item \textbf{Oublier le signe moins dans $E_n = -\dfrac{E_0}{n^2}$.} Les énergies des niveaux sont négatives (électron lié). Écrire $E_2 = +3{,}40$~eV est une erreur de modèle, pas seulement de calcul. En revanche, l'énergie du photon, elle, est toujours positive : $E_{\text{photon}} = |E_f - E_i|$.
\item \textbf{Confondre émission et absorption.} Émission = transition vers un niveau \emph{inférieur} (flèche descendante, l'atome perd de l'énergie) ; absorption = transition vers un niveau \emph{supérieur}, possible uniquement si le photon a exactement la bonne énergie. Inverser les indices dans $\dfrac{1}{n_f^2} - \dfrac{1}{n_i^2}$ change le signe du résultat : toujours vérifier que l'énergie du photon est positive.
\item \textbf{Négliger les conversions d'unités.} Les longueurs d'onde sont souvent données en nm : il faut convertir en mètres avant d'appliquer $E = \dfrac{hc}{\lambda}$. De même, une énergie en eV doit être convertie en joules ($\times 1{,}60\times 10^{-19}$) avant tout calcul avec $h$ et $c$ en SI. Une erreur de facteur $10^{9}$ sur $\lambda$ fausse tout l'exercice.
\item \textbf{Croire qu'un photon de $10{,}0$~eV peut exciter l'atome « partiellement ».} L'absorption est tout ou rien : si l'énergie du photon ne correspond à aucun écart de niveaux et ne permet pas l'ionisation, le photon n'est pas absorbé. Ne jamais écrire que l'atome « garde une partie » de l'énergie.
\item \textbf{Mal rédiger la conclusion.} Le correcteur attend une phrase complète avec valeur numérique, unité et interprétation physique : « $\lambda = 656$~nm, cette radiation appartient au domaine visible (rouge), c'est la raie H$\alpha$ de la série de Balmer ». Un résultat brut sans unité ni domaine spectral perd des points, même si le calcul est juste.
\end{enumerate}
\end{attention}

\newpage

\coursec{Exercices}

\courssub{Je vérifie mes connaissances}

\begin{exercice}[QCM : le photon]
Pour chaque affirmation, choisir la (ou les) bonne(s) réponse(s). On donne $h = \num{6,63e-34}$ J·s et $c = \num{3,00e8}$ m/s.
\begin{enumerate}
\item L'énergie d'un photon de fréquence $\nu$ vaut :
\begin{enumerate}
\item $E = \dfrac{h}{\nu}$ \quad b) $E = h\nu$ \quad c) $E = \dfrac{\nu}{h}$
\end{enumerate}
\item Quand la longueur d'onde d'une radiation augmente, l'énergie du photon associé :
\begin{enumerate}
\item augmente \quad b) diminue \quad c) reste constante
\end{enumerate}
\item Un photon de longueur d'onde $\lambda = \num{500}$ nm a une énergie voisine de :
\begin{enumerate}
\item \num{2,5} eV \quad b) \num{0,25} eV \quad c) \num{25} eV
\end{enumerate}
\item L'énergie d'un photon s'exprime dans le système international en :
\begin{enumerate}
\item watt \quad b) joule \quad c) électronvolt
\end{enumerate}
\end{enumerate}
\end{exercice}

\begin{exercice}[Vrai ou faux, en justifiant]
Répondre par vrai ou faux et justifier chaque réponse par une phrase ou un calcul.
\begin{enumerate}
\item L'énergie de l'atome d'hydrogène peut prendre n'importe quelle valeur.
\item Le niveau fondamental de l'atome d'hydrogène correspond à $n = 0$.
\item L'énergie du niveau $n = 2$ de l'hydrogène vaut $E_2 = -\num{3,40}$ eV.
\item Ioniser un atome d'hydrogène, c'est le faire passer de l'état fondamental à un état excité.
\item Un atome d'hydrogène ne peut absorber un photon que si l'énergie de celui-ci correspond exactement à la différence d'énergie entre deux niveaux.
\end{enumerate}
\end{exercice}

\begin{exercice}[Définitions]
Définir en une ou deux phrases chacun des termes suivants :
\begin{enumerate}
\item photon ;
\item état fondamental et état excité d'un atome ;
\item énergie d'ionisation de l'atome d'hydrogène ;
\item spectre d'émission et spectre d'absorption.
\end{enumerate}
\end{exercice}

\begin{exercice}[Calculs directs]
On donne $E_0 = \num{13,6}$ eV, $h = \num{6,63e-34}$ J·s, $c = \num{3,00e8}$ m/s et $\num{1}$ eV $= \num{1,60e-19}$ J.
\begin{enumerate}
\item Calculer les énergies $E_1$, $E_2$ et $E_4$ des niveaux $n = 1$, $n = 2$ et $n = 4$ de l'atome d'hydrogène.
\item Un atome d'hydrogène passe du niveau $n = 3$ au niveau $n = 2$. Calculer, en eV puis en joule, l'énergie du photon émis.
\item Calculer la fréquence puis la longueur d'onde de ce photon.
\end{enumerate}
\end{exercice}

\courssub{Je m'entraîne}

\begin{exercice}[Énergie d'ionisation depuis un état excité]
On donne $E_0 = \num{13,6}$ eV et $\num{1}$ eV $= \num{1,60e-19}$ J.
\begin{enumerate}
\item Calculer l'énergie minimale nécessaire pour ioniser un atome d'hydrogène initialement dans l'état $n = 3$.
\item Un photon d'énergie \num{2,0} eV frappe cet atome. Peut-il l'ioniser ? Justifier par un calcul.
\item Même question avec un photon d'énergie \num{1,0} eV. Que peut-il se passer alors ?
\end{enumerate}
\end{exercice}

\begin{exercice}[Limite de la série de Lyman]
On donne $h = \num{6,63e-34}$ J·s, $c = \num{3,00e8}$ m/s, $E_0 = \num{13,6}$ eV et $\num{1}$ eV $= \num{1,60e-19}$ J.
\begin{enumerate}
\item Rappeler à quelles transitions correspond la série de Lyman.
\item Calculer la longueur d'onde de la transition $\infty \rightarrow 1$ (limite de la série).
\item Calculer la longueur d'onde de la transition $2 \rightarrow 1$, la plus grande de la série.
\item Dans quel domaine spectral se situent toutes les raies de Lyman ? Expliquer pourquoi aucune raie de cette série n'est visible à l'œil nu.
\end{enumerate}
\end{exercice}

\begin{exercice}[Absorption d'un photon]
Un atome d'hydrogène, initialement dans son état fondamental, absorbe un photon de longueur d'onde $\lambda = \num{97,3}$ nm. On donne $h = \num{6,63e-34}$ J·s, $c = \num{3,00e8}$ m/s, $E_0 = \num{13,6}$ eV et $\num{1}$ eV $= \num{1,60e-19}$ J.
\begin{enumerate}
\item Calculer l'énergie du photon absorbé, en joule puis en eV.
\item Déterminer le niveau d'énergie $n$ atteint par l'atome après absorption.
\item À quelle série de raies appartient cette transition ? Dans quel domaine spectral se situe-t-elle ?
\end{enumerate}
\end{exercice}

\begin{exercice}[Identification d'un spectre]
Le spectre d'émission d'une lampe à décharge présente quatre raies dans le visible : $\lambda_1 = \num{410}$ nm, $\lambda_2 = \num{434}$ nm, $\lambda_3 = \num{486}$ nm et $\lambda_4 = \num{656}$ nm. On donne $h = \num{6,63e-34}$ J·s, $c = \num{3,00e8}$ m/s, $E_0 = \num{13,6}$ eV et $\num{1}$ eV $= \num{1,60e-19}$ J.
\begin{enumerate}
\item Calculer l'énergie, en eV, du photon associé à chacune des quatre raies.
\item Calculer les différences d'énergie $E_n - E_2$ pour $n = 3$, $n = 4$, $n = 5$ et $n = 6$ dans l'atome d'hydrogène.
\item Montrer que le spectre observé est bien celui de l'hydrogène et identifier la transition correspondant à chaque raie.
\item À quelle série ces raies appartiennent-elles ? Justifier qu'elles sont visibles.
\end{enumerate}
\end{exercice}

\courssub{Je me prépare au Bac}

\begin{exercice}[Type Bac : diagramme énergétique et raie H$\alpha$]
\textit{Données :} $E_0 = \num{13,6}$ eV ; $h = \num{6,63e-34}$ J·s ; $c = \num{3,00e8}$ m/s ; $\num{1}$ eV $= \num{1,60e-19}$ J.

Dans un laboratoire de l'Université de Yaoundé I, des étudiants étudient le spectre d'émission d'une lampe à hydrogène.
\begin{enumerate}
\item Rappeler l'expression de l'énergie $E_n$ du niveau $n$ de l'atome d'hydrogène. Que représente le nombre $n$ ?
\item Calculer les énergies $E_1$, $E_2$, $E_3$ et $E_4$ des quatre premiers niveaux, en eV.
\item Représenter le diagramme énergétique de l'atome d'hydrogène (échelle : 1 cm pour 2 eV) en y plaçant ces quatre niveaux ainsi que le niveau d'ionisation.
\item Les étudiants observent une raie rouge intense appelée raie H$\alpha$, correspondant à la transition $3 \rightarrow 2$.
\begin{enumerate}
\item Calculer l'énergie du photon émis lors de cette transition, en eV puis en joule.
\item En déduire la longueur d'onde de la raie H$\alpha$.
\item Citer la couleur correspondante et vérifier la cohérence avec l'observation.
\end{enumerate}
\item Calculer l'énergie minimale d'ionisation de l'atome d'hydrogène pris dans son état fondamental. La convertir en joule.
\end{enumerate}
\end{exercice}

\begin{exercice}[Type Bac : désexcitation depuis le niveau $n = 4$]
\textit{Données :} $E_0 = \num{13,6}$ eV ; $h = \num{6,63e-34}$ J·s ; $c = \num{3,00e8}$ m/s ; $\num{1}$ eV $= \num{1,60e-19}$ J ; domaine du visible : de \num{400} nm à \num{800} nm.

Lors d'une décharge électrique dans une enceinte contenant de l'hydrogène, des atomes sont portés au niveau d'énergie $n = 4$.
\begin{enumerate}
\item Calculer l'énergie $E_4$ du niveau $n = 4$.
\item Énumérer, à l'aide d'un diagramme énergétique, toutes les transitions possibles lors de la désexcitation de ces atomes vers l'état fondamental (transitions directes et en cascade). Combien de raies d'émission distinctes peut-on observer au maximum ?
\item Calculer la plus grande longueur d'onde émise. À quelle transition correspond-elle et dans quel domaine spectral se situe-t-elle ?
\item Calculer la plus petite longueur d'onde émise. À quelle transition correspond-elle et dans quel domaine spectral se situe-t-elle ?
\item Parmi toutes ces raies, lesquelles appartiennent au domaine visible ? Justifier par le calcul.
\item À quelles séries (Lyman, Balmer, Paschen) appartiennent les différentes raies obtenues ?
\end{enumerate}
\end{exercice}

\begin{exercice}[Type Bac : photo-ionisation de l'hydrogène]
\textit{Données :} $E_0 = \num{13,6}$ eV ; $h = \num{6,63e-34}$ J·s ; $c = \num{3,00e8}$ m/s ; $\num{1}$ eV $= \num{1,60e-19}$ J ; masse de l'électron $m_e = \num{9,11e-31}$ kg.

Dans la haute atmosphère, le rayonnement solaire ultraviolet peut ioniser les atomes d'hydrogène présents, phénomène exploité pour expliquer certaines propriétés de l'ionosphère qui affectent les télécommunications.
\begin{enumerate}
\item Un photon d'énergie $E = \num{15,0}$ eV frappe un atome d'hydrogène au repos dans son état fondamental.
\begin{enumerate}
\item Montrer que l'atome peut être ionisé.
\item Expliquer ce que devient l'excédent d'énergie du photon.
\item Calculer, en eV puis en joule, l'énergie cinétique de l'électron éjecté.
\item En déduire la vitesse de l'électron éjecté (on supposera les effets relativistes négligeables).
\end{enumerate}
\item Calculer la longueur d'onde du photon incident de \num{15,0} eV. À quel domaine spectral appartient-il ?
\item Déterminer la longueur d'onde maximale qu'un photon doit posséder pour pouvoir ioniser un atome d'hydrogène pris dans son état fondamental.
\item Un photon d'énergie \num{12,1} eV frappe maintenant un atome d'hydrogène à l'état fondamental. Montrer qu'il ne peut pas ioniser l'atome. Que se passe-t-il alors ? Déterminer le niveau atteint et la série de la raie correspondante.
\end{enumerate}
\end{exercice}

\newpage

\coursec{Corrigés des exercices}

\begin{corrige}[Exercice 1 — QCM : le photon]
\begin{enumerate}
\item \textbf{Réponse b)} : $E = h\nu$ (relation de Planck--Einstein). Vérification par analyse dimensionnelle : $[h] = \text{J}\cdot\text{s}$ et $[\nu] = \text{s}^{-1}$, donc $[h\nu] = \text{J}$ : c'est bien une énergie. Les réponses a) et c) donneraient des J$\cdot$s$^2$ et des J$^{-1}\cdot$s$^{-2}$, sans signification énergétique.
\item \textbf{Réponse b)} : diminue. En effet $E = \dfrac{hc}{\lambda}$ : l'énergie est \emph{inversement proportionnelle} à la longueur d'onde. Si $\lambda$ augmente, $E$ diminue. C'est pourquoi les UV ($\lambda$ petit) sont plus énergétiques — et plus dangereux — que l'infrarouge ($\lambda$ grand).
\item \textbf{Réponse a)} : $E(\text{eV}) = \dfrac{1240}{\lambda(\text{nm})} = \dfrac{1240}{500} = \num{2,48} \approx \SI{2,5}{eV}$. Vérification en SI : $E = \dfrac{6,63\times 10^{-34}\times 3,00\times 10^{8}}{500\times 10^{-9}} \approx \SI{3,98e-19}{J} \approx \SI{2,49}{eV}$.
\item \textbf{Réponse b)} : le joule. Le watt est l'unité SI de \emph{puissance} ($\SI{1}{W} = \SI{1}{J/s}$) ; l'électronvolt est une unité commode en physique atomique mais \emph{hors} Système International.
\end{enumerate}
\textbf{Point de vigilance :} au Bac, une réponse de QCM non justifiée peut coûter des points si la consigne l'exige ; l'analyse dimensionnelle est une justification rapide et toujours appréciée.
\end{corrige}

\begin{corrige}[Exercice 2 — Vrai ou faux, en justifiant]
\begin{enumerate}
\item \textbf{Faux.} L'énergie de l'atome d'hydrogène est \cle{quantifiée} : elle ne peut prendre que les valeurs discrètes $E_n = -\dfrac{E_0}{n^2}$ avec $n \in \mathbb{N}^*$. Toute valeur intermédiaire (par exemple $-\SI{5,0}{eV}$) est interdite pour un état lié.
\item \textbf{Faux.} L'état fondamental correspond à $n = 1$, d'énergie $E_1 = -\SI{13,6}{eV}$ (énergie minimale, état le plus stable). La valeur $n = 0$ est exclue : $n \in \mathbb{N}^*$, et $E_0/0^2$ n'a d'ailleurs aucun sens mathématique.
\item \textbf{Vrai.} $E_2 = -\dfrac{13,6}{2^2} = -\dfrac{13,6}{4} = \SI{-3,40}{eV}$.
\item \textbf{Faux.} Ioniser un atome, c'est \emph{arracher} son électron : transition de l'état fondamental ($n=1$) vers $n \to \infty$ ($E_\infty = \SI{0}{eV}$), l'électron devenant libre. Le passage vers un état excité ($n = 2, 3, \dots$) est une simple \emph{excitation} : l'électron reste lié au noyau.
\item \textbf{Vrai.} C'est la \cle{condition de résonance} : l'atome ne peut absorber un photon que si $h\nu = E_{n_f} - E_{n_i}$ \emph{exactement}, avec $n_f > n_i$. Sinon, le photon traverse le gaz sans interaction : l'atome lui est transparent.
\end{enumerate}
\textbf{Point de vigilance :} « vrai » ou « faux » sans justification ne rapporte aucun point ; chaque réponse doit s'appuyer sur la formule $E_n = -E_0/n^2$ ou sur la condition $\Delta E = h\nu$.
\end{corrige}

\begin{corrige}[Exercice 3 — Définitions]
\begin{enumerate}
\item \textbf{Photon :} grain d'énergie indivisible (quantum) associé à une onde électromagnétique ; il se propage dans le vide à la vitesse $c$, a une masse nulle et transporte l'énergie $E = h\nu = \dfrac{hc}{\lambda}$.
\item \textbf{État fondamental :} état de plus basse énergie de l'atome ($n = 1$ pour l'hydrogène, $E_1 = -\SI{13,6}{eV}$) ; c'est son état le plus stable. \textbf{État excité :} tout état d'énergie supérieure ($n \geq 2$) ; il est instable et l'atome s'en désexcite spontanément en émettant un ou plusieurs photons.
\item \textbf{Énergie d'ionisation :} énergie minimale qu'il faut fournir à l'atome dans son état fondamental pour en extraire l'électron et l'amener au repos à l'infini ($n \to \infty$). Pour l'hydrogène : $E_{\text{ion}} = E_\infty - E_1 = 0 - (-13,6) = \SI{13,6}{eV}$.
\item \textbf{Spectre d'émission :} ensemble de raies lumineuses colorées sur fond noir, émises par un gaz excité ; chaque raie correspond à une désexcitation $n_i \to n_f$ ($n_i > n_f$). \textbf{Spectre d'absorption :} ensemble de raies sombres sur fond continu, obtenu quand une lumière blanche traverse un gaz froid ; chaque raie sombre correspond à l'absorption d'un photon lors d'une transition $n_i \to n_f$ ($n_i < n_f$). Les raies d'émission et d'absorption d'un même gaz se situent aux \emph{mêmes} longueurs d'onde.
\end{enumerate}
\end{corrige}

\begin{corrige}[Exercice 4 — Calculs directs]
\begin{enumerate}
\item Application directe de $E_n = -\dfrac{E_0}{n^2}$ :
\begin{align*}
E_1 &= -\frac{13,6}{1^2} = \SI{-13,6}{eV} \\
E_2 &= -\frac{13,6}{2^2} = -\frac{13,6}{4} = \SI{-3,40}{eV} \\
E_4 &= -\frac{13,6}{4^2} = -\frac{13,6}{16} = \SI{-0,85}{eV}
\end{align*}
\item Transition $3 \to 2$ : c'est une \textbf{émission} ($n_i > n_f$). D'abord $E_3 = -\dfrac{13,6}{9} \approx \SI{-1,51}{eV}$.
\[
\Delta E = E_3 - E_2 = (-1,51) - (-3,40) = \SI{1,89}{eV}
\]
Conversion en joules : $E = 1,89 \times 1,60\times 10^{-19} \approx \SI{3,02e-19}{J}$.
\item Fréquence : $\nu = \dfrac{E}{h} = \dfrac{3,02\times 10^{-19}}{6,63\times 10^{-34}} \approx \SI{4,56e14}{Hz}$.

Longueur d'onde : $\lambda = \dfrac{c}{\nu} = \dfrac{3,00\times 10^{8}}{4,56\times 10^{14}} \approx \SI{6,58e-7}{m} = \SI{658}{nm}$.
(Avec la valeur approchée $hc \approx \SI{1240}{eV\cdot nm}$ : $\lambda = 1240/1,89 \approx \SI{656}{nm}$ ; l'écart vient des arrondis sur $E_3$ et sur les constantes. C'est la raie rouge H$\alpha$ du visible.)
\end{enumerate}
\textbf{Point de vigilance :} penser à convertir les eV en joules \emph{avant} d'utiliser $\nu = E/h$ avec $h$ en J$\cdot$s ; mélanger eV et joules est l'erreur la plus fréquente.
\end{corrige}

\begin{corrige}[Exercice 5 — Énergie d'ionisation depuis un état excité]
\begin{enumerate}
\item L'énergie du niveau $n = 3$ est $E_3 = -\dfrac{13,6}{9} \approx \SI{-1,51}{eV}$. Ioniser, c'est amener l'électron à $E_\infty = \SI{0}{eV}$. L'énergie minimale à fournir est donc :
\[
E_{\text{ion}} = E_\infty - E_3 = 0 - (-1,51) = \SI{1,51}{eV}
\]
Remarque : elle est bien inférieure aux $\SI{13,6}{eV}$ nécessaires depuis le fondamental, car l'électron excité est déjà moins lié.
\item Photon d'énergie $E = \SI{2,0}{eV}$. Comme $2,0 > 1,51$ eV, \textbf{oui, l'atome peut être ionisé}. L'excédent $2,0 - 1,51 = \SI{0,49}{eV}$ est emporté par l'électron éjecté sous forme d'énergie cinétique.
\item Photon d'énergie $E = \SI{1,0}{eV}$. Comme $1,0 < 1,51$ eV, \textbf{l'ionisation est impossible}. Reste à tester une éventuelle excitation depuis $n = 3$ :
\begin{align*}
3\to 4 &: \quad \Delta E = -0,85 - (-1,51) = \SI{0,66}{eV} \\
3\to 5 &: \quad \Delta E = -0,544 - (-1,51) \approx \SI{0,97}{eV} \\
3\to 6 &: \quad \Delta E = -0,378 - (-1,51) \approx \SI{1,13}{eV}
\end{align*}
Aucune différence ne vaut exactement $\SI{1,0}{eV}$ (et les niveaux suivants demandent encore plus). \textbf{Conclusion : le photon n'est pas absorbé} ; il traverse le gaz sans interaction.
\end{enumerate}
\textbf{Point de vigilance :} pour un photon non ionisant, il ne suffit pas de dire « l'énergie est trop faible » : il faut vérifier qu'elle ne correspond à \emph{aucune} transition excitante possible depuis l'état initial.
\end{corrige}

\begin{corrige}[Exercice 6 — Limite de la série de Lyman]
\begin{enumerate}
\item La série de Lyman rassemble toutes les transitions qui se terminent sur le niveau fondamental : $n_i \to n_f = 1$ avec $n_i = 2, 3, 4, \dots, \infty$.
\item Limite de série $\infty \to 1$ : $\Delta E = E_\infty - E_1 = 0 - (-13,6) = \SI{13,6}{eV}$, soit $E = 13,6 \times 1,60\times 10^{-19} = \SI{2,18e-18}{J}$.
\[
\lambda_{\text{lim}} = \frac{hc}{E} = \frac{6,63\times 10^{-34}\times 3,00\times 10^{8}}{2,18\times 10^{-18}} \approx \SI{9,14e-8}{m} \approx \SI{91,2}{nm}
\]
C'est la \emph{plus courte} longueur d'onde de la série (transition la plus énergétique).
\item Transition $2 \to 1$ (plus petite énergie de la série, donc plus grande longueur d'onde) :
\[
\Delta E = E_2 - E_1 = (-3,40) - (-13,6) = \SI{10,2}{eV}
\]
\[
\lambda = \frac{hc}{\Delta E} = \frac{1240}{10,2} \approx \SI{121,6}{nm}
\]
\item Toutes les raies de Lyman ont des longueurs d'onde comprises entre $\SI{91,2}{nm}$ et $\SI{121,6}{nm}$, donc toutes inférieures à $\SI{400}{nm}$ : elles appartiennent à l'\textbf{ultraviolet}. Or l'œil humain n'est sensible qu'au domaine visible ($\approx 400$ à $\SI{800}{nm}$) : aucune raie de Lyman n'est visible à l'œil nu ; il faut un détecteur UV.
\end{enumerate}
\textbf{Point de vigilance :} « limite de série » signifie $n_i \to \infty$, ce qui donne la longueur d'onde \emph{minimale} de la série (et non maximale) ; ne pas inverser.
\end{corrige}

\begin{corrige}[Exercice 7 — Absorption d'un photon]
\begin{enumerate}
\item Énergie du photon absorbé :
\[
E = \frac{hc}{\lambda} = \frac{6,63\times 10^{-34}\times 3,00\times 10^{8}}{97,3\times 10^{-9}} \approx \SI{2,04e-18}{J}
\]
Conversion : $E = \dfrac{2,04\times 10^{-18}}{1,60\times 10^{-19}} \approx \SI{12,75}{eV}$.
\item L'atome part du fondamental $E_1 = -\SI{13,6}{eV}$. Après absorption, conservation de l'énergie :
\[
E_f = E_1 + E_{\text{photon}} = -13,6 + 12,75 = \SI{-0,85}{eV}
\]
On cherche $n$ tel que $E_n = -\dfrac{13,6}{n^2} = -0,85$ eV :
\[
n^2 = \frac{13,6}{0,85} = 16 \quad \Rightarrow \quad n = 4
\]
$n$ est bien un entier : l'absorption est possible et l'atome atteint le niveau \textbf{$n = 4$}. (Si $n$ n'avait pas été entier, le photon n'aurait pas pu être absorbé.)
\item La transition $1 \to 4$ aboutit sur le niveau final $n_f = 1$ : elle appartient à la \textbf{série de Lyman} (observée ici en absorption). Comme $\lambda = \SI{97,3}{nm} < \SI{400}{nm}$, elle se situe dans l'\textbf{ultraviolet}.
\end{enumerate}
\textbf{Point de vigilance :} en absorption, l'énergie du photon s'\emph{ajoute} algébriquement à l'énergie initiale (négative) de l'atome : $E_f = E_i + h\nu$. Le test « $n$ entier » est la justification attendue.
\end{corrige}

\begin{corrige}[Exercice 8 — Identification d'un spectre]
\begin{enumerate}
\item Avec $E(\text{eV}) = \dfrac{1240}{\lambda(\text{nm})}$ :
\begin{align*}
\lambda_1 = \SI{410}{nm} &\Rightarrow E_1 = 1240/410 \approx \SI{3,02}{eV} \\
\lambda_2 = \SI{434}{nm} &\Rightarrow E_2 = 1240/434 \approx \SI{2,86}{eV} \\
\lambda_3 = \SI{486}{nm} &\Rightarrow E_3 = 1240/486 \approx \SI{2,55}{eV} \\
\lambda_4 = \SI{656}{nm} &\Rightarrow E_4 = 1240/656 \approx \SI{1,89}{eV}
\end{align*}
\item Avec $E_2 = -\SI{3,40}{eV}$ et $E_n = -13,6/n^2$ :
\begin{align*}
E_3 - E_2 &= -1,51 - (-3,40) = \SI{1,89}{eV} \\
E_4 - E_2 &= -0,85 - (-3,40) = \SI{2,55}{eV} \\
E_5 - E_2 &= -0,544 - (-3,40) \approx \SI{2,86}{eV} \\
E_6 - E_2 &= -0,378 - (-3,40) \approx \SI{3,02}{eV}
\end{align*}
\item On constate une correspondance \emph{exacte} entre les énergies des photons observés et les différences $E_n - E_2$ :
\begin{center}
\begin{tabularx}{\linewidth}{|l|X|X|X|}\hline
\rowcolor{popblueL}
\textbf{Raie} & \textbf{Énergie (eV)} & \textbf{Transition} & \textbf{Nom} \\ \hline
$\lambda_4 = \SI{656}{nm}$ & 1,89 & $3 \to 2$ & H$\alpha$ (rouge) \\ \hline
$\lambda_3 = \SI{486}{nm}$ & 2,55 & $4 \to 2$ & H$\beta$ (bleu-vert) \\ \hline
$\lambda_2 = \SI{434}{nm}$ & 2,86 & $5 \to 2$ & H$\gamma$ (violet) \\ \hline
$\lambda_1 = \SI{410}{nm}$ & 3,02 & $6 \to 2$ & H$\delta$ (violet extrême) \\ \hline
\end{tabularx}
\end{center}
Le spectre observé est donc bien celui de l'\textbf{hydrogène} : ce sont les quatre raies visibles de la série de Balmer.
\item Ces quatre raies se terminent toutes sur $n_f = 2$ : elles appartiennent à la \textbf{série de Balmer}. Leurs longueurs d'onde ($410$ à $\SI{656}{nm}$) sont toutes comprises dans l'intervalle du visible ($400$ à $\SI{800}{nm}$), ce qui justifie qu'on les observe à l'œil ou avec un spectroscope ordinaire.
\end{enumerate}
\textbf{Point de vigilance :} la démarche d'identification est un classique du Bac : calculer les énergies des photons, calculer les écarts candidats entre niveaux, puis conclure par la coïncidence des valeurs (aux arrondis près).
\end{corrige}

\begin{corrige}[Exercice 9 — Type Bac : diagramme énergétique et raie H$\alpha$]
\begin{enumerate}
\item $E_n = -\dfrac{E_0}{n^2}$ avec $E_0 = \SI{13,6}{eV}$. Le nombre $n$ est le \textbf{nombre quantique principal} : un entier naturel non nul ($n \in \mathbb{N}^*$) qui étiquette les états stationnaires de l'électron et fixe leur énergie.
\item $E_1 = -\SI{13,6}{eV}$ ; $E_2 = -\dfrac{13,6}{4} = \SI{-3,40}{eV}$ ; $E_3 = -\dfrac{13,6}{9} \approx \SI{-1,51}{eV}$ ; $E_4 = -\dfrac{13,6}{16} = \SI{-0,85}{eV}$.
\item Diagramme à l'échelle 1 cm $\leftrightarrow$ 2 eV, avec l'origine $E = 0$ (niveau d'ionisation $n = \infty$) en haut :
\begin{itemize}
\item $n = 1$ : $-13,6$ eV, soit $6,8$ cm sous l'origine ;
\item $n = 2$ : $-3,40$ eV, soit $1,7$ cm sous l'origine ;
\item $n = 3$ : $-1,51$ eV, soit $0,76$ cm sous l'origine ;
\item $n = 4$ : $-0,85$ eV, soit $0,43$ cm sous l'origine.
\end{itemize}
Les traits horizontaux se resserrent vers le haut : l'écart entre niveaux consécutifs diminue quand $n$ croît. Chaque trait est légendé ($n$ et $E_n$).
\item Raie H$\alpha$, transition $3 \to 2$ :
\begin{enumerate}
\item $\Delta E = E_3 - E_2 = (-1,51) - (-3,40) = \SI{1,89}{eV}$, soit $E = 1,89 \times 1,60\times 10^{-19} \approx \SI{3,02e-19}{J}$.
\item $\lambda = \dfrac{hc}{E} = \dfrac{6,63\times 10^{-34}\times 3,00\times 10^{8}}{3,02\times 10^{-19}} \approx \SI{6,58e-7}{m} \approx \SI{658}{nm}$ (valeur de référence : $\SI{656}{nm}$).
\item $\lambda \approx \SI{656}{nm}$ correspond au \textbf{rouge} : c'est cohérent avec l'observation d'une « raie rouge intense ».
\end{enumerate}
\item Ionisation depuis le fondamental : $E_{\text{ion}} = E_\infty - E_1 = 0 - (-13,6) = \SI{13,6}{eV}$, soit $13,6 \times 1,60\times 10^{-19} \approx \SI{2,18e-18}{J}$.
\end{enumerate}
\textbf{Barème indicatif (sur 5 pts) :} 1) 0,5 pt ; 2) 1 pt ; 3) 1 pt (échelle respectée + resserrement des niveaux) ; 4a) 0,75 pt ; 4b) 0,75 pt ; 4c) 0,25 pt ; 5) 0,75 pt.
\textbf{Points de vigilance :} citer la formule avant de l'appliquer ; sur le diagramme, respecter l'échelle et placer le zéro en haut ; pour le photon émis, écrire $\Delta E = E_{n_i} - E_{n_f} > 0$ ; convertir les eV en joules avant d'utiliser $h$ en unités SI ; conclure chaque question par une phrase avec unité.
\end{corrige}

\begin{corrige}[Exercice 10 — Type Bac : désexcitation depuis le niveau $n = 4$]
\begin{enumerate}
\item $E_4 = -\dfrac{13,6}{4^2} = -\dfrac{13,6}{16} = \SI{-0,85}{eV}$.
\item Les atomes au niveau $n = 4$ peuvent redescendre vers le fondamental par des sauts directs ou en cascade, en passant par les niveaux intermédiaires $n = 3$ et $n = 2$. Sur le diagramme, on trace toutes les flèches descendantes possibles entre les niveaux 4, 3, 2, 1 :
\[
4\to1, \quad 4\to2, \quad 4\to3, \quad 3\to1, \quad 3\to2, \quad 2\to1
\]
soit \textbf{6 transitions distinctes}, donc au maximum \textbf{6 raies d'émission} différentes (formule générale : $\dfrac{n(n-1)}{2} = \dfrac{4\times 3}{2} = 6$).
\item La plus grande longueur d'onde correspond à la plus petite énergie émise : transition $4 \to 3$.
\[
\Delta E = E_4 - E_3 = (-0,85) - (-1,51) = \SI{0,66}{eV}
\]
\[
\lambda_{\max} = \frac{1240}{0,66} \approx \SI{1879}{nm} \approx \SI{1,88}{\mu m}
\]
Domaine : \textbf{infrarouge} ($\lambda > \SI{800}{nm}$).
\item La plus petite longueur d'onde correspond à la plus grande énergie : transition $4 \to 1$.
\[
\Delta E = E_4 - E_1 = (-0,85) - (-13,6) = \SI{12,75}{eV}
\]
\[
\lambda_{\min} = \frac{1240}{12,75} \approx \SI{97,3}{nm}
\]
Domaine : \textbf{ultraviolet} ($\lambda < \SI{400}{nm}$).
\item Cherchons les raies dont $\lambda \in [\SI{400}{nm} ; \SI{800}{nm}]$, c'est-à-dire dont l'énergie est comprise entre $\dfrac{1240}{800} = \SI{1,55}{eV}$ et $\dfrac{1240}{400} = \SI{3,10}{eV}$ :
\begin{itemize}
\item $4\to2$ : $\Delta E = (-0,85)-(-3,40) = \SI{2,55}{eV}$, $\lambda = 1240/2,55 \approx \SI{486}{nm}$ : \textbf{visible} (H$\beta$, bleu-vert) ;
\item $3\to2$ : $\Delta E = (-1,51)-(-3,40) = \SI{1,89}{eV}$, $\lambda = 1240/1,89 \approx \SI{656}{nm}$ : \textbf{visible} (H$\alpha$, rouge) ;
\item $4\to3$ : $\SI{1879}{nm}$ (IR) ; $4\to1$, $3\to1$, $2\to1$ : $\lambda \leq \SI{122}{nm}$ (UV).
\end{itemize}
\textbf{Seules 2 raies sont visibles}, ce sont précisément les raies de Balmer accessibles depuis $n = 4$.
\item Classement par série :
\begin{itemize}
\item \textbf{Lyman} ($n_f = 1$) : $4\to1$, $3\to1$, $2\to1$ — 3 raies, toutes UV ;
\item \textbf{Balmer} ($n_f = 2$) : $4\to2$, $3\to2$ — 2 raies, visibles ;
\item \textbf{Paschen} ($n_f = 3$) : $4\to3$ — 1 raie, IR.
\end{itemize}
\end{enumerate}
\textbf{Barème indicatif (sur 6 pts) :} 1) 0,5 pt ; 2) 1,5 pt (diagramme + dénombrement) ; 3) 1 pt ; 4) 1 pt ; 5) 1,5 pt ; 6) 0,5 pt.
\textbf{Points de vigilance :} ne pas oublier les transitions en cascade (un atome en $n=4$ peut descendre en plusieurs sauts, peuplant les niveaux 3 et 2 qui se désexcitent à leur tour) ; « plus grande longueur d'onde » = « plus petite énergie » (proportionnalité inverse) ; justifier l'appartenance au visible par le \emph{calcul} de $\lambda$, pas seulement par le nom de la série.
\end{corrige}

\begin{corrige}[Exercice 11 — Type Bac : photo-ionisation de l'hydrogène]
\begin{enumerate}
\item Photon $E = \SI{15,0}{eV}$ sur un atome à l'état fondamental :
\begin{enumerate}
\item L'énergie d'ionisation depuis $n = 1$ vaut $E_{\text{ion}} = E_\infty - E_1 = \SI{13,6}{eV}$. Comme $E = \SI{15,0}{eV} > \SI{13,6}{eV}$, le photon transporte assez d'énergie pour arracher l'électron : \textbf{l'atome peut être ionisé}.
\item L'énergie du photon sert d'abord à ioniser ($\SI{13,6}{eV}$) ; l'\textbf{excédent} est emporté par l'électron éjecté sous forme d'\textbf{énergie cinétique} (c'est l'analogue atomique de l'effet photoélectrique).
\item $E_c = E - E_{\text{ion}} = 15,0 - 13,6 = \SI{1,4}{eV}$, soit $E_c = 1,4 \times 1,60\times 10^{-19} = \SI{2,24e-19}{J}$.
\item En mécanique classique : $E_c = \dfrac{1}{2}m_e v^2$, donc
\[
v = \sqrt{\frac{2E_c}{m_e}} = \sqrt{\frac{2\times 2,24\times 10^{-19}}{9,11\times 10^{-31}}} = \sqrt{4,92\times 10^{11}} \approx \SI{7,0e5}{m/s}
\]
Vérification de l'hypothèse : $v/c \approx \num{2,3e-3} \ll 1$, l'approximation non relativiste est bien valide.
\end{enumerate}
\item $\lambda = \dfrac{hc}{E} = \dfrac{1240}{15,0} \approx \SI{82,7}{nm}$. Comme $\lambda < \SI{400}{nm}$ (et même $< \SI{100}{nm}$), ce photon appartient à l'\textbf{ultraviolet lointain}.
\item Pour ioniser depuis le fondamental, il faut $E_{\text{photon}} \geq \SI{13,6}{eV}$, c'est-à-dire $\lambda \leq \dfrac{hc}{E_{\text{ion}}}$ :
\[
\lambda_{\max} = \frac{1240}{13,6} \approx \SI{91,2}{nm}
\]
Tout photon de longueur d'onde inférieure ou égale à $\SI{91,2}{nm}$ peut ioniser l'hydrogène fondamental (c'est la limite de la série de Lyman).
\item Photon $E = \SI{12,1}{eV}$ : comme $12,1 < 13,6$ eV, \textbf{l'ionisation est impossible}. Voyons si une excitation est possible : après absorption éventuelle, l'atome aurait l'énergie
\[
E_f = E_1 + E = -13,6 + 12,1 = \SI{-1,5}{eV}
\]
Or $E_3 = -\dfrac{13,6}{9} \approx \SI{-1,51}{eV}$ : la coïncidence est excellente (l'écart de $\SI{0,01}{eV}$ vient des arrondis des données). Le photon est donc \textbf{absorbé} et l'atome est porté au niveau \textbf{$n = 3$} (transition $1 \to 3$). Cette raie appartient à la \textbf{série de Lyman} (raie Lyman-$\beta$, en absorption) ; sa longueur d'onde est $\lambda = \dfrac{1240}{12,1} \approx \SI{102,5}{nm}$, dans l'ultraviolet. L'atome excité se désexcitera ensuite en réémettant des photons (cascade $3\to2$ puis $2\to1$, ou $3\to1$ direct).
\end{enumerate}
\textbf{Barème indicatif (sur 6 pts) :} 1a) 0,5 pt ; 1b) 0,5 pt ; 1c) 0,75 pt ; 1d) 0,75 pt ; 2) 0,75 pt ; 3) 1 pt ; 4) 1,75 pt.
\textbf{Points de vigilance :} comparer explicitement l'énergie du photon à l'énergie d'ionisation \emph{avant} de conclure ; préciser que l'excédent devient de l'énergie \emph{cinétique} ; citer la condition d'application de $E_c = \frac{1}{2}m_ev^2$ (régime non relativiste) et la vérifier a posteriori ; pour la question 4, ne pas s'arrêter à « pas d'ionisation » : tester l'absorption résonnante vers un état excité, c'est la justification attendue.
\end{corrige}

\newpage

\coursec{Fiche bilan}

\begin{popfigure}
\begin{tikzpicture}[mindmap, grow cyclic, every node/.style={concept, font=\small}, text=white,
  root concept/.append style={concept color=popdark, font=\small\bfseries},
  level 1/.append style={level distance=3.4cm, sibling angle=60, font=\footnotesize},
  level 2/.append style={level distance=2.1cm, sibling angle=45, font=\scriptsize}]
\node[root concept]{Niveaux d'énergie de l'atome H}
  child[concept color=popblue]{ node {Photon}
    child { node {$E=h\nu=\dfrac{hc}{\lambda}$} }
    child { node {$h=\SI{6,63e-34}{J.s}$} }
  }
  child[concept color=poporange]{ node {Niveaux $E_n$}
    child { node {$E_n=-\dfrac{E_0}{n^2}$} }
    child { node {$E_0=\SI{13,6}{eV}$} }
  }
  child[concept color=popgreen]{ node {États}
    child { node {fondamental $n=1$} }
    child { node {excités $n\geq 2$} }
    child { node {ionisation $n\to\infty$} }
  }
  child[concept color=poppurple]{ node {Transitions}
    child { node {émission : photon émis} }
    child { node {absorption : photon absorbé} }
    child { node {$h\nu=|E_f-E_i|$} }
  }
  child[concept color=poppink]{ node {Spectres}
    child { node {raies d'émission} }
    child { node {raies d'absorption} }
    child { node {spectres discontinus} }
  }
  child[concept color=popgold]{ node {Séries}
    child { node {Lyman $n\to1$ (UV)} }
    child { node {Balmer $n\to2$ (visible)} }
    child { node {Paschen $n\to3$ (IR)} }
  };
\end{tikzpicture}
\legende{Carte mentale de la leçon : les six idées clés autour des niveaux d'énergie de l'atome d'hydrogène.}
\end{popfigure}

\begin{bilan}
\textbf{Définitions clés}
\begin{itemize}
  \item \cle{Photon} : grain d'énergie lumineuse, sans masse, se propageant à la célérité $c=\SI{3,00e8}{m.s^{-1}}$.
  \item \cle{Quantification} : l'énergie de l'atome ne peut prendre que des valeurs discrètes $E_n$ ; les échanges d'énergie se font par \emph{quanta}.
  \item \cle{État fondamental} ($n=1$) : état de plus basse énergie, $E_1=-\SI{13,6}{eV}$. \cle{États excités} : $n\geq 2$. \cle{Ionisation} : l'électron est arraché ($n\to\infty$, $E_\infty=0$) ; énergie d'ionisation $E_i=\SI{13,6}{eV}$.
  \item \cle{Transition} : passage d'un niveau $E_i$ à un niveau $E_f$ avec émission ($E_i>E_f$) ou absorption ($E_i<E_f$) d'un photon.
  \item \cle{Spectre de raies} : spectre discontinu de raies colorées (émission) ou de raies noires sur fond continu (absorption), signature de l'élément.
\end{itemize}

\textbf{Formules à connaître par cœur}
\begin{center}
\begin{tabularx}{\linewidth}{|l|X|X|}\hline
\rowcolor{popblueL} \textbf{Grandeur} & \textbf{Formule} & \textbf{Remarques} \\ \hline
Énergie du photon & $E=h\nu=\dfrac{hc}{\lambda}$ & $h=\SI{6,63e-34}{J.s}$ ; $\lambda$ en m \\ \hline
Niveaux de H & $E_n=-\dfrac{E_0}{n^2}$, $E_0=\SI{13,6}{eV}$ & $n\in\mathbb{N}^*$ ; énergies négatives \\ \hline
Transition & $h\nu=\dfrac{hc}{\lambda}=|E_f-E_i|$ & émission si $E_i>E_f$ \\ \hline
Longueur d'onde & $\dfrac{1}{\lambda}=\dfrac{E_0}{hc}\left|\dfrac{1}{n_f^2}-\dfrac{1}{n_i^2}\right|$ & formule de Ritz--Balmer \\ \hline
Conversion & $\SI{1}{eV}=\SI{1,60e-19}{J}$ & indispensable aux calculs \\ \hline
\end{tabularx}
\end{center}

\textbf{Séries spectrales de l'hydrogène}
\begin{itemize}
  \item \cle{Lyman} : retour vers $n=1$ $\Rightarrow$ ultraviolet ; plus grande raie : $2\to1$, $\lambda\approx\SI{121,6}{nm}$.
  \item \cle{Balmer} : retour vers $n=2$ $\Rightarrow$ visible ; raie H$_\alpha$ rouge : $3\to2$, $\lambda\approx\SI{656,3}{nm}$.
  \item \cle{Paschen} : retour vers $n=3$ $\Rightarrow$ infrarouge.
\end{itemize}

\textbf{Méthodes express}
\begin{enumerate}
  \item Écrire $E_i$ et $E_f$ avec $E_n=-13,6/n^2$ (en eV), calculer $\Delta E=|E_f-E_i|$.
  \item En déduire $\lambda=\dfrac{hc}{\Delta E}$ (convertir $\Delta E$ en joules !).
  \item Comparer $\lambda$ au domaine visible ($400$--$\SI{800}{nm}$) pour identifier la série et la couleur.
  \item Absorption : seuls les photons d'énergie \emph{exactement} égale à un écart de niveaux sont absorbés.
  \item Ionisation depuis le niveau $n$ : il faut fournir au minimum $E_{\text{ion}}=\dfrac{13,6}{n^2}\ \text{eV}$.
\end{enumerate}
\end{bilan}

\begin{aretenir}[Checklist avant le Bac]
\begin{itemize}
  \item[$\square$] Je sais énoncer la relation de Planck--Einstein $E=h\nu=hc/\lambda$ avec les unités SI.
  \item[$\square$] Je convertis sans erreur eV $\leftrightarrow$ J ($\SI{1}{eV}=\SI{1,60e-19}{J}$).
  \item[$\square$] Je calcule $E_n=-13,6/n^2$ pour tout $n$ et je place le niveau sur un diagramme énergétique.
  \item[$\square$] Je sais que les énergies des niveaux sont \textbf{négatives} et croissent vers $0$ quand $n$ augmente.
  \item[$\square$] Je distingue état fondamental, états excités et ionisation ($E_\infty=0$).
  \item[$\square$] Je sais qu'une émission correspond à $E_i>E_f$ (photon émis) et une absorption à $E_i<E_f$.
  \item[$\square$] Je calcule $\lambda=hc/|\Delta E|$ d'une transition et je vérifie l'ordre de grandeur.
  \item[$\square$] Je connais les séries de Lyman ($n\to1$, UV), Balmer ($n\to2$, visible) et Paschen ($n\to3$, IR).
  \item[$\square$] Je sais interpréter un spectre de raies d'émission et un spectre d'absorption (raies noires aux \emph{mêmes} longueurs d'onde).
  \item[$\square$] Je sais que le spectre est discontinu parce que les niveaux d'énergie sont quantifiés.
  \item[$\square$] Je sais que l'énergie d'ionisation depuis l'état $n$ vaut $13,6/n^2$ eV.
  \item[$\square$] Je rédige proprement : schéma du diagramme, flèche de transition, calcul avec unités et chiffres significatifs.
\end{itemize}
\end{aretenir}

\findecours

\end{document}
